% before: sha-tate-shafarevich-group
The local error terms of the window bound, at the places where the characters of a host field are twisted, are measured by Galois cohomology classes that are unramified outside $\ell$ and locally trivial at $\ell$. For the cyclotomic twist $-k$ these classes form the group below. Its length is the quantity that the rest of the chapter bounds.

% before: sha-plus-minus
Over a CM field, complex conjugation acts on the Tate-Shafarevich groups, and the two eigenspaces behave differently. On $\mathbf Q(\zeta_\ell)$ the minus part corresponds to class groups with odd characters, which are controlled by Bernoulli numbers, and the plus part to the real class group of Vandiver's conjecture. The sign $(-1)^k$ accounts for the action of complex conjugation on the coefficients.

% before: sha-plus-minus-length
Away from $\ell=2$ the involution can be diagonalized over $\mathbf Z_\ell$. So the length of each Tate-Shafarevich group is the sum of the lengths of its two parts, and the parts can be bounded separately.

% before: sha-bounded-field
The bounds are stated as weighted sums over primes and twists, with the weight $\log\ell$ that the window bound requires. Separate versions for a single prime and for the two eigenspaces are needed because the three parts of the proof produce them.

% before: sha-bound-from-parts
A bound at the prime $2$, a bound for the minus parts and a bound for the plus parts at odd primes combine into the full bound. The prime $2$ enters with weight $\log2$, as in the weighted sum.

% before: sha-vandiver-weakening
The host fields of Chapter~\ref{chapter-host-fields} are defined at each scale $x$ from a primitive character $\psi$, a $2$-power exponent $e(x)$ and a pair of primes $(q,d)$. The window bound needs the Tate-Shafarevich bound for twists up to a power of $x$, on a window of scales $[Y,Y^2]$ that can be found inside every long enough range $[X^{1/J},X]$. The bound allowed, $e^{x^{\varepsilon}/2}$, is much larger than anything polynomial in the conductor, which leaves room for the exceptional twists of the Iwasawa argument.

% before: sha-abelian-field
The arguments for the minus part and for the prime $2$ apply to composita in which the field of $\psi$ is replaced by an arbitrary abelian field. Abelian fields are described by subgroups of $(\mathbf Z/m)^\times$, and the composita by the same congruence conditions as host fields.

% before: sha-composita
Host fields are the composita built from the field of $\psi$, so the general results apply to them. The proofs also need that host fields are CM fields with conductor bounded in terms of $e$, $N_\psi$, $q$ and $d$, and that the only odd primes dividing their degree are those dividing the degree of $\psi$.

% before: sha-bernoulli-bound
The minus part splits further into eigenspaces for the characters of the Galois group, when $\ell$ does not divide its order. Each eigenspace is bounded by the valuation of a generalized Bernoulli number, by the Euler system of Gauss sums.

% before: sha-minus-prime-to-degree
For primes $\ell$ not dividing the degree, the Bernoulli bounds for all characters can be added over $\ell$. The product of the Bernoulli numbers over a Galois-stable set of characters is a nonzero integer, so the weighted sum over all $\ell$ is bounded by the logarithm of its size.

% before: sha-stickelberger-elements
For the odd primes dividing the degree the eigenspace decomposition is not available, and the minus part is bounded through an explicit annihilator. The annihilator comes from Stickelberger elements: a higher one, whose characters give generalized Bernoulli numbers, and a twisted one of large level with integer coefficients.

% before: sha-stickelberger-annihilator
Stickelberger's theorem kills the class groups of cyclotomic fields, and the classes of the Tate-Shafarevich group become homomorphisms on the subgroup fixing the $\ell$-power roots of unity, where class field theory applies. The twisted elements of level $N\ell^m$ converge $\ell$-adically to the higher Stickelberger element, which therefore kills the Tate-Shafarevich group up to the factor $\varphi(N)$.

% before: sha-stickelberger-determinant
On the minus part complex conjugation acts by a sign, so an annihilating element of the group ring can be replaced by its determinant on the corresponding quotient of the group ring, which is a scalar. For the higher Stickelberger element this determinant is a product of generalized Bernoulli numbers with nonzero factors, and its size is bounded explicitly.

% before: sha-torsion-length
An exponent for the minus part bounds its length once the number of generators is bounded. That number is the length of the $\ell$-torsion of $H^1$, which Kummer theory over $E(\mu_\ell)$ bounds by units and the class number.

% before: sha-minus-exponent
The odd primes dividing the degree of a host field divide the order of $\psi$, so there are finitely many of them, independently of the scale. For these a cruder bound suffices, polynomial in the degree and in the twist: an exponent that kills the minus part and a bound for the number of generators.

% before: sha-minus-part
The two lemmas together bound the minus parts of the composita for all odd primes. The bound is polynomial in the degree, the twist and the conductor, which is $\exp(O(x^{\varepsilon/4}))$ on host fields.

% before: sha-ferrero-washington
The prime $2$ and the primes dividing the degree are treated through Iwasawa theory in cyclotomic towers. Its input is that the class groups along a cyclotomic $\mathbf Z_\ell$-extension of an abelian field have bounded $\ell$-rank.

% before: sha-tower-data
A host field at scale $x$ sits in a layer of the $2$-power cyclotomic tower over a fixed field. Along a tower, the Tate-Shafarevich groups of all layers and all twists are controlled by one Iwasawa module with a bounded number of generators. The definition records the two bounds that are used: one for twists away from finitely many exceptional $\ell$-adic numbers, and one for all twists in terms of $H^2$ of a fixed layer.

% before: sha-iwasawa-module
The Tate-Shafarevich groups of all layers of a cyclotomic tower are controlled by a single module over the Galois group of the tower. It is the Galois group of the maximal unramified abelian pro-$\ell$ extension of the top of the tower in which the places above $\ell$ split completely.

% before: sha-iwasawa-descent
The Tate-Shafarevich group of a layer is recovered from the Iwasawa module by taking coinvariants under the Galois group of the tower over the layer, twisted by the cyclotomic character. Only the length of a cokernel is needed later.

% before: sha-cokernel-growth
Going up one step in the tower raises the power of the generator by a factor $\ell$. Once the generator is close enough to $1$, each step adds at most the number of generators of the module to the length of the cokernel.

% before: sha-exceptional-centres
For a fixed layer, the cokernel can be large only for twists $k$ that are $\ell$-adically close to finitely many numbers. These numbers come from the eigenvalues of the generator on the module, and there are at most as many of them as generators.

% before: sha-tower-from-module
Descent, growth and the exceptional numbers together give the two conditions of a tower bound. The comparison with $H^2$ in the second condition comes from Poitou-Tate duality at a fixed layer.

% before: sha-tower-bounds
The tower bounds hold for the fields of abelian field data containing $\mu_\ell$ at odd $\ell$. At $\ell=2$ the base field depends on the auxiliary primes $q$ and $d$, which vary with the scale; the number of generators must not.

% before: sha-kummer-layers
The bound of Soulé type is proved by passing to cyclotomic extensions in which Kummer theory turns classes into homomorphisms on class groups. Counting these homomorphisms reduces to counting elements of the field with prescribed divisors up to $\ell^\nu$-th powers, which the homology of $GL_N$ of the ring of integers controls.

% before: sha-bar-chains
Soulé's method detects elements of a number field by evaluating homology classes of $GL_N$ of its ring of integers at primes of degree one. These classes are built by explicit operations on chains, so the inhomogeneous chains, the prism operator for conjugation and an explicit cycle of the cyclic group are fixed first. The cycle $z_{q,e}$ represents a generator of the homology of $\mathbf Z/q$ in degree $2e+2$.

% before: sha-bar-identities
The constructions below combine chains by hand and must know which combinations are cycles or boundaries. The prism operator is a chain homotopy between conjugation and the identity, and it is additive in the conjugating element up to boundaries.

% before: sha-bar-relation
The homology of $\mathbf Z/q\times\mathbf Z/q$ in degree $2e+2$ has fewer generators than the $2e+4$ images of $z_{q,e}$ under the maps $\Delta_t$. A relation among them, with a unit coefficient, lets a combination of prisms over different conjugating elements be moved into a single arithmetic group.

% before: sha-eta
The relation of Lemma~\ref{sha-bar-relation} has an unknown index $t_0$ with a unit coefficient. Divisors weighted by the function $\eta$ below are seen by the projection $\pi_{t_0}$ as the coefficients of the polynomial whose roots are the powers of $1+\ell^c$, and by the other projections as zero, which isolates the term of index $t_0$.

% before: sha-additive-family
The homology classes are evaluated at a prime of degree one through linear forms on the homology of $GL$ of the residue field. These forms must be compatible with direct sums and must take prescribed values on classes coming from the scalars, and over the ring of integers they give the compatible families below.

% before: sha-transfer
A linear form on the homology of the multiplicative group of a finite field extends to all general linear groups over that field. The extension is a transfer through the action on lines, and it is additive because the number of lines in a fibre is a power of the order of the field.

% before: sha-cocycle
The linear form to be extended is the cup product of a two-cocycle on the cyclic group $k^\times$ with a homomorphism. Its value on the prism of the cycle $z_{q,e}$ is computed by a telescoping sum.

% before: sha-theta
Extending the cup product by Lemma~\ref{sha-transfer} gives forms on all general linear groups whose values on prisms of scalars are given by the determinant. This is where the residue of an element of the number field will appear.

% before: sha-primitive
Composed with reduction modulo a prime, these forms give families over the ring of integers that are compatible with stabilization and multiply by $d$ under the $d$-fold block sum. These two properties are what the rank bound for the homology of $GL_N(\mathcal O_L)$ needs.

% before: sha-steinitz
The construction moves lattices of the form $\bigoplus_j\mathfrak a_j$ into one another by linear maps with prescribed determinant. This is possible when the total divisors differ by the divisor of the determinant.

% before: sha-detect
A count of elements up to $\ell^\nu$-th powers needs a way to recognize $\ell^\nu$-th powers from their residues. By Chebotarev's density theorem it suffices to look at primes of degree one.

% before: sha-class-cycle
The central construction attaches to each element $\alpha$ with divisor of the form $(\gamma-a)\mathfrak d$ a cycle in the homology of $GL_N(\mathcal O_L)$ whose value, at every prime of degree one away from finitely many, is a fixed multiple of the residue of the twisted norm of $\alpha$. The cycle is the image of a sum of prisms of lattice automorphisms, made into a cycle of the stabilizer of $\mathcal O_F^J$ by the relation of Lemma~\ref{sha-bar-relation}.

% before: sha-borel-rank
The number of possible values of these evaluations is bounded by the rank of the homology of $GL_N(\mathcal O_L)$ in odd degree. After removing products of lower-degree classes, this rank is $r_2(L)$ by Borel's computation.

% before: sha-few-values
With the rank of Borel's theorem and homological stability for the torsion, the values of a compatible family on cycles of a fixed degree can be counted. The factor $\ell^k$ is the price of removing the decomposable classes.

% before: sha-few-norms
Combining the cycles of Proposition~\ref{sha-class-cycle} with the counting of Lemma~\ref{sha-few-values}, a set of elements with distinct twisted norms modulo $\ell^\nu$-th powers is bounded by the number of values of one compatible family. Lemma~\ref{sha-detect} is what makes the evaluations injective.

% before: sha-count-homs
The homomorphisms on the class group that come from Tate-Shafarevich classes are compared with the elements whose twisted norms were counted. The comparison is a computation in the group of divisors and the unit group, in which the unit theorem as a Galois module contributes the factor $(\ell^\nu)^{r_2(L)}$ that also appears in Lemma~\ref{sha-few-norms}.

% before: sha-kummer-descent
Kummer theory turns classes of the Tate-Shafarevich group into homomorphisms on class groups of the layers $L(\mu_{\ell^\nu})$ with a prescribed action of the generator $\gamma$. A uniform bound for the number of such homomorphisms bounds the length of the group.

% before: sha-enlarging
The counting needs more roots of unity and more complex places than the given field may have, compared with the twist $k$. Passing to a cyclotomic layer of degree about $k$ provides them, at the cost of explicit bounds on degree and discriminant, and can only increase the Tate-Shafarevich group.

% before: sha-convex-nerve
Following Soulé, the torsion in the homology of $GL_N(\mathcal O_L)$ is bounded through an action on an acyclic complex with finitely many orbits of cells and finite stabilizers. The complex is the nerve of a cover of a convex set by closed convex sets, and its acyclicity reduces, by cutting the convex set along a hyperplane, to covers in which one set contains a point common to all the others.

% before: sha-polyhedron-nerve
The convex set will be a cone of linear forms, covered by the normal cones at the vertices of a locally finite polyhedron. A linear form that is positive on the cone containing the polyhedron has compact sublevel sets on it, and this gives the finiteness that the nerve lemma requires.

% before: sha-hermitian-polyhedron
For $GL_N(\mathcal O_L)$ the polyhedron is that of Voronoi's theory of perfect forms, for tuples of Hermitian forms indexed by the complex places of $L$. Its vertices have finitely many minimal vectors spanning $L^N$, and this makes the stabilizers finite and the orbits of tuples of vertices countable in terms of the degree and the discriminant.

% before: sha-torsion-from-action
An exact complex of permutation modules with finitely many orbits and bounded stabilizers bounds the torsion of group homology. Shapiro's lemma reduces the permutation modules to finite groups, whose homology is killed by their order, and the lowest degree is a computation with lattices.

% before: sha-gl-torsion
The count of Lemma~\ref{sha-few-norms} involves the torsion of the homology of $GL_N(\mathcal O_L)$ in degrees about $2k$. Lemmas~\ref{sha-hermitian-polyhedron} and~\ref{sha-torsion-from-action} bound it explicitly in terms of the degree and the discriminant of $L$, as Step 1 of the proof of Theorem~\ref{sha-soule-bound} requires.

% before: sha-soule-bound
At the exceptional twists the tower bound refers to $H^2(G_{L,S_\ell},\mathbf Z_\ell(k+1))$ of a fixed layer, and a bound for this group is needed that is effective in the twist $k$. The bound below is doubly exponential in a power of $k$, so it is useful only when $k$ is small compared with the scale.

% before: sha-window-pigeonhole
The exceptional twists depend on the tower and on the range of scales, but there are at most $R$ of them. A pigeonhole argument over nested ranges of scales finds a window on which none of them lies between $x^{\eta}$ and $x^{C_t}$, so every exceptional twist that matters is small enough for the bound of Soulé type.

% before: sha-window-constants
The numbers of exceptional twists for the prime $2$ and for the bad primes determine how far down the pigeonhole may have to go. Window constants record these numbers and the constant of the bound of Soulé type, and controlled windows are the ranges of scales on which the bounds at $2$ and at the bad primes hold.

% before: sha-prime-two
This is the part of the argument that uses windows of scales. Combining the tower bounds, the bound of Soulé type and the pigeonhole gives, for every range $[X^{1/J},X]$, a window on which the Tate-Shafarevich groups at $2$ of the composita, and those at the bad primes of fixed layers, are bounded by $\exp(x^{4\varepsilon/5})$.

% before: sha-character-parts
The plus part of a host field, like the minus part, splits according to the characters of its Galois group over $\mathbf Q$. By Shapiro's lemma each piece becomes a Tate-Shafarevich group over $\mathbf Q$ of a twist of $\mathbf Q_\ell/\mathbf Z_\ell(-k)$ by a Dirichlet character. These groups over $\mathbf Q$ are what the Eisenstein-congruence arguments bound.

% before: sha-plus-by-characters
The plus part of a host field is bounded by the sum of its character pieces. When $\ell$ divides the degree the decomposition is not exact, and a loss polynomial in the conductor and the twist is added for each character.

% before: sha-control
For the trivial character with $\ell-1\mid k$, the classes become homomorphisms after restriction to the cyclotomic $\mathbf Z_\ell$-extension. Restriction to the subgroup fixing the $\ell$-power roots of unity loses no information when the twist is nonzero, because the remaining group is procyclic and the coefficients are divisible.

% before: sha-trivial-character
For the trivial character and twists $k$ with $\ell-1\mid2k$, the group over $\mathbf Q$ vanishes for elementary reasons. These are the twists at which the cyclotomic character is congruent to the trivial or to a quadratic character modulo $\ell$, where the Eisenstein-congruence method does not apply.

% before: sha-ordinary-deformation
For a nontrivial character, a nonzero class in the group over $\mathbf Q$ gives a nonsplit extension, and hence a residually reducible representation of type $\Psi\oplus1\oplus\Psi^{-1}$ of the group $\mathcal G_3$ attached to an imaginary quadratic field. Its ordinary deformations are the objects of the automorphy theorem that follows.

% before: sha-hecke-situation
The automorphy of ordinary points in Theorem~\ref{sha-residually-reducible-automorphy} is proved over CM fields, through a universal ring of ordinary determinants lifting a sum of distinct characters. A determinant is defined for every lift of a reducible residual representation, including lifts that do not come from a representation over the base ring. The places that may ramify are those above $\ell$, some auxiliary places, and one place $v_0$ at which the level of the automorphic forms is the principal congruence subgroup.

% before: sha-hecke-arcs
The Hecke ideal is controlled through two kinds of points of the universal ring: arithmetic points, which carry the weights of automorphic forms, and arcs, which are one-parameter families of representations in characteristic $\ell$. Patching bounds the ring near an arc whose representation is irreducible on the subgroup where the cyclotomic character is trivial, or is a sum of two such pieces with no maps between them. Thick primes are the non-semisimple points through which such arcs are joined.

% before: sha-hecke-data
Hida's Hecke algebra on a definite unitary group needs three choices: a type at the auxiliary places that cuts out the inertial behaviour of the residual characters, a weight character of the torus at $\ell$ through which the Iwasawa algebra acts, and residual eigenvalues for the Hecke operators. Each is determined by the residual datum.

% before: sha-hecke-algebra
The Hecke algebra must be finite over the Iwasawa algebra and detected by classical points, so that relations proved at classical points hold on the whole algebra. Both come from Hida's theory of ordinary forms, which makes the ordinary part free of finite rank over the Iwasawa algebra.

% before: sha-hecke-local-ordinary
At a place $v\mid\ell$ the Hecke operators $U_{v,j}$ come from the diagonal matrices $\alpha_j$, acting on vectors fixed by an Iwahori-type subgroup on which the diagonal units act through a character. The flags at $\ell$ of the Galois representations of the Hecke algebra are read off from the local representations that carry such vectors in the image of every power of $U_z$.

% before: sha-hecke-local-principal
An irreducible representation with a nonzero ordinary vector, for a character of the diagonal units with distinct components, is a principal series. The Jacquet module identifies the ordinary vectors with a line in it, and the action of the torus on that line gives the eigenvalues of the operators $U_{\alpha_j}$.

% before: sha-hecke-local-rec
The local Langlands correspondence sends a principal series to the sum of its characters, and the characteristic polynomials of a Weil-Deligne parameter determine a supercuspidal support made of characters. Both follow from the laws of the correspondence and Zelevinsky's classification.

% before: sha-hecke-local-global
At a place above $\ell$, the Galois representation of an automorphic representation is compatible with the local Langlands correspondence when the local parameter is a sum of characters without monodromy. For a principal series whose characters are distinct on the units this condition holds, so the Weil-Deligne representation on the Galois side is the sum of those characters.

% before: sha-hecke-local-characters
The characters $\psi_{\mathbf T,v,i}$ are defined by their values on inertia and at a uniformizer, through the weight character and the operators $U_{v,j}$. Once the principal series at $v$ is known, these values identify each $\psi_{\mathbf T,v,i}$ at a classical point with a principal series character times an algebraic character.

% before: sha-hecke-local-flag
A de Rham representation whose Weil-Deligne representation is a sum of characters distinct on inertia, with Hodge-Tate weights increasing in the order of the characters, is triangular. Its invariant lines are found one at a time: after a twist, the character with the smallest weight becomes trivial and all weights become nonnegative, and the fundamental exact sequence turns the corresponding crystalline period into an invariant vector.

% before: sha-hecke-determinant
To compare the Hecke algebra with the universal ring, the Hecke algebra must carry an ordinary determinant. The Galois representations of the classical points glue to a determinant over the whole Hecke algebra, because the classical points of a fixed weight separate its elements.

% before: sha-hecke-ideal-speh
The Galois representations at the arithmetic points on the zero locus of the Hecke ideal come from automorphic representations of $U(n)$, whose base changes to $GL_n$ need not be cuspidal. The Speh representations in the base change determine the shape of these Galois representations, and when all of them have length $1$, the Jacquet-Shalika bound separates the absolute values of any two Frobenius eigenvalues by less than the norm of the place.

% before: sha-hecke-ideal-inertia
A point of the universal ring whose determinant comes from a semisimple representation constrains that representation on inertia at the auxiliary places. The relations of the universal ring make it trivial where the residual characters are trivial, and the residual characters then determine it through their Teichm\"uller lifts. This is how a local component at an auxiliary place comes to contain the type $\tau_0$.

% before: sha-hecke-ideal-forms
An automorphic representation of the definite unitary group contributes to the Hecke module when it contains classical forms of the right weight, type and level. The local conditions at the auxiliary places, at $v_0$ and outside $S$ are the conditions for such forms to exist.

% before: sha-hecke-ideal-unitary
A contributing automorphic representation of $U(n)$ needs a Galois representation to compare with the given point. It comes from the base change to $GL_n$ and the Galois representations attached to the cuspidal pieces, and at the split places where the automorphic representation is spherical its Frobenius eigenvalues are given by the Satake parameters. Its Hodge-Tate multisets are given by the weight at infinity, and it has them at the split places above $\ell$ where a local component is plain.

% before: sha-hecke-ideal-ordinary-vector
At a place above $\ell$, a contributing representation must carry an ordinary vector for the character of the diagonal units given by the point. A principal series has one, and the eigenvalues of the operators $U_{\alpha_j}$ on it are computed from its characters, in the converse direction to Lemma~\ref{sha-hecke-local-principal}.

% before: sha-hecke-ideal-eigenform
An ordinary eigenform with the residual eigenvalues inside a contributing representation produces a point of the Hecke algebra. Hida's control theorem identifies the classical ordinary forms at the point with the specializations of the $\Lambda$-adic forms, so the eigenform gives a character of the Hecke algebra on $\Lambda$-adic forms, which factors through the localization at the residual eigenvalues.

% before: sha-hecke-ideal-attached
For the automorphic points of the universal ring, the Galois representation is only known to match a cuspidal representation at almost all places. It then has the same Frobenius characteristic polynomials as the representation attached to the cuspidal representation by Definition~\ref{gal-galois-representations}, so the two are conjugate, and the local-global compatibility of the latter holds for the former.

% before: sha-hecke-ideal-spherical
The descent of a cuspidal representation to $U(n)$ at the right level needs spherical local components where the Galois representation is unramified. For a plain local component the local Langlands correspondence gives a sum of characters, so an unramified parameter forces unramified characters and a spherical vector.

% before: sha-hecke-ideal-above
At a place above $\ell$, the flag of an ordinary point determines the local component of the cuspidal representation. The characters of the flag give the characters of a principal series, the Weil-Deligne representation shows that the local component is a subquotient of it, and the condition $\psi_i\ne\psi_j\varepsilon$, with increasing weights, keeps the principal series irreducible.

% before: sha-hecke-ideal-twist-local
Twisting a point of the zero locus by a character of $\ell$-power order twists its automorphic representation by a character of the determinant. The conditions that make a representation contribute to the Hecke algebra are local, so each must survive such a twist: the local Langlands correspondence must commute with it, and spherical, principal series and typed local components must keep their form. A representation with a spherical twist, or with a monodromy-free parameter of the given inertial type, is good, so the twist law applies to it.

% before: sha-hecke-ideal-twist-global
A Galois character of odd order with $\beta(cgc^{-1})=\beta(g)^{-1}$ is the square of a character with the same property, and the Hecke character of the square root, composed with the determinant, is a character of the adelic unitary group. Multiplying the automorphic forms of a representation by it gives the twisted representation, with local components twisted by the Hecke character of $\beta$ at the split places.

% before: sha-hecke-ideal-plain
The windowed condition asks for plain local components whose characters lie in a window of width one. For a cuspidal representation, plainness gives the window: after a twist that makes the central character unitary, the local components are unitary and generic, and their exponents lie strictly between $-1/2$ and $1/2$. Plainness itself passes between conjugate places of a conjugate self-dual representation and between two representations with the same Galois representation.

% before: sha-hecke-ideal-cuspidal
The Galois representation of an automorphic representation of $U(n)$ is built from the cuspidal pieces of its base change, one summand for each segment of their Speh representations. When it is irreducible there is a single summand, so the base change itself is cuspidal and supplies the cuspidal representation required for classicality. Any irreducible representation with the same Frobenius eigenvalues is conjugate to it, by the Chebotarev density theorem, and so is attached to the same cuspidal representation.

% before: sha-hecke-ideal-jacquet
The windowed condition above $\ell$ needs the local components at $\ell$ of points that need not be of distinct type. An eigenvector of the $U$-operators with nonzero eigenvalues has a nonzero image in the Jacquet module, and a character of the torus through which that image survives makes the representation a subquotient of a principal series, by Frobenius reciprocity. The eigenvalues of the vector determine the characters of this principal series.

% before: sha-hecke-ideal-ordinary-point
At an arithmetic point of the Hecke algebra the ordinary eigenform lies in an automorphic representation of $U(n)$, and the local lemma, applied at each place above $\ell$, places its local components there in principal series. Local class field theory then ties the characters of these principal series to the ordinary characters of the point. No condition on the type of the point is needed.

% before: sha-hecke-ideal-triangular
At a point of the zero locus with an irreducible Galois representation, the ordinary relations hold as identities of matrices. If one element separates the diagonal characters, its eigenspaces give a basis in which these identities make the representation upper triangular. The projections onto the eigenspaces are polynomials in that element, by Lagrange interpolation.

% before: sha-hecke-ideal-flag
Condition ($\alpha$) and the plainness above $\ell$ start from a flag at the chosen place above each place $v\mid\ell$ of $L^+$. Irreducibility makes the trace form on the image nondegenerate, so the trace relations of the universal pair become identities of matrices. Since the weight is increasing, an element of inertia separates the ordinary characters.

% before: sha-hecke-ideal-conjugate
The windowed condition concerns every place above $\ell$, and the flag is known at one place above each place of $L^+$. A conjugate self-dual cuspidal representation makes its Galois representation polarized, so at the conjugate place the restriction to the decomposition group is a twisted dual of the restriction at the chosen place. Twisted duality reverses and inverts the diagonal characters, and it preserves the relations $\kappa_i=\kappa_j\varepsilon$ and their powers.

% before: sha-hecke-ideal-exponents
At the chosen place above $\ell$, the local components of the cuspidal representation are subquotients of a principal series whose characters are tied to the ordinary characters of the point. Through this tie, two linked characters of the principal series would give $\psi_i=\psi_j\varepsilon$. A relation $\psi_i^m=(\psi_j\varepsilon)^m$ would give two characters whose absolute values at a uniformizer differ by the factor $q_w$, which no window of width one contains.

% before: sha-hecke-ideal-tw-base
Property (7) needs the Hecke module as a free module over a power series ring in weight variables, supported on the zero locus of the Hecke ideal. These data do not depend on the auxiliary places, and every construction at an augmented level refers to them.

% before: sha-hecke-ideal-tw-iwahori
At the auxiliary places the Hecke module is studied through the vectors fixed by subgroups of the Iwahori subgroup. The subgroup cut out by the diagonal entries of order prime to $\ell$ leaves a quotient that is an $\ell$-group, whose group ring supplies the variables $y_{j,i}$ of the patching. The operators at these places are double coset operators of diagonal elements.

% before: sha-hecke-ideal-tw-torus
The double coset operators of the dominant diagonal elements multiply as the elements do, and they are invertible once the size of the residue field is. They therefore extend to an action of the whole diagonal torus, and the actions at different places commute. The operators $U$ and the diamond operators are values of these actions.

% before: sha-hecke-ideal-tw-diamond
Two facts about the level cut out by the elements of order prime to $\ell$ are used. The quotient of the Iwahori subgroup by it is a product of cyclic $\ell$-groups generated by diagonal units, which act through the torus action. The sum over this quotient carries the torus action at the smaller level to the torus action at the Iwahori level.

% before: sha-hecke-ideal-tw-dimension
The relations of the Taylor-Wiles data hold only up to nilpotence, and the exponent comes from a bound on the dimension of the fixed vectors in an irreducible representation. The bound follows from the results of Borel and Casselman on Jacquet modules and of Bernstein and Zelevinsky on principal series.

% before: sha-hecke-ideal-tw-local
At an auxiliary place, the eigenvalues of the torus action on the fixed vectors of a local component determine the Frobenius eigenvalues of the Galois representation, through the local Langlands correspondence. When these eigenvalues are distinct and no two are in ratio $q_w$, tame inertia acts on each Frobenius eigenspace through the character given by the torus action at the unit $\delta$. These local relations become the relations of the Taylor-Wiles data.

% before: sha-hecke-ideal-tw-level-module
The construction at an augmented level starts from a module of ordinary forms with an action of the groups at the auxiliary places. Its vectors fixed by the maximal compact subgroups are the module at the original level, and its points are the automorphic representations it contains, with their Galois representations. The local conditions record the local relations of Lemma~\ref{sha-hecke-ideal-tw-local} at all points.

% before: sha-hecke-ideal-tw-kill
A polynomial in the operators at an auxiliary place that vanishes at every point of the module is nilpotent. Its $n!$-th power vanishes, because at each point the operators act on the fixed vectors of an irreducible local component, a space of dimension at most $n!$.

% before: sha-hecke-ideal-tw-operator-module
The passage from automorphic forms to the Taylor-Wiles data goes through a module with commuting operators whose relations are stated in a polynomial ring. This intermediate object separates the representation theory at the auxiliary places from the commutative algebra that produces the patching data.

% before: sha-hecke-ideal-tw-comparison
The relation between the augmented level and the original level is a comparison of the Iwahori-fixed vectors at the auxiliary places with the module at the original level, tensored with the algebra that splits the characteristic polynomials of Frobenius. The comparison holds up to multiplication by the element $f$, which separates the roots and their multiples by $q_{w_j}$.

% before: sha-hecke-ideal-tw-operators
With the local conditions, a module at an augmented level gives an operator module. The diamond operators come from the quotient of the Iwahori subgroups by their pro-$\ell$ parts, and the relations come from the local relations at the points, with the exponent of Lemma~\ref{sha-hecke-ideal-tw-kill}.

% before: sha-hecke-ideal-tw-nilpotent
The patching data live on the part of a module where chosen operators act topologically nilpotently. This part is a direct summand, and on it the operators define an action of a power series ring.

% before: sha-hecke-ideal-tw-free
The freeness in property (7) is over the group ring of the diamond operators, a quotient of a power series ring in the variables $y_{j,i}$. A direct summand of a projective module over this local ring is free.

% before: sha-hecke-ideal-tw-quotient
Dividing by the diamond variables brings the smaller level back to the Iwahori level. On the topologically nilpotent parts, the sum over the diamond group identifies the two.

% before: sha-hecke-ideal-tw-splitting
The universal splitting algebra of the characteristic polynomials of Frobenius has the power series quotient $B[u]$ as a direct factor, cut out by the lifts of the residual eigenvalues. So comparison maps through the splitting algebra restrict to comparison maps through $B[u]$.

% before: sha-hecke-ideal-tw-patching
The commutative algebra of the last four lemmas turns an operator module with its relations into the data of property (7) at one level. Its ring is a power series ring over the ring of the operator module, in variables for the operators $U$ and the diamond operators.

% before: sha-hecke-ideal-tw-divided
The comparison of the Iwahori-fixed vectors with the spherical vectors at the auxiliary places runs through the Iwahori-Hecke algebra, whose commutative part is a polynomial ring in $n$ variables. Passing from that polynomial ring to its symmetric polynomials, which act on spherical vectors, uses the divided difference, an antisymmetrization divided by the Vandermonde product, and pairs of dual families, which express the polynomial ring over the symmetric polynomials.

% before: sha-hecke-ideal-tw-dual
Dual families exist in every rank, and the divided difference is linear over the symmetric polynomials. The two expansion identities let a polynomial in the Hecke operators be moved across a tensor product over the symmetric polynomials, which is how the comparison maps become linear.

% before: sha-hecke-ideal-tw-bernstein
The double coset operators of the diagonal elements and of the permutation matrices generate the Iwahori-Hecke algebra of $GL_n(F)$. The Bernstein presentation describes it through a polynomial ring in the diagonal operators and the finite Hecke algebra, and its basis shows that the polynomial ring acts faithfully on the sum of the operators of the permutations.

% before: sha-hecke-ideal-tw-identities
The comparison uses three identities in the Iwahori-Hecke algebra. The symmetric polynomials commute with the finite Hecke algebra, the products of the sum of the permutation operators with a pair of dual families add up to $N_q$, and that sum turns a polynomial on either side into a symmetric one, its divided difference after multiplication by $N_q$.

% before: sha-hecke-ideal-tw-trace
On Iwahori-fixed vectors the sum of the permutation operators is the trace to the spherical vectors. So the identities of the Hecke algebra hold for the operators $q^{i-1}\Theta(e_i(x))$ and this trace on any module without torsion by integers, and on spherical vectors the trace is multiplication by the index of the Iwahori subgroup, which can be cancelled.

% before: sha-hecke-ideal-tw-nil
The argument of Lemma~\ref{sha-hecke-ideal-tw-kill} applies to any polynomial in the operators $\Theta_j(t)$ and any level at the place $w_j$. A polynomial that vanishes at every character occurring in $\sigma$ acts nilpotently, with exponent $n!$, on the fixed vectors of a sum of copies of $\sigma$; it is needed here for the full Iwahori subgroup.

% before: sha-hecke-ideal-tw-scalar
On the spherical vectors of an irreducible representation, a symmetric polynomial in the operators $q^{i-1}\Theta(e_i(x))$ acts by a scalar, since these vectors form at most a line. When the polynomial takes one value at all the characters that occur, the scalar is that value, which ties the action of the symmetric functions to the Satake parameters.

% before: sha-hecke-ideal-tw-symmetric
At the points of a module at an augmented level, the local conditions identify the Satake parameters at $w_j$ with the Frobenius eigenvalues of the Galois representation. On the spherical vectors the elementary symmetric functions of the $U^w_{j,i}$ therefore act by the coefficients of the characteristic polynomial of the determinant at $\phi_j$.

% before: sha-hecke-ideal-tw-split
The comparison maps are built by commutative algebra from the traces at the places $w_j$ and the action of the symmetric functions on the spherical vectors. One map extends the inclusion of the spherical vectors by the operators $U_{j,i}$; the other applies the traces after multiplication by a dual family. Their composites are the products of the $N_{b_j}$, which divide $f$.

% before: sha-hecke-ideal-tw-augmented-idempotent
The Hecke modules below are localized at systems of eigenvalues of commuting operators. Over a complete local ring the localization is cut out by an idempotent of the algebra of the operators. So it is a direct summand, preserved by every endomorphism that commutes with the operators.

% before: sha-hecke-ideal-tw-augmented-fewer
At an augmented level there are fewer Hecke operators away from $\ell$, since the places of $S'\setminus S$ are left out. The localization at fewer operators agrees with the full one as soon as every eigensystem of the smaller family with the right congruences has them for the whole family.

% before: sha-hecke-ideal-tw-augmented-isotypic
A point of an augmented module needs a representation on which each group $GL_n(L_{w_j})$ acts through a single irreducible representation. A representation of a product of groups that is semisimple for each factor has quotients with this property, and they can be chosen to keep a given nonzero vector.

% before: sha-hecke-ideal-tw-augmented-projection
The specialization of an augmented module at a point of $\Lambda$ is spanned by eigenspaces for the action of the universal ring. A projection onto one eigenspace that commutes with the group action and keeps a given vector picks out the eigensystem of a single point.

% before: sha-hecke-ideal-tw-augmented-setup
Hida's ordinary forms at levels deepened at the places $w_j$ are functions on the adelic points of $U(n)$ with values in $\Lambda^m$. The setting fixes these forms, the levels, the Hecke operators and the localized modules $M_K$, whose sum $P$ carries the action of the groups at the $w_j$.

% before: sha-hecke-ideal-tw-augmented-parts
The points of $P$ come from classical automorphic forms. They are described through the points of $\Lambda$ at which the forms specialize, the spaces of classical forms in a single automorphic representation, and the pairs of a determinant and characters that the Hecke operators determine.

% before: sha-hecke-ideal-tw-augmented-torsion
The levels must be neat for the ordinary forms to be free. Neatness comes from the place $v_0$, where the level lies in the principal congruence subgroup: for an odd prime unramified in the field, this subgroup has no element of finite order other than $1$.

% before: sha-hecke-ideal-tw-augmented-finite-order
An element of $U(n)(L^+)$ that fixes a point of the adelic quotient is, place by place, conjugate into an integral group. Since the group of real points is compact, such an element has finite order, which is the other half of the argument for neatness.

% before: sha-hecke-ideal-tw-augmented-uop
The ordinary projector is built from the operator $U_z$, so its interaction with the groups $N_0$, $T_0$ and $\bar N_c$ at the places above $\ell$ is needed. These are computations with coset representatives, valid in any representation over any ring.

% before: sha-hecke-ideal-tw-augmented-forms
The operators $U_\alpha$ act on $\Lambda$-adic forms and preserve the ordinary ones, and the forms of a given level form a complete module. Repeated application of $U_z$ makes a form invariant modulo $\mathfrak m_\Lambda^k$ under a fixed group $\bar N_c$, which reduces questions modulo $\mathfrak m_\Lambda^k$ to finitely many values.

% before: sha-hecke-ideal-tw-augmented-sections
For a subgroup satisfying the freeness condition, a form is determined modulo $\mathfrak m_\Lambda^k$ by its values at finitely many points, and every choice of values at these points comes from a form. So the forms modulo $\mathfrak m_\Lambda^k$ are approximated by a finite free module.

% before: sha-hecke-ideal-tw-augmented-projector
Hida's finiteness theorem has an algebraic core. An operator on a complete filtered module over a complete local ring, approximated modulo each step of the filtration by maps through a finite free module, has a finite free ordinary part, and the limit of its factorial powers is a projector onto it.

% before: sha-hecke-ideal-tw-augmented-ordinary
The $\Lambda$-adic forms fixed by a subgroup of finite level satisfying the freeness condition meet the hypotheses of the algebraic statement, with the operator $U_z$. So their ordinary part is finite and free over $\Lambda$, with an ordinary projector.

% before: sha-hecke-ideal-tw-augmented-spherical
The Hecke operators at a split place and at its conjugate act through the same group $U(n)(L^+_v)$. They commute because the spherical Hecke operators of $GL_n$ commute, by Gelfand's argument with the transpose.

% before: sha-hecke-ideal-tw-augmented-hida-classes
The finite level of $\mathcal K_0K$ comes from the finiteness of class numbers of the definite unitary group $U(n)$. The form needed allows any open subgroup of the finite adelic points, with the real points, which form a compact group, absorbed into the double cosets.

% before: sha-hecke-ideal-tw-augmented-hida-weight
The specialization at an arithmetic point compares $\Lambda$-adic forms with vector-valued forms of an algebraic weight. Of the algebraic representation only a few properties are used: integrality on the Iwahori subgroup and, after a twist by a torus character, on the dominant elements, a highest weight vector in the coordinate $0$, and continuity. The control map reads off the coordinate $0$ after untwisting a form by the representation.

% before: sha-hecke-ideal-tw-augmented-hida-classical
A classical form of a fixed level is determined by its values at finitely many points, and its values elsewhere are integral combinations of these. This gives the finite dimension of the classical forms and the integrality of the eigenvalues of the operators $U'_\alpha$ and of sums of translations.

% before: sha-hecke-ideal-tw-augmented-hida-base-change
An arithmetic point with values in $\mathcal O'$ specializes the ordinary $\Lambda$-adic forms to ordinary $\mathcal O'$-adic forms. The ordinary forms over $\Lambda$ and over $\mathcal O'$ are both approximated, modulo the maximal ideal, by the ordinary projections of the same finitely many sections, so by Nakayama's lemma the specialization is a base change.

% before: sha-hecke-ideal-tw-augmented-hida-inputs-separation
The local statement at the places above $\ell$ passes from a space of forms to an irreducible representation $\sigma$ by equivariant maps. When the space is a sum of images of $\sigma$, such maps detect every nonzero element, since the space is then a direct sum of copies of $\sigma$.

% before: sha-hecke-ideal-tw-augmented-hida-inputs-embedding
The weight datum of a given weight is defined over a ring of integers that contains the images of the local rings $\mathcal O_{L_{\tilde v}}$ under the embeddings into $\bar{\mathbf Q}_\ell$. The contraction of the coordinates other than the coordinate $0$ by $\chi(z)^{-1}\xi(z)$ also needs the uniformizers to map into the maximal ideal.

% before: sha-hecke-ideal-tw-augmented-hida-inputs-antidiagonal
The lowest weight vector of the dual Weyl module is singled out by evaluating bideterminants at the antidiagonal matrix. This evaluation transforms under upper triangular matrices by a character, and it vanishes on the bideterminants of all but one shape of rows.

% before: sha-hecke-ideal-tw-augmented-hida-inputs-dominance
The weights of the bideterminants dominate the lowest weight, for every nonincreasing pairing. Equality for the pairing with $(n-1,\dots,1,0)$ singles out the bideterminant whose rows are the last ones, which is the lowest weight vector.

% before: sha-hecke-ideal-tw-augmented-hida-inputs-tableaux
The integral structure of the algebraic representation comes from the standard basis of bideterminants. In this basis the congruence subgroups act by matrices congruent to $1$, the diagonal torus acts diagonally, and the coordinate of the lowest weight vector is a character of the upper triangular matrices.

% before: sha-hecke-ideal-tw-augmented-hida-inputs-matrix
The conditions of a weight datum are conditions on the matrix coefficients of the representation in one basis. Ordering the basis so that the lowest weight vector comes first turns such conditions into a weight datum, with the torus character read off at that vector.

% before: sha-hecke-ideal-tw-augmented-hida-inputs-lattices
The comparison through the control map keeps track of integrality by untwisting a classical form by the algebraic representation. The untwisted values of a form lying in a fixed $\mathcal O'$-module define the lattices used in the comparison.

% before: sha-hecke-ideal-tw-augmented-hida-inputs-lattice
The operator $U'_z$ contracts the coordinates other than the coordinate $0$ by the maximal ideal. So if the untwisted values of a form lie in a module $\Lambda$ and its control map lies in a module containing $\mathfrak m'\Lambda$, the untwisted values of its image under $U'_z$ lie in the second module.

% before: sha-hecke-ideal-tw-augmented-hida-inputs-levels
Raising the level at $\ell$ by one step is undone by one application of $U'_z$. This lets classical forms of deeper level be brought back to the level $c$ without leaving the classical forms.

% before: sha-hecke-ideal-tw-augmented-hida-inputs-values
A classical form of level $c$ is determined by its untwisted values at the representatives of finitely many double cosets, and every family of values occurs. The integrality of these values gives the integrality of all the untwisted values.

% before: sha-hecke-ideal-tw-augmented-hida-inputs-lift
The control map of an integral classical form satisfies the conditions defining $\mathcal O'$-adic forms. The weight datum supplies these conditions through the row $0$ of the representation and its congruences on the groups $\bar N_c$.

% before: sha-hecke-ideal-tw-augmented-hida-inputs-fitting-projector
The ordinary part of an endomorphism of a finite free module over a complete local ring is cut out by an idempotent. It is built modulo each power of the maximal ideal and then passed to the limit, so the conditions it satisfies are stated for an arbitrary ideal.

% before: sha-hecke-ideal-tw-augmented-hida-inputs-fitting-naturality
An ordinary projector is determined by the endomorphism, and maps that intertwine two endomorphisms intertwine their ordinary projectors. This makes the projectors modulo the powers of the maximal ideal compatible, and identifies the image of a projector with the intersection of the images of the powers of the endomorphism.

% before: sha-hecke-ideal-tw-augmented-hida-inputs-fitting-length
Modulo a power of the maximal ideal, the module has finite length. There Fitting's decomposition into the kernel and the image of a high power of the endomorphism gives an ordinary projector.

% before: sha-hecke-ideal-tw-augmented-hida-inputs-fitting-factorial
On a finite module the powers of an endomorphism repeat, so its factorial powers are eventually constant and idempotent. This identifies the ordinary projector modulo each power of the maximal ideal with a factorial power of the endomorphism.

% before: sha-hecke-ideal-tw-augmented-hida-inputs-fitting
On a finite free module over a complete local ring, an endomorphism splits off the part on which it acts bijectively. After extending scalars to an algebraically closed field this part spans the generalized eigenspaces whose eigenvalues are units integral over the ring.

% before: sha-hecke-ideal-tw-augmented-hida-inputs-ordinary-lattice
Through the bijection of the classical forms with their values at finitely many points, the ordinary part of $U'_z$ acquires a free integral structure. Its elements are periodic for $U'_z$ modulo every power of the maximal ideal, and every integral classical form is approximated by them under the powers $U_z'^{j!}$.

% before: sha-hecke-ideal-tw-augmented-hida-inputs-automorphic-sum
The decomposition of classical forms into automorphic parts starts from the automorphic forms on the definite unitary group. Its adelic quotient is compact, so these forms decompose into automorphic representations.

% before: sha-hecke-ideal-tw-augmented-hida-inputs-automorphic-projections
The decomposition is used through projections. A finite-dimensional space of automorphic forms is split by finitely many projections onto automorphic representations, which commute with right translation and so preserve every condition defined by translations.

% before: sha-hecke-ideal-tw-augmented-hida-inputs-automorphic-forms
A classical form of weight $\xi$ is a function on the finite adelic points with values in the algebraic representation. Untwisting it by the representation at the infinite places and pairing with linear forms turns it into automorphic forms on the adelic points, and every compatible family of automorphic forms arises in this way.

% before: sha-hecke-ideal-tw-augmented-hida-inputs-automorphic-near
The level of a classical form and its invariance under an Iwahori group at $\ell$ are conditions at finitely many places. So the elements of the finite adelic points close to $1$ act on such a form trivially, which makes the attached automorphic forms invariant under an open subgroup.

% before: sha-hecke-ideal-tw-augmented-hida-inputs-automorphic-local
The data of an automorphic part at the split places come from the local components of an automorphic representation. Restricted to the local group at a split place, the representation is a sum of copies of its local component, and at an unramified place the Hecke operators act on the spherical vectors by scalars given by the Satake parameters.

% before: sha-hecke-ideal-tw-augmented-hida-inputs-conjugate-dual
A classical form gives a map into an automorphic representation from the dual of an algebraic representation, while weights are defined through the algebraic representations themselves. On the compact unitary group, complex conjugation of the matrix entries is the inverse transpose, and it relates the dual of $\nabla(\lambda)$ to $\nabla(\lambda^\vee)$ in two ways: through the cofactor substitution, and through the Fock pairing.

% before: sha-hecke-ideal-tw-augmented-hida-inputs-automorphic-weight
At the places above $\ell$, the Hodge-Tate multisets of an automorphic part come from the weight at infinity of its automorphic representation. A nonzero classical form of weight $\xi$ in that representation gives a map from the dual of the algebraic representation into it, equivariant at infinity, and so fixes the weight at each embedding.

% before: sha-hecke-ideal-tw-augmented-hida-inputs
Four further inputs enter the specialization and the decomposition. They are the algebraic representation of a given weight with its integral structure, the comparison of classical and $\mathcal O'$-adic ordinary forms through the control map, the decomposition of classical forms into automorphic parts, and a statement on ordinary vectors in principal series at the places above $\ell$.

% before: sha-hecke-ideal-tw-augmented-hida-control
At an arithmetic point, an ordinary $\Lambda$-adic form has at most one classical ordinary form with the same control map at the levels of a directed family. This defines the specialization map, and the comparison through the control map makes its image span the classical ordinary forms, with the operators $U_\alpha$ and the sums of translations carried along.

% before: sha-hecke-ideal-tw-augmented-hida-decomposition
A classical form of weight $\xi$ lies in the sum of the automorphic parts, but the summands need not be classical of the same level and type. Averaging over the compact groups at the places $w_j$ and above $\ell$, against the eigencharacters of the form, moves the summands into the classical forms, and the ordinary part follows by taking generalized eigenspaces of $U'_z$.

% before: sha-hecke-ideal-tw-augmented-hida-eigen
The Hecke algebra of a localized module acts on the span of the specializations of the module, once the specializations of a basis are linearly independent. The action is the base change of the action on the module, and a decomposition of the span of the specializations of the whole ordinary module passes to the smaller span through the idempotent of the localization.

% before: sha-hecke-ideal-tw-augmented-hida
The construction at the augmented levels uses five inputs from Hida theory on $U(n)$, at every level $\mathcal K_0K$. The levels are neat; the ordinary forms are finite and free with commuting operators and have an ordinary projector; and at an arithmetic point they specialize onto classical forms that decompose into automorphic parts.

% before: sha-hecke-ideal-tw-augmented-projective
Patching needs the invariants of a deeper level to be projective over the group ring of the quotient of the two levels. At a neat level this follows from a fundamental domain for the action of the quotient, together with the ordinary projector. The same splitting identifies the ordinary forms of the larger level with the traces from the deeper one.

% before: sha-hecke-ideal-tw-augmented-sum
The localized modules $M_K$ grow as the levels $K$ shrink, and right translation by $G_w$ carries one into another. Their union $P$ is therefore a representation of $G_w$, whose vectors fixed by each $K$ form $M_K$, and these are projective over the finite quotients of levels.

% before: sha-hecke-ideal-tw-augmented-unique
A Hecke algebra carries at most one pair of a determinant and characters attached to its operators. This makes the maps from the universal ring compatible when the level changes and when operators are left out.

% before: sha-hecke-ideal-tw-augmented-compatible
An injective map between localized modules that commutes with the operators restricts the larger Hecke algebra to the smaller one. By the uniqueness of attached pairs, the maps from the universal ring are compatible with this restriction.

% before: sha-hecke-ideal-tw-augmented-action
The maps from the universal ring of the augmented level to the Hecke algebras of the $M_K$ glue to an action on the union $P$. This action commutes with $G_w$, because translation by $G_w$ maps one localized module into another.

% before: sha-hecke-ideal-tw-augmented-eigensystems
At an arithmetic point of $\Lambda$, every eigensystem of the Hecke algebra of the ordinary forms occurs in an automorphic part. This attaches a Galois representation to the eigensystem, with Frobenius characteristic polynomials given by its Hecke eigenvalues.

% before: sha-hecke-ideal-tw-augmented-base
At the level $K_w$ the localized module is the module $M$ of the original level. Its localization at the operators outside $S'$ agrees with the one outside $S$, and on it the universal ring of the augmented level acts through $\beta$.

% before: sha-hecke-ideal-tw-augmented-points
Every nonzero element of $P$ is detected by a point of the module. Its specialization at a suitable point of $\Lambda$, projected onto one eigensystem of the universal ring and reduced to single local components at the $w_j$, gives the representation, the Galois representation and the local components that a point requires.

% before: sha-hecke-ideal-tw-augmented
The modules at the augmented levels come from Hida's ordinary forms on $U(n)$ at levels that are deeper at the places $w_j$. Their sum over these levels carries the action of the groups at the $w_j$, the universal ring of the augmented level acts on it through the Hecke algebras, and its spherical part at the $w_j$ is the module $M$ of the original level.

% before: sha-hecke-ideal
The Hecke ideal is the kernel of the map from the universal ring to the Hecke algebra. Its zero locus contains the automorphic points of the universal ring, and its irreducible arithmetic points are automorphic, so a point is automorphic once it is shown to lie on this zero locus. The same construction at levels augmented by auxiliary places gives the Taylor-Wiles data used to bound the universal ring.

% before: sha-patching-normalization
The patching argument works with arcs in a normalized form: the symmetric form is the identity, the representation is block diagonal, and the inertia at the auxiliary places acts trivially. Each nice arc can be brought to this form after a finite extension of its coefficients, without changing its determinant.

% before: sha-patching-frobenius-cocycles
A Selmer basis and a finite detecting set both come from finite generation. The Selmer group and the dual cocycles are modules of continuous cocycles with values in a free module of finite rank over $\kappa[[T]]$. Such a module is finitely generated when the group has finitely many continuous homomorphisms to $\mathbf Z/\ell$ on each open subgroup, by reduction modulo $T$ and Nakayama's lemma.

% before: sha-patching-frobenius-lattice
A Selmer basis is a family of Selmer classes that is independent modulo coboundaries and spans the Selmer group modulo coboundaries up to powers of $T$. Over a discrete valuation ring such a family lifts a basis of the free part of the quotient.

% before: sha-patching-frobenius-selmer
With these two lemmas every normalized arc has a Selmer basis, as soon as the Galois group satisfies the finiteness condition. The Selmer group and its coboundaries are submodules, and the Selmer group consists of continuous cocycles.

% before: sha-patching-frobenius-forms
The pairings with the elements of a set $\Sigma$ are linear forms on the dual cocycles. On a finitely generated module over a noetherian ring, finitely many of a family of linear forms already cut out the common kernel of all of them.

% before: sha-patching-frobenius-finite
The density argument shows that the set of elements of $\Gamma'$ with separable characteristic polynomials detects the dual Selmer group. Taylor-Wiles places are chosen for finitely many elements, so a finite detecting subset is needed.

% before: sha-patching-frobenius-extension
The eigenvalues of the chosen elements must lie in the coefficient ring of the arc, and the residue fields of all arcs must have the same size. Both are reached by a finite extension of $\kappa[[T]]$, and the Selmer basis and the detection survive the extension, through the coordinates of the extension as a free module.

% before: sha-patching-frobenius
Taylor-Wiles places are chosen with Frobenius elements close to finitely many elements $\sigma_j$ of $\Gamma'$. The $\sigma_j$ must have distinct eigenvalues on each arc and must detect the dual Selmer group, so that the added places kill the dual Selmer group in the limit.

% before: sha-patching-places
The Chebotarev density theorem gives places whose Frobenius elements have inverses approximating the $\sigma_j$ and whose norms are congruent to $1$ modulo $\ell^N$. The patched rings and modules also need the relation between Frobenius and tame inertia at these places, and control of the new ramification.

% before: sha-patching-tangent-modules
The tangent bound counts Selmer classes of the first-order deformation $\rho_1$ modulo $T^m$ at the augmented levels. They take values in a module of matrices over the reduction of $R_g$, with local conditions at the places above $\ell$ given by the triangular form of $\rho_1$, and the count compares them with cocycles of the twisted dual module $V^*$.

% before: sha-patching-tangent-counts
The count rests on published results in Galois cohomology. The Greenberg-Wiles formula relates a Selmer group to the dual Selmer group, Tate's local formulas give the sizes of local cocycle groups, and local class field theory describes the abelian quotients of the decomposition groups above $\ell$.

% before: sha-patching-tangent-involution
The Selmer group is defined by cocycles of $\Gamma_{N,0}$ with a symmetry under $c$, while the Greenberg-Wiles formula concerns cocycles of $\Gamma_N$. Since $c$ has order $2$ and $2$ is invertible, the two are related by a count of cocycles vanishing at $c$.

% before: sha-patching-tangent-symmetry
A symmetric cocycle of $\Gamma_{N,0}$ is the restriction of a unique cocycle of $\Gamma_N$ vanishing at $c$. This identifies the Selmer group, and the dual cocycles, with cocycles of $\Gamma_N$.

% before: sha-patching-tangent-selmer
With these identifications the Selmer group of $\rho_1$ and its coboundaries are bounded by cocycles of $\Gamma_N$ with local conditions, up to the fixed points of $c$ and of $\Gamma_N$. The local conditions are the ordinary cocycles at the places above $\ell$, the cocycles vanishing on $I'_u$ at the auxiliary places, and no condition at $v_0$ and at the Taylor-Wiles places.

% before: sha-patching-tangent-fixed
Several terms of the count are invariants of a group on a module modulo $T^m$. Their number grows like $|\kappa|^{mr}$, where $r$ is the rank of the invariants over $\kappa[[T]]$, by the Artin-Rees lemma.

% before: sha-patching-tangent-invariants
The global invariants and the local invariants at $v_0$ have rank $0$, by the irreducibility built into a normalized arc and by (N4). At a Taylor-Wiles place the invariants are a centralizer of rank at most $n$, which gives the term $nq$ of the exponent.

% before: sha-patching-tangent-trace
The local counts and the Greenberg-Wiles formula need a pairing on $A_m$ with values in $\mathbf Z/\ell$ whose kernels are trivial. The trace to $\mathbf F_\ell$ of the coefficient of $T^{m-1}$ gives one, and it is compatible with the action of $\mathbf Z_\ell$ through the cyclotomic character.

% before: sha-patching-tangent-local
The local conditions at the places above $\ell$ are the ordinary cocycles. After conjugating to triangular form they are filtered by height, with homomorphisms on the diagonal and cocycles for the quotients of the diagonal characters above it.

% before: sha-patching-tangent-skew
The complex conjugations enter the Greenberg-Wiles formula through their fixed points. Each fixes at least as many elements as there are antisymmetric matrices in $W_m$, which is the source of the term $m\nu d$ that cancels at the end.

% before: sha-patching-tangent-solutions
Detection is a statement over $\kappa[[T]]$, while the counts are made modulo $T^m$ for every $m$. The passage between them is a fact of commutative algebra: if every solution of a linear system over a complete discrete valuation ring has a multiple in a given submodule, then modulo $t^m$ the solutions exceed the reductions of that submodule by a factor bounded independently of $m$.

% before: sha-patching-tangent-detection
Applied to dual cocycles, the previous lemma turns detection over $\kappa[[T]]$ into a bound modulo $T^m$ that is uniform in $m$. The linear system records that a map takes values in $V^*$, the cocycle relation, the symmetry under $c$ and the vanishing of the pairings with the $\sigma_j$.

% before: sha-patching-tangent-descent
The bound of the previous lemma concerns cocycles of $\Gamma_0$, and the Greenberg-Wiles count produces cocycles of $\Gamma_N$. A cocycle of $\Gamma_N$ vanishing at $c$ and at the tame generators $\tau_{N,j}$ vanishes on the kernel of $\Gamma_N\to\Gamma$, since that kernel is generated by the conjugates of the $\tau_{N,j}$ and by $P(H_N)$, so it comes from $\Gamma_0$.

% before: sha-patching-tangent-dual
The dual term is where the Taylor-Wiles places act. A dual cocycle that is a coboundary at the places $w_{N,j}$ descends, for large $N$, to a dual cocycle of $\Gamma_0$ whose pairings with the $\sigma_j$ vanish, and detection bounds these.

% before: sha-patching-tangent
The tangent space of the deformation problem at an arc is bounded by the Greenberg-Wiles formula. With the dual Selmer group killed by the added places, the bound has exponent $n[L^+:\mathbf Q]+nq$ in $m$, the number of variables of the patched power series ring.

% before: sha-patching-obstruction-count
The obstruction group is a second cohomology group with local conditions. An exact sequence relates it to the Selmer group, the local first cohomology modulo the local conditions and the cokernel of global to local second cohomology. When the global and local Euler characteristic formulas bound the first two, the sequence bounds its size by the Selmer group, the cokernel and the local factors $r_w$.

% before: sha-patching-obstruction
A bound on the tangent space alone does not make the deformation ring a power series ring near the arc; the obstructions to lifting first-order deformations must also be bounded. Duality bounds the relevant second cohomology group by the rank $g$ of the Selmer group, up to a power of $T$ independent of the level.

% before: sha-patching-chain-patched
The patching method turns modules over the rings of the augmented levels, free over quotients of a power series ring, into information about the support of the module at the base level near chosen primes. It is stated here in the form used in Alpöge's proof: near each prime the support of $X$ on the locus where the patching variables vanish is cut out by an ideal all of whose components have the dimension of the power series ring.

% before: sha-patching-chain-limit
To patch, the Galois groups of the augmented levels must be compared with each other along a single profinite group. A free profinite group covers all of them, since the groups $H_N$ need a bounded number of generators, and a limit along an ultrafilter gives a limit level with limits of the Frobenius and tame elements.

% before: sha-patching-chain-ring
The ring on which the patched modules live is a deformation ring for the limit level, enlarged by the patching variables and by elements that play the role of the eigenvalues at the Frobenius elements. Its relations are imposed later, as an ideal, so that the ring itself has a universal property.

% before: sha-patching-chain-maps
The ring of patching tuples maps to each ring of the Taylor-Wiles data and to the base ring with its eigenvalues adjoined. These maps carry the patching variables to the variables of the Taylor-Wiles data and to zero.

% before: sha-patching-chain-relations
The relations of an ordinary pair, together with the trace relations at the limit Frobenius and tame elements, cut out an ideal that the map to the base ring kills. Modulo this ideal the universal determinant descends to the limit level.

% before: sha-patching-chain-hypotheses
The hypotheses of the patching lemma hold for the ring of patching tuples. Each comes from the corresponding property of the Taylor-Wiles data at the levels $N$, passed to the limit along the ultrafilter.

% before: sha-patching-chain-points
The patching lemma needs primes of the ring of patching tuples at which the comparison element $f$ is a unit. The arcs give such primes, because their eigenvalues at the $\sigma_j$ are distinct.

% before: sha-patching-chain-completion
The patching lemma works in completions of the ring of patching tuples at the primes below the arcs. Near an arc such a completion has the fraction field of the arc as residue field, and it is an algebra over a complete discrete valuation ring with that residue field, built from the weight variable that the arc sends to $T$.

% before: sha-patching-chain-chart
The bound at an arc ends with a dimension count on a determinantal locus: a matrix $z$ and a vector $m$ with $zm=0$. When the entries are part of a regular system of parameters, the primes containing the relations and $\det z$ all contain one prime, the generic kernel read off on affine charts.

% before: sha-patching-chain-incidence
The relations $zm=0$ cut out a locus of codimension at least the length of $m$. So a prime of large dimension on which $m$ does not vanish is the generic kernel itself, and this passes along a surjection onto an integral extension.

% before: sha-patching-chain-square-zero
At an arc of type E, the representation near the arc lives over an extension of the completion by a module whose products fall into the maximal ideal. Such an extension is again a complete local ring with the same residue field.

% before: sha-patching-chain-cotangent
The bounds at an arc give a presentation of the completion and of the module $N$ up to the ideal $J$: few generators, and $g^2$ relations whose coefficients are independent modulo the uniformizer. Lifting the presentation to a power series ring puts the determinantal locus of the previous lemmas inside the completion.

% before: sha-patching-chain-stable
Lemma~\ref{sha-patching-chain-cotangent} needs the products of elements of $N$ to stay outside the prime $Z$. At a stable prime the Gram determinants of the traces do not all vanish, and the span of the representation over $B\oplus N$ turns this into a product of elements of $N$ outside the prime.

% before: sha-patching-chain-datum
The bound at an arc is proved over a quotient of the patched ring, with the limit group $\Gamma_\infty$ in place of $\Gamma$. The arc enters through a homomorphism to $\kappa[[T]]$ together with the Taylor-Wiles elements and the relations they satisfy. The construction of the representation over the completion uses only these data.

% before: sha-patching-chain-cayley-factor
The arc itself must satisfy the relations of its arc datum. Each generator of the relation ideal is a Cayley-Hamilton or ordinary relation that holds for the matrices of the arc, and the closure in the definition is harmless because the arc is continuous.

% before: sha-patching-chain-cayley-closure
The algebra of an arc datum is a quotient of a group ring by a closed ideal, and the algebra $E_0$ that defines thick primes is built in the same way. A determinant descends to such a quotient through the elements it kills, and the closure over powers of the maximal ideal and open subgroups keeps the quotient continuous.

% before: sha-patching-chain-cayley-extension-bound
The finiteness of these quotients comes from a bound on their dimension over the residue field. That dimension is controlled by the spaces of cocycles linking two residual characters, which describe the off-diagonal parts of the algebra.

% before: sha-patching-chain-cayley-bound
The cocycle spaces are finite because the groups involved have few homomorphisms to $\mathbf Z/\ell$. This holds for topologically finitely generated groups, such as the limit group of an arc datum.

% before: sha-patching-chain-cayley-lift
Modulo its radical, a Cayley-Hamilton algebra with distinct residual characters is a product of copies of the residue field. Its coordinate idempotents lift to the algebra, by Hensel's lemma in the finite commutative subalgebra generated by one element.

% before: sha-patching-chain-cayley-gma
With the lifted idempotents, such an algebra has a decomposition like that of the matrix algebra of size $n$, with diagonal idempotents of rank one. Its radical is the kernel of the reduction to the diagonal.

% before: sha-patching-chain-cayley-trace
In such an algebra the kernel of the determinant is detected by traces. Each corner cut out by an idempotent is a line, so an element whose traces against all elements vanish is killed.

% before: sha-patching-chain-cayley-chars
Over the residue field the determinant is that of a sum of characters. The projection to the diagonal describes both its kernel and its characteristic polynomials.

% before: sha-patching-chain-cayley-dimension
A finite-dimensional algebra with such idempotents has its radical bounded by the cocycle spaces. The off-diagonal corners of a linear form vanishing on the square of the radical are cocycles, and the radical is nilpotent of order at most $n$.

% before: sha-patching-chain-cayley-field
Over the residue field, the group ring modulo the Cayley-Hamilton ideal and a small open subgroup is finite-dimensional. The bound on its dimension does not depend on the open subgroup, which is what the passage to the limit uses.

% before: sha-patching-chain-cayley-level
The spanning set over the residue field lifts to the coefficient ring by Nakayama's lemma. The Cayley-Hamilton ideal of the reduction is generated by the reduction of the Cayley-Hamilton ideal, so the relations lift as well.

% before: sha-patching-chain-cayley-limit
A bound on the number of spanning elements that does not depend on the open subgroup passes to the closure. The quotients modulo powers of the maximal ideal are finite, so the open subgroups stabilize at each level, and the completeness of the coefficient ring gives the limit.

% before: sha-patching-chain-cayley-finite
Together these give the finiteness of the quotient algebra for any set of relations, with the images of the group elements varying continuously. The same statement gives the finite spanning sets needed for thick primes.

% before: sha-patching-chain-cayley-descent
A continuous determinant kills the closure of any set it kills, as one sees by reduction modulo powers of the maximal ideal. So it descends to the quotient by the closure, where it is Cayley-Hamilton.

% before: sha-patching-chain-cayley-kills-trace
For a determinant whose residual characters have an extension bound, being killed is a condition on traces. The relations of an ordinary pair and of an arc datum are stated as conditions on traces, so this is how they are seen to be killed.

% before: sha-patching-chain-cayley-star
The polarization of an arc gives an anti-involution of the group ring. It preserves the determinant, its kernel and its Cayley-Hamilton ideal, and it preserves closures because the cyclotomic character is close to $1$ on small open subgroups.

% before: sha-patching-chain-cayley-kills
The relations of an arc datum are killed by its determinant and stable under the involution. So the relation ideal is the closure of the Cayley-Hamilton ideal and a set of killed relations, which is the form the previous lemmas need.

% before: sha-patching-chain-cayley
The representation of Lemma~\ref{sha-patching-chain-cotangent} is built from the algebra generated by the group modulo the relations of an arc datum. After completion at the arc, this algebra is finite, carries a Cayley-Hamilton determinant and an involution from the polarization, and maps to the matrices over $\kappa((T))$ through the arc.

% before: sha-patching-chain-idempotents-lift
An element that is idempotent modulo a nilpotent ideal of a complete algebra lifts to an idempotent. The lift is a limit of polynomials in the element, which is what makes it compatible with an involution and with a given idempotent in the next lemma.

% before: sha-patching-chain-idempotents-orthogonal
The matrix units of the residual representation are lifted one at a time. Each lift is taken inside the complement of the previous ones and symmetrized with the involution, which needs $2$ to be a unit.

% before: sha-patching-chain-idempotents-corner
A lifted idempotent of rank one cuts out a corner of the algebra. The Cayley-Hamilton property, applied over the dual numbers, shows that this corner is a line spanned by the idempotent, with the trace as coordinate.

% before: sha-patching-chain-idempotents
To turn the algebra into matrices one needs idempotents and units lifting the matrix units of the residual representation. Completeness lifts idempotents, and the square roots needed to normalize the units exist because $2$ is invertible.

% before: sha-patching-chain-gma
With these idempotents and units, the algebra embeds into matrices over a commutative ring. For two blocks the ring is the base ring plus one off-diagonal corner, with the product of two corner elements given by a trace pairing; this pairing vanishes residually, since the residual representation is block diagonal.

% before: sha-patching-chain-semisimple
Thick primes are defined by the failure of semisimplicity. Away from the annihilator of the off-diagonal part the algebra of the last lemma becomes a product of matrix algebras, so such primes contain that annihilator.

% before: sha-patching-chain-representation-determinant
The determinant of the representation near the arc is read off from the idempotents. Amitsur's formula reduces the comparison to the corners $e_iEe_j$: a diagonal corner is a line by Lemma~\ref{sha-patching-chain-idempotents-corner}, and an off-diagonal one has the characteristic polynomial of $0$.

% before: sha-patching-chain-representation
Putting the last four lemmas together gives the representation used in Step 4 of Lemma~\ref{sha-patching-chain}. Its span over the completion has the block shape that Lemma~\ref{sha-patching-chain-stable} needs, and its reduction is the arc.

% before: sha-patching-chain-selmer
The presentation of the completion comes from Selmer groups of the arc over the limit group. They are needed both over $\kappa[[T]]$, where the bounds at the levels $N$ pass to the limit, and over its fraction field, where the tangent directions of the completion give cocycles with bounded denominators.

% before: sha-patching-chain-limit-bounds-kernel
At each level the kernel of the map to $\Gamma$ is generated by tame inertia at the Taylor-Wiles places. This passes to the limit group, which is what lets odd Selmer classes for $\Gamma_\infty$ descend to $\Gamma$, where the Selmer basis spans them.

% before: sha-patching-chain-limit-bounds-transfer
Modulo a fixed power of $T$, Selmer sets are finite sets of locally constant cocycles. Along the patching ultrafilter, such a cocycle factors through the limit group exactly when it factors through the groups at the levels $N$, and the local conditions match because the local maps converge.

% before: sha-patching-chain-limit-bounds-artin-rees
The passage from bounds modulo powers of $T$ to a rank bound over $\kappa[[T]]$ needs a uniform loss of precision. The Artin-Rees lemma gives it for a family of linear maps between finite modules.

% before: sha-patching-chain-limit-bounds-closed
Second-order classes over $\kappa[[T]]$ are built from compatible classes modulo powers of $T$. The conditions other than the local one pass to the limit coefficient by coefficient, and the matrices of the local condition converge along an ultrafilter, as the residue field is finite.

% before: sha-patching-chain-limit-bounds-ultrafilter
Classes modulo increasing powers of $T$ need not be compatible. Since the Selmer sets modulo a fixed power are finite, an ultrafilter selects compatible reductions, and the last lemma turns them into a class over $\kappa[[T]]$.

% before: sha-patching-chain-limit-bounds
The bounds of Lemmas~\ref{sha-patching-tangent} and~\ref{sha-patching-obstruction} hold at each level $N$ for large $N$. The patching ultrafilter turns them into bounds for the limit group: a rank bound for first-order classes, the lifting of even classes to second order, and the generation of the odd Selmer group by the Selmer basis.

% before: sha-patching-chain-laurent-scaling
Bounded denominators let a single power of $T$ move a Selmer cocycle over the fraction field to one over $\kappa[[T]]$. The relations that define the Selmer groups are linear in the cocycle, with coefficients over $\kappa[[T]]$, so they survive this scaling.

% before: sha-patching-chain-laurent
The cocycles that come from the completion have values over the fraction field of $\kappa[[T]]$. Since their denominators are bounded, a power of $T$ brings them back over $\kappa[[T]]$, where the bounds of the last lemma apply.

% before: sha-patching-chain-tangent
The cotangent space of the completion is compared with Selmer groups through derivations at the arc. Additive maps on the off-diagonal part play the same role for the odd directions, and the functionals that kill $J$ detect the relations needed in Lemma~\ref{sha-patching-chain-cotangent}.

% before: sha-patching-chain-duality
The spaces of the last definition are dual to quotients of the completion and of its off-diagonal part. These dualities turn statements about all tangent directions into statements about elements, which is how the elements $X_k$ and $m_i$ are produced.

% before: sha-patching-chain-cocycles-ordinary
The local condition at the places above $\ell$ comes from the ordinary relations. Over a square-zero thickening, such as the dual numbers, Lagrange idempotents of a separating element of the group algebra put a deformation in triangular form, and moving them back to those of the reduction gives the coboundary that corrects the cocycle.

% before: sha-patching-chain-cocycles-selmer
A deformation of the arc over a ring of the shape $R_g$, with the relations of the completion, gives Selmer cocycles by taking coefficients. The coefficients of $\epsilon'_k$ give even cocycles and those of $\epsilon_k$ odd ones, and the local condition is the previous lemma.

% before: sha-patching-chain-cocycles-bounded
The cocycles coming from the completion must have bounded denominators. The values of the representation, pushed to a finite algebra over the fraction field, lie in a finitely generated module over $\kappa[[T]]$, and the continuity of the representation makes them locally constant modulo powers of $T$.

% before: sha-patching-chain-cocycles-test
Tangent directions and additive maps on the off-diagonal part combine into a homomorphism from $B\oplus N$ to a ring of the shape $R_g$. Composed with the representation it is a deformation of the arc with all the relations of Lemma~\ref{sha-patching-chain-cocycles-selmer}.

% before: sha-patching-chain-cocycles
Each tangent direction gives a map from $\Gamma_{\infty,0}$ to matrices by differentiating the representation of Lemma~\ref{sha-patching-chain-representation}. These maps are Selmer cocycles with bounded denominators, so the bounds of Lemma~\ref{sha-patching-chain-laurent} apply to them.

% before: sha-patching-chain-linear
The relations among the $X_k$ come from linear algebra once two counts are known. The lifting of even classes bounds the codimension of the lifting families, and the rank bound bounds the families that lift in every coordinate.

% before: sha-patching-chain-rank-lift-triangular
The local condition for first-order classes is taken in the basis given by $w'_v$. In that basis the first-order deformation is upper triangular on the decomposition group with the characters of the arc on its diagonal, because the Selmer basis satisfies its own local condition with odd matrices.

% before: sha-patching-chain-rank-lift-invariants
A matrix that is triangular in the basis $w_v$ after conjugation by every element of the decomposition group, up to a triangular error, is itself triangular. The diagonal characters of the arc at $v$ are pairwise distinct, which separates the entries below the diagonal one at a time.

% before: sha-patching-chain-rank-lift
An even class that lifts to second order in every coordinate gives a first-order class. Its first-order coordinates are read off from the second-order lifts.

% before: sha-patching-chain-rank-basis
The Selmer basis placed in one first-order coordinate gives first-order classes. Products of the first-order variables vanish, so the conditions reduce to those for the basis.

% before: sha-patching-chain-rank-invariant
An even antisymmetric matrix fixed by the representation is zero. The span of the arc contains multiples of the diagonal matrix units, and $2$ is invertible.

% before: sha-patching-chain-rank-injective-generators
A homomorphism from $\mathcal A$ to a finite algebra over the fraction field is determined on $\hat{\mathcal R}$ by its values at the generators of $\mathcal R$. Agreement with $f$ on the generators spreads to $\mathcal R$ by Krull's intersection theorem, and then to its completion, because the homomorphism kills a power of the maximal ideal.

% before: sha-patching-chain-rank-injective-constant
When such a homomorphism carries the representation to a conjugate of the arc, it takes scalar values at the generators. The characteristic polynomials are conjugation invariant, and the ordinary and Taylor-Wiles relations, evaluated at eigenvectors and on diagonals, force the characters and the elements $U_{j,i}$ and $y_{j,i}$ to be scalars.

% before: sha-patching-chain-rank-injective
A tangent direction whose even class is a coboundary is zero. Its homomorphism to the dual numbers then gives a deformation conjugate to the arc, and the completion is generated by traces and characters.

% before: sha-patching-chain-rank
The second count of Lemma~\ref{sha-patching-chain-linear} is a rank bound. A tangent direction whose cocycle lifts in every coordinate contributes a first-order class, and the classes of the Selmer basis contribute $g^2$ more.

% before: sha-patching-chain-functionals-decomposition
A functional that kills $JN$ is a derivation in its dependence on the base. Expanding that derivation in a basis of the tangent space writes the functional through tangent directions and additive maps on $N$.

% before: sha-patching-chain-functionals-lift-conjugate
The odd coefficients of a deformation are combinations of the Selmer basis only up to coboundaries. Conjugating by a matrix that is the identity plus odd first-order terms removes the coboundaries and leaves the even coefficients unchanged.

% before: sha-patching-chain-functionals-lift-push
When the odd coefficients are exact combinations of the Selmer basis, the deformation can be rewritten over $R_g$ with the Selmer basis itself as its odd part. The change of variables along the coefficients preserves products, because it matches the pairing of odd and even first-order terms that produces $\eta$.

% before: sha-patching-chain-functionals-lift-second
A deformation over $R_g$ whose odd part is the first-order deformation differs from it by a square-zero term. That difference is a second-order Selmer cocycle, and its local condition at $\ell$ is the general form of Lemma~\ref{sha-patching-chain-cocycles-ordinary}.

% before: sha-patching-chain-functionals-lift
The decomposition gives a deformation over a ring of the shape $R_g$ whose odd coefficients are combinations of the Selmer basis up to coboundaries. Such a deformation is, after conjugation and a change of variables, a deformation of the first-order representation, and its second-order part is a Selmer class that lifts the even coefficients.

% before: sha-patching-chain-functionals
The relations among the $m_i$ are tested against the functionals that kill $J$. Each such functional produces a family of tangent directions, and this family lifts because the $m_i$ are dual to the Selmer basis.

% before: sha-patching-chain-dual-basis
The odd Selmer group is spanned by the Selmer basis, by the last statement of Lemma~\ref{sha-patching-chain-laurent}. Reading off the coefficients of the cocycles of the odd directions gives elements of the off-diagonal part, and they generate it by Nakayama's lemma.

% before: sha-patching-chain-presentation
These lemmas give the input of Lemma~\ref{sha-patching-chain-cotangent} at an arc. The elements $X_k$ generate the maximal ideal modulo $J$, and the $g^2$ relations come from the linear algebra applied to the families produced by the functionals.

% before: sha-patching-chain-thick
Thickness is a failure of semisimplicity over a field, and it has to be moved from a prime of $B$ to a prime of a completion of the patched ring. Non-semisimplicity of a finite-dimensional quotient survives extension of the field, and passes to an algebra that maps onto the quotient.

% before: sha-patching-chain-ideals
Thickness of a prime of $B$ has to be compared with the failure of semisimplicity of the algebra of an arc datum over the patched ring. This needs the relation ideal of the arc datum to map into the ideal that defines $E_0$. It is checked on generators, once closures are known to behave well under maps of groups and under base change.

% before: sha-patching-chain
Patching combines the Taylor-Wiles data of the Hecke ideal with the bounds at the arcs. The patched module has large support, the bounds make the patched ring small near each arc, and together they put every stable or thick prime through an arc in the support, which is the zero locus of the Hecke ideal.

% before: sha-patching-chain-burnside
Stability is defined by a nonvanishing determinant of traces, which is the form the patching argument uses. Over an algebraic closure of the residue field it is the same as the absolute irreducibility of the representation with the given determinant, which is how stable points arise.

% before: sha-patching-bound
The patching bound is stated for nice arcs, before normalization. It spreads the zero locus of the Hecke ideal from a stable prime below an arc to every prime below it, and along chains of endoscopic arcs joined by thick primes.

% before: sha-descent-datum
The descent step of the proof of Theorem~\ref{sha-residually-reducible-automorphy} needs an abelian CM field over which the ordinary points become representations of a group with controlled cohomology, together with seeds that come from modular forms. Its construction uses only a short list of properties of the characters $\eta$ and $\Psi$, which Eisenstein data and the data of Definition~\ref{sha-lifts-problem} both have.

% before: sha-descent-auxiliary
The field is cut out by four open subgroups of $\operatorname{Gal}(\bar{\mathbf Q}/\mathbf Q)$. A real quadratic field ramified at a prime outside $M$ separates characters, a prime $w$ that splits completely in the fields of $\eta$ and $\bar\Psi$ carries the seeds, and the layers of the cyclotomic $\mathbf Z_\ell$-extension make the local degrees above $\ell$ large. The Ferrero-Washington theorem bounds the class groups along these layers.

% before: sha-descent-subgroup
The intersection of these subgroups, and its intersection with $\ker\eta$, fix an abelian CM field and its real subfield. The properties needed later are read off on the Galois side: the field is unramified over its real subfield, split there at the primes dividing $M$, has large local degrees above $\ell$, and the residual characters stay distinct on it.

% before: sha-descent-link
Over the real subfield, the ordinary points factor through the quotient by the tame inertia on which $\eta$ and $\bar\Psi$ are trivial. This quotient, with its local data at the places above $\ell$, is a link, and its representations with the ordinary relations at $\ell$ are the linked points.

% before: sha-descent-link-exists
The quotient by tame inertia exists for the field of Lemma~\ref{sha-descent-subgroup}. The places above $\ell$ are chosen as conjugates of $\mathfrak B$ by elements of $\ker\eta$, so that the ordinary condition at $\mathfrak B$ can be moved to each of them.

% before: sha-descent-residual
The automorphy of linked points needs bounds on the cohomology of the link with local conditions at the places above $\ell$. These are the residual conditions; a constant $C_0$ bounds the relevant Selmer groups, and the link is large when its local degrees exceed what $C_0$ requires.

% before: sha-descent-strict
Two of the residual conditions ask for nonzero strict Selmer classes for the characters $\bar\Psi^{-1}\eta$ and $\bar\Psi^{-2}\eta$. They are produced over $\mathbf Q$ by the Greenberg-Wiles formula, with all classes allowed at the auxiliary prime $w$, and then restricted to the field.

% before: sha-descent-link-local-tate
The remaining residual conditions concern the Galois group of a totally real field modulo tame inertia, and they come from class field theory and from the cohomology of Galois groups with finite coefficients. Tate's local theorems count the cohomology of a decomposition group with coefficients in a finite module killed by a power of $\ell$: its Euler characteristic, the duality with the dual module twisted by the cyclotomic character, and the vanishing of the third cohomology.

% before: sha-descent-link-local
For a vector space over a finite field of characteristic $\ell$, these theorems become statements about dimensions. The duality computes the second cohomology from the linear forms twisted by $\omega$, and the vanishing of the third cohomology makes the second cohomology right exact.

% before: sha-descent-link-unramified
At a place not above $\ell$, the decomposition group modulo inertia is procyclic. So its quotient by a subgroup of inertia of index prime to $\ell$ has no second cohomology with coefficients of characteristic $\ell$, and cocycles that factor through such a quotient can be discarded.

% before: sha-descent-link-global-tate
The group of a link is a quotient of a Galois group with restricted ramification, whose cohomology with finite coefficients is described by the theorems of Tate and Poitou. For a finite module killed by a power of $\ell$, they give the finiteness of the cohomology, its Euler characteristic, and the exact sequence that compares it with the cohomology of the decomposition groups and of the dual module twisted by the cyclotomic character.

% before: sha-descent-link-global
For a vector space over a finite field of characteristic $\ell$, the Euler characteristic formula over a totally real field and the end of the Poitou-Tate sequence become statements about dimensions. They will be applied to a Galois group with restricted ramification that maps onto the group of a link.

% before: sha-descent-link-greenberg-wiles
The bound on the first cohomology of the link with trivial coefficients comes from the Greenberg-Wiles formula. It compares a Selmer group with the dual Selmer group for the orthogonal local conditions.

% before: sha-descent-link-hermite
Condition (5) asks for finiteness of the characters of order $\ell$ of each open subgroup of the link's group. Over a number field that contains the $\ell$-th roots of unity, the characters with bounded ramification are finite in number, by Hermite's theorem.

% before: sha-descent-link-hilbert
Condition (10) asks that the characters of $\Gamma_0$ to finite abelian $\ell$-groups that kill inertia above $\ell$ have bounded order. Over the CM field $L$ such a character is unramified everywhere, and class field theory bounds its order by the class number.

% before: sha-descent-link-cone
The groups $H^i_b$ of a link sit in a long exact sequence with the cohomology of $\Gamma$ and of the decomposition groups with coefficients in $W/b_v$. Counting dimensions along it reduces condition (11) to a bound for the cohomology of $\Gamma$ and the local Euler characteristic.

% before: sha-descent-link-inflation
The group of a link is a quotient of a Galois group with restricted ramification by the closed normal subgroup generated by some tame inertia. The five-term exact sequence bounds the change in the first and second cohomology by the equivariant homomorphisms on the generators.

% before: sha-descent-link-tame
At a place not above $\ell$, the equivariant homomorphisms on a subgroup of inertia of index prime to $\ell$ are counted by the second cohomology of the decomposition group. This is what lets the bound of Lemma~\ref{sha-descent-link-inflation} cancel against the local terms of the Poitou-Tate sequence at these places.

% before: sha-descent-link-generators
Condition (9) asks, at each place $v$ above $\ell$, for $d_v$ elements of inertia that control every homomorphism from the decomposition group to a finite abelian $\ell$-group. The characters of order $\ell$ of the decomposition group form a space of dimension $d_v+1$, and Alpöge finds the elements by linear algebra in this space and a Frattini argument.

% before: sha-descent-link-strict
The Selmer groups with the strict condition at the places above $\ell$ are bounded by the class group invariant $r_\ell$ of an open subgroup on which the character is trivial. A strict class restricts to an everywhere unramified homomorphism on this subgroup, and a class that restricts to zero is a coboundary, since the index is prime to $\ell$.

% before: sha-descent-link-setting
The other conditions concern the quotient of the Galois group of a number field by the tame inertia in an open normal subgroup of index prime to $\ell$. A realized link is such a quotient, for the field $L^+$ and the subgroup $H_{\eta,\Psi}$.

% before: sha-descent-link-finite
Condition (5) holds for every such quotient. An open subgroup of the quotient contains, up to finite index, the image of the Galois group of a number field, and the characters of order $\ell$ that factor through the quotient are unramified outside a fixed finite set of places.

% before: sha-descent-link-level
The theorems of Galois cohomology apply to Galois groups with restricted ramification. Since $K$ is open, only finitely many places ramify in its fixed field, and the quotient is a quotient of such a Galois group by tame inertia at finitely many places.

% before: sha-descent-link-h2
The global theorems for the Galois group with restricted ramification give a bound for the cohomology of the quotient. The tame inertia killed at the extra places costs at most the local second cohomology there, which the end of the Poitou-Tate sequence pays back.

% before: sha-descent-link-character
The first cohomology of a character of the quotient is bounded through the Greenberg-Wiles formula for the dual character, with the strict condition above $\ell$ and the tame condition at the extra places. The Selmer group on the other side of the formula consists of classes that are unramified on an open subgroup, and Lemma~\ref{sha-descent-link-strict} bounds it by $C_0$.

% before: sha-descent-link-cohomology
For a realized link, the field $L^+$ and the subgroup $H_{\eta,\Psi}$ satisfy the hypotheses of these lemmas. They give the remaining residual conditions, with the class group bound $C_0$ for the Selmer groups.

% before: sha-descent-residual-holds
With the auxiliary subgroups and the class group bound in place, all the residual conditions hold. The layer $m$ is chosen with $\ell^m>3C_0+7$, so that the link is large.

% before: sha-descent-seed
The automorphy of linked points starts from seeds: representations of rank two and three that come from modular forms congruent to Eisenstein series. A seed datum records the characters and the auxiliary prime from which they are built.

% before: sha-descent-seed-exists
The character $\chi'$ of a seed datum is a Teichm\"uller lift of $\bar\Psi\omega^{-j}$, and the character $\theta$ of order $\ell$ comes from the auxiliary prime. The congruence condition on $w$ makes it split in the fields of $\eta$ and $\chi'$.

% before: sha-descent-construction
Together these give the field, the base element, the link and the seed datum of the descent step. The base element $g_1$ is taken in a decomposition group above $\ell$, where $\bar\Psi^2$ is nontrivial.

% before: sha-descent-eisenstein
An Eisenstein datum gives a descent datum, with its imaginary quadratic field as the field of $\eta$. The conditions follow from the definition of an Eisenstein datum.

% before: sha-descent-points
The ordinary points factor through the link, because they are trivial on the tame inertia on which $\eta$ and $\Psi$ are trivial. Moving the ordinary condition at $\mathfrak B$ to each place $A_v$ gives the local characters of a linked point.

% before: sha-hecke-glue
The determinant over the Hecke algebra is built from the Galois representations at its classical points. Representations at the points of a compact ring glue to a determinant over the ring once the points separate the elements of the ring and the characteristic polynomials lie in the ring on a dense set of group elements.

% before: sha-patching-blocks
Normalizing an arc makes its symmetric form the identity over the power series ring. In each block this needs a stable lattice on which the form is unimodular, which exists after the substitution $T\mapsto T^2$, and square roots of units, which exist after enlarging the residue field.

% before: sha-patching-selmer
The tangent and obstruction bounds at an arc concern the deformations that join its two blocks. Their Selmer classes take values in the off-diagonal blocks, and the dual of the first-order deformation they define is described by cocycles with values in a module over a ring of dual numbers.

% before: sha-patching-twisted
Detecting the dual Selmer group comes down to showing that certain cocycles over the fraction field of the coefficients are coboundaries. This is linear algebra over a field: the representation spans all block diagonal matrices on the kernel of the cyclotomic character, and the cyclotomic character is nontrivial.

% before: sha-patching-detection
The elements $\sigma_j$ that fix the Taylor-Wiles places must detect the dual Selmer group of each arc. The first step is that the whole of $\Gamma'$ detects it, which follows from Lemma~\ref{sha-patching-twisted} after clearing denominators.

% before: sha-patching-modules
The obstruction bound turns a bound on the size of a second cohomology group into a relation among obstruction classes. Over a discrete valuation ring, if the quotients of a finite module by the powers $t^e$ have at most $|\kappa|^{ge+C}$ elements, then any $g+1$ of its elements satisfy a relation with a unit coefficient after multiplication by $t^C$.

% before: sha-automorphy-eisenstein-series
The seeds of rank two and three in the proof of Theorem~\ref{sha-residually-reducible-automorphy} come from cusp forms congruent to Eisenstein series. Those cusp forms are built from the Eisenstein series of level $\Gamma_1(N)$ attached to residue classes, whose constant terms are values of the functions $L_k$ at roots of unity and which $\Gamma_0(N)$ permutes.

% before: sha-automorphy-eulerian
The constant terms of these Eisenstein series must be controlled $\ell$-adically. At roots of unity of order prime to $\ell$ the values of $L_k$ are integral, and three combinations of them separate two residue classes that are not equal or opposite.

% before: sha-automorphy-unit-form
An Eisenstein series becomes a cusp form after subtracting a multiple of a form whose constant term is $1$ at $\infty$ and which vanishes at the cusps not equivalent to $\infty$ under $\Gamma_0(N)$. Alpöge builds such a form with integral coefficients from products of Eisenstein series of small weight and forms of level one, using the separation of the previous lemma.

% before: sha-automorphy-constant-term
The multiple to subtract is the constant term of the Eisenstein series with character $\chi$, and the cusp form is congruent to the Eisenstein series only if this constant term is divisible by the maximal ideal. An auxiliary prime $w\equiv1\pmod\ell$ and a character of order $\ell$ modulo $w$ achieve this through a distribution relation for $L_k$.

% before: sha-automorphy-deligne-serre
The cusp form obtained this way has Hecke eigenvalues congruent to those of the Eisenstein series, but it need not be an eigenform. The lifting lemma of Deligne and Serre replaces it by an eigenform with the same residual eigenvalues.

% before: sha-automorphy-eisenstein-eigenform
The three lemmas combine into eigenforms of every large weight of the right parity whose eigenvalues are congruent to $1+\chi(p)p^{k-1}$. Their Galois representations are therefore residually reducible, with constituents $1$ and $\chi_{\mathrm{cyc}}^{k-1}$ times the character of $\chi$.

% before: sha-automorphy-ordinary-forms
The Galois representation of such an eigenform has a $U_\ell$-eigenvalue that is a unit, so it is ordinary at $\ell$. Its shape at the primes away from and above $\ell$ comes from local-global compatibility, and its irreducibility from the two distinct Hodge-Tate weights together with the bound of Jacquet and Shalika.

% before: sha-automorphy-eigenform-representation
The eigenforms of Proposition~\ref{sha-automorphy-eisenstein-eigenform} satisfy the hypotheses of the previous lemma. The Chebotarev density theorem spreads the congruence of their Hecke eigenvalues to a congruence of traces at every element of the Galois group.

% before: sha-automorphy-eisenstein-galois
The seeds of rank two are twists of these representations by characters of a seed datum, so the Eisenstein congruences are needed in Galois-theoretic form. The proposition gives them with the determinant, the residual constituents and the local shape that the seeds require.

% before: sha-automorphy-has-conjugates
The seeds of rank three are built from the adjoint of a modular form, and the adjoint lift needs a condition on its Frobenius eigenvalues away from $\ell$. That condition is proved by comparing an eigenform with its images under automorphisms of $\mathbf C$, so the Eisenstein congruences must produce eigenforms whose conjugates are again cusp forms of the same level.

% before: sha-automorphy-conjugation
Each ingredient of the construction of Proposition~\ref{sha-automorphy-eisenstein-eigenform} is compatible with automorphisms of $\mathbf C$. An automorphism permutes the Eisenstein series of level $\Gamma_1(N)$ through its action on the $N$-th roots of unity, and it carries the constant term of the Eisenstein combination for a character to that for the conjugate character.

% before: sha-automorphy-deligne-serre-conjugates
The cusp forms with conjugate forms make up a space that is stable under the Hecke operators. The lemma of Deligne and Serre can therefore be applied inside this space, and the eigenform it produces keeps the property.

% before: sha-automorphy-eisenstein-conjugates
With these compatibilities, the congruence between a cusp form and an Eisenstein series can be carried out for every conjugate character at once. The resulting eigenforms are those of Proposition~\ref{sha-automorphy-eisenstein-eigenform}, with conjugate forms.

% before: sha-automorphy-petersson
The comparison of eigenvalues uses two classical facts about eigenforms with primitive character. Away from the level the Hecke eigenvalue is real up to the value of the character, and at the primes dividing the level its absolute value is known.

% before: sha-automorphy-half-step
Two Frobenius eigenvalues of the adjoint differ by a factor $p$ in absolute value when the eigenvalues of the form differ by a factor $p$ or $\sqrt p$. The second case is excluded by an algebraic argument: it forces the ratio of the roots to be $\sqrt p$, and an automorphism of $\mathbf C$ that changes the sign of $\sqrt p$ then equates a negative number with a positive one.

% before: sha-automorphy-hecke-roots
For a prime not dividing the level, the roots of the Hecke polynomial of an eigenform with conjugate forms satisfy the hypotheses of the previous lemma. The conjugate of the eigenform under an automorphism changing the sign of $\sqrt p$ is again a normalized eigenform, and both satisfy the first relation of Lemma~\ref{sha-automorphy-petersson}.

% before: sha-automorphy-frobenius-ratios
The Galois representation of such an eigenform satisfies both conditions on its Frobenius eigenvalues at every prime other than $\ell$. At the primes not dividing the level they come from the Hecke polynomial and the bound of Jacquet and Shalika; at the primes dividing it, both eigenvalues have the absolute value of the $U_p$-eigenvalue.

% before: sha-automorphy-tate-free
The adjoint lift needs, at the places away from $\ell$ where inertia acts through a finite group, a Frobenius element whose eigenvalues are not related by a power of the norm of the place, even up to roots of unity. At such places it makes the local components of the descended representation principal series of the required kind.

% before: sha-automorphy-adjoint-tate-free
The eigenvalues of the adjoint at a Frobenius element are $1$ and the two ratios of the eigenvalues of the representation. The two inequalities on the eigenvalues of the representation exclude every relation $x=q_wy$ between the absolute values of these three numbers.

% before: sha-automorphy-eisenstein-tate-free
Combining these, the Eisenstein congruences can be chosen so that the adjoint is Tate-free. The proposition adds this, and the inequalities behind it, to the conclusions of Proposition~\ref{sha-automorphy-eisenstein-galois}, and it supplies the representations from which the seeds of rank three are built.

% before: sha-automorphy-twist-glue
The restrictions of these representations to a CM field are made conjugate self-dual by a twist, a character of the Galois group of the CM field whose product with its conjugate is prescribed. On the ideles, such a character extends a character of the totally real subfield, with a correction on the units that must have finite order. That extension is a question about abelian groups with a divisible target, and it is settled first in that generality.

% before: sha-automorphy-twist-artin
Characters of Galois groups and characters of ideles are compared through global class field theory. Besides the Artin map, three consequences are used: a character trivial on complex conjugations ignores the infinite components, a character of the ideles trivial on the multiplicative group of the field descends to the Galois group, and Frobenius elements determine continuous characters.

% before: sha-automorphy-twist-units
The hypotheses of the gluing lemma concern units. For the ideles, a power of every unit is congruent to $1$ modulo a given integer; for a CM field, the quotient of a global number by its conjugate has absolute value $1$ at every complex place, so a unit of this kind is a root of unity, and it is $1$ if it is congruent to $1$ modulo an integer at least $3$.

% before: sha-automorphy-twist-places
The identity between the two Galois characters is checked at Frobenius elements, so it needs the behaviour of places in a quadratic extension. These facts describe the split and inert places, the image of a local uniformizer in the ideles of the larger field, and the effect of conjugation on decomposition and inertia groups.

% before: sha-automorphy-twist-transfer
For a subgroup of index two, a character of the subgroup times its conjugate is the restriction of a character of the whole group, the composite with the transfer. For index two the transfer has an explicit formula, which gives this character directly.

% before: sha-automorphy-twist-character
These combine to give the twisting character for a CM field over its totally real subfield. The given character of the larger Galois group becomes a character of the ideles of the totally real field, which the gluing lemma extends to the ideles of the CM field, and class field theory returns it to the Galois side; the transfer identity is then checked at Frobenius elements.

% before: sha-automorphy-algebraic-twist
To make the restriction to the CM field of a representation from Proposition~\ref{sha-automorphy-eisenstein-tate-free} conjugate self-dual, it is twisted by a character whose product with its conjugate is prescribed. The twist must be algebraic, unramified away from the integer of the link, and a power of the cyclotomic character up to finite order at the places above $\ell$, so that the twisted representation stays ordinary.

% before: sha-automorphy-adjoint-input
The twisted representation of rank two has a list of properties that the seeds of rank three need: automorphy, a polarization, ordinarity at $\ell$, control of ramification and a Tate-free adjoint. Its adjoint, twisted by the inverse of the cyclotomic character, is to be shown automorphic for $GL_3$, through definite unitary groups in two and three variables, with principal series local components at chosen split places.

% before: sha-automorphy-unitary-comparison
The descent to the definite unitary group in two variables gives a representation with the Satake parameters of the cuspidal representation at almost all split places. To recover the local components at every split place and the weights at infinity, the base change of the descended representation is written as a sum of Speh representations, and Galois representations and the theorem of Jacquet and Shalika show that it has only one part.

% before: sha-automorphy-unitary-two
The first step descends the automorphic representation of $GL_2$ to the definite unitary group in two variables. Its local components at the split places are plain and compatible with the Galois representation, and its weights at infinity differ by $m-1$.

% before: sha-automorphy-local-ratios
The adjoint lift needs the local components at the chosen places to be principal series whose ratio of characters avoids $|\cdot|^{\pm1}$, together with its square. Away from $\ell$ this comes from the Frobenius eigenvalues, through the Tate-free condition or the choice of the place; above $\ell$ it comes from the ordinary characters, whose weights differ by $m\ge2$.

% before: sha-automorphy-adjoint-lift
Alpöge's adjoint lift passes from the unitary group in two variables to the unitary group in three variables, with local components $(a,1,a^{-1})$ for the ratio $a$ at the chosen places. It allows a second outcome, a representation in two variables with the squared ratios, which has to be excluded separately.

% before: sha-automorphy-adjoint-branches
The two outcomes of the adjoint lift are compared with the Galois representation. The first gives a cuspidal representation of $GL_3$ attached to the twisted adjoint. The second would give a Galois representation of rank two whose ratios of eigenvalues at Frobenius elements are the squares of those of the given one on a dense set, and this is excluded.

% before: sha-automorphy-plain-at
At a chosen place dividing the integer of the link, the ratio $a$ and its square avoid $|\cdot|^{\pm1}$. So no two of $a$, $1$ and $a^{-1}$ differ by $|\cdot|$, and the cuspidal representation from the first outcome has plain local components there.

% before: sha-automorphy-windowed-glue
The windowed condition for $GL_3$ is a uniform bound on the supercuspidal supports at all finite places. It follows once every place has, among the cuspidal representations attached to the Galois representation, one that is unramified or plain there, since strong multiplicity one makes them all the same.

% before: sha-automorphy-adjoint-automorphy
Combining these gives the automorphy of the twisted adjoint, with the windowed condition. A chosen place dividing the integer of the link, run through every such place, supplies the plain local components there, and the places away from that integer are unramified.

% before: sha-automorphy-lift-forms
The adjoint lift compares automorphic forms on the definite unitary groups in two and three variables. They are taken with values in a representation at infinity, of a level that is an Iwahori subgroup or a deeper subgroup at finitely many split places, and as eigenforms of the double coset operators there.

% before: sha-automorphy-lift-comparison
The adjoint lift rests on a comparison of eigenforms. An eigenform on $U_2$ of deep level at a fixed set of places and of Iwahori level at further places gives an eigenform on $U_3$ with the adjoint eigenvalues, or an eigenform on $U_2$ whose eigenvalues are squared.

% before: sha-automorphy-lift-forms-facts-reps-projections
The coordinates of the forms of a given weight and level are automorphic forms on a definite unitary group, whose automorphic quotient is compact. The space of these automorphic forms is then a direct sum of automorphic representations, and a finite-dimensional part of it is cut by finitely many equivariant projections.

% before: sha-automorphy-lift-forms-facts-reps-subquotient-form
Both cases of the next lemma end in the same way: a vector of the Jacquet module on which the torus acts through generalized eigenvalues gives an eigenform on the whole module. For a smooth representation of finite length of a commutative group, this comes from a subrepresentation maximal among those that miss the vector, and from an irreducible quotient, which is a character.

% before: sha-automorphy-lift-forms-facts-reps-subquotient-theta
On the vectors fixed by $\mathrm{Iw}_C$, the torus acts through the double coset operators of its dominant elements. In the Jacquet module these operators become the action of the torus, scaled by the modulus character, since each such double coset is a union of translates by upper triangular unipotent matrices.

% before: sha-automorphy-lift-forms-facts-reps-subquotient-injective
The projection to the Jacquet module is injective on the vectors fixed by $\mathrm{Iw}_C$. Conjugation by a suitable dominant element contracts any upper triangular unipotent matrix into $\mathrm{Iw}_C$, and the double coset operator of that element is invertible.

% before: sha-automorphy-lift-forms-facts-reps-subquotient-jacquet
A common generalized eigenvector of the torus on the vectors fixed by $\mathrm{Iw}_C$ thus gives a generalized eigenvector in the Jacquet module. Frobenius reciprocity then places the representation in a principal series whose characters are read off from the eigenvalues.

% before: sha-automorphy-lift-forms-facts-reps-subquotient
An eigenvector of the double coset operators of an Iwahori-type subgroup determines the principal series that contains a representation. Its image in the Jacquet module is an eigenvector of the diagonal torus, with the character read off from the eigenvalues.

% before: sha-automorphy-lift-forms-facts-reps
Forms of a fixed weight and level decompose according to finitely many automorphic representations, and a form can be projected to one of them without losing its eigenvalues. An eigenvector of the double coset operators at a split place then places a local component of that representation inside the corresponding principal series.

% before: sha-automorphy-lift-forms-facts-operators-product
The double coset operators of the Iwahori subgroup multiply like the dominant elements of the diagonal torus. This rests only on the factorization of the Iwahori subgroup into upper triangular, diagonal and lower triangular parts, and on the way conjugation by a dominant element contracts the first and expands the last.

% before: sha-automorphy-lift-forms-facts-span-local
In a principal series whose characters are trivial on the diagonal part of the level, the ordinary vector of the Hida theory is fixed by the whole level. Its eigenvalues for the operators $U_\alpha$ then give its eigenvalues for all the double coset operators of strictly dominant diagonal matrices, by the product formula for these operators.

% before: sha-automorphy-lift-forms-facts-span-isotypic
Under the group at one split place, an automorphic representation of a definite unitary group is a sum of copies of its local component there, because the rest of the group commutes with that place and the representation is irreducible. A space of functions with values in a finite-dimensional space and coefficients in the representation therefore contains a copy of the local component.

% before: sha-automorphy-lift-forms-facts-span
Going back from an automorphic representation of $U_2$ to a form, one starts from a vector of the right weight at infinity and improves it one place at a time. At each place a translate of the form is chosen in the local component, either an eigenvector of the double coset operators or a spherical vector, without disturbing the other places.

% before: sha-automorphy-lift-forms-facts-operators
An eigenform for the operators $U_t$ is refined to an eigenform for all the double coset operators of the Iwahori subgroup. This uses that the forms of a level make a finite-dimensional space on which these operators commute and multiply like the dominant elements of the torus.

% before: sha-automorphy-lift-forms-facts
To turn eigenforms into automorphic representations, and back, the proof of the adjoint lift needs a few facts about forms of a fixed weight and level. Only finitely many automorphic representations carry such forms, so a single one can be chosen for infinitely many places at once.

% before: sha-automorphy-unitary-galois
The two alternatives of the adjoint lift are compared with the Galois side through the representations attached to automorphic representations of definite unitary groups. Besides the compatibility at the split places, the proof uses their pairings and, for irreducible representations, the cuspidal representation of the base change.

% before: sha-automorphy-spread
A closed set of Galois elements that contains the Frobenius elements at almost all places of a large enough set of split places contains every element outside a closed exceptional set. This is the Chebotarev density theorem, with the complex conjugation to pass between the two places above a split place.

% before: sha-automorphy-regular-elements
The arguments with traces need, in every open subgroup, an element at which the adjoint input is far from scalar. Such elements come from inertia above $\ell$, where the cyclotomic character has infinite order.

% before: sha-automorphy-traces
Four facts about representations of profinite groups decide the alternatives of the adjoint lift. They rest on Burnside's theorem and the dimension of the algebra spanned by the fourth symmetric power.

% before: sha-automorphy-link-situation
The propositions (c) and (d) in the proof of the theorem below concern a link to an abelian CM field and the linked points over it. The following statements fix that setting once.

% before: sha-automorphy-realization
A link is made into an instance of the Hecke ideal by a realization. It identifies the residual representation of the link, twisted by an algebraic character, with a residual datum of the Hecke ideal, and so relates the universal rings on the two sides.

% before: sha-automorphy-realization-twists
The twists that make linked points conjugate self-dual are the inverse of the cyclotomic character in rank three and a Hecke character of the quadratic field in rank two. Both have the properties a realization needs.

% before: sha-automorphy-realization-exists-determinant
A realization contains the universal ordinary determinant ring of its residual datum, so that ring has to exist. It is cut out of a completed tensor product of universal rings for the determinant and for the local characters by the three kinds of relations of an ordinary pair. The residual pair shows that the quotient is not zero.

% before: sha-automorphy-realization-exists-linked
A realization also contains the universal ring of linked points. Mazur's universal lifting ring of the diagonal residual representation carries the polarization and the base element; at each place above $\ell$, an element of the decomposition group with distinct residual eigenvalues lets Hensel's lemma fix a basis in which the ordinary condition becomes triangularity.

% before: sha-automorphy-realization-exists
Realizations exist, with an auxiliary place whose residue characteristic avoids any given finite set of primes. Two realizations with different auxiliary places give the windowed condition everywhere.

% before: sha-automorphy-realization-seed-normal
A matrix that is skew for a symmetric form, with the trace of its square twice a nonzero square, can be diagonalized with eigenvalues $\zeta,-\zeta$ or $\zeta,0,-\zeta$ in a basis in which the form is antidiagonal. This is the linear algebra behind the normal form of a linked point.

% before: sha-automorphy-realization-seed-point
A polarized irreducible representation of the group of the link whose twist has the determinant of a point and the flags of its characters can be conjugated into a linked point. The element $g_0$ gives the skew matrix, and Hensel's lemma gives the square root that the normal form needs.

% before: sha-automorphy-realization-seed
The seeds of rank two and three give points of the universal ordinary determinant rings on the zero locus of every ideal with the property of automorphic points. This is how the automorphy of the seeds enters.

% before: sha-automorphy-realization-dimension
The dimension of the components of the Hecke ideal is read off from its finiteness over the Iwasawa algebra. When that algebra is a power series ring over a discrete valuation ring, every component through the ideal has the dimension of the algebra.

% before: sha-automorphy-realization-arcs
Condition (N4) on the arcs is a statement at the auxiliary place alone. An element of the decomposition group there at which the cyclotomic character is nontrivial modulo $\ell$, and at which the residual characters agree, forces it.

% before: sha-automorphy-realization-ideal
For a realization, the Hecke ideal and the patching bound provide an ideal of the universal ordinary determinant ring. Its quotient is finite over the Iwasawa algebra, its components have the expected dimension, and its arithmetic points are classical.

% before: sha-automorphy-realization-components-transfer-primes
The ring of linked points is studied through its completions at primes and through primes of these completions of a few kinds. Off-triangular primes are where the linked point is generically irreducible, endoscopic primes are where it splits into blocks of dimensions two and one, and nice primes of dimension one carry arcs to which the patching bounds apply.

% before: sha-automorphy-realization-components-transfer-twist
A linked point becomes a point of the universal ordinary determinant ring after it is pulled back to the Galois group of the residual datum and twisted by the character of the realization. The conditions on ordinary pairs follow from the conditions on linked points and the properties of the twist.

% before: sha-automorphy-realization-components-transfer-dense
To show that the map from the universal ordinary determinant ring lands in the subalgebra of weight zero, it suffices to control a set of topological generators. A set of elements on which maps to complete local rings are determined generates a dense subalgebra.

% before: sha-automorphy-realization-components-transfer-coefficients
The characteristic polynomials of the twisted linked point have coefficients that are sums of products of entries of total weight zero. So they lie in the subalgebra of weight zero.

% before: sha-automorphy-realization-components-transfer-onto
For the converse inclusion, the subalgebra of weight zero must be recovered from traces and characters. A diagonal matrix in the span of the images of the linked point gives the matrix units, and with them every product of weight zero.

% before: sha-automorphy-realization-components-transfer-arc-curve
An arc has values in a power series ring over a finite field. A prime of dimension one containing $\ell$ gives a one-dimensional complete local domain of characteristic $\ell$, and its normalization is such a power series ring.

% before: sha-automorphy-realization-components-transfer-arc-triangular
The flags of an arc come from a single element of each decomposition group whose characters take values with distinct residues. The Lagrange idempotents of that element conjugate the whole representation into upper triangular form.

% before: sha-automorphy-realization-components-transfer-arc-flag
Residual condition (7) supplies such an element at every place above $\ell$. So the restriction of the universal linked point to each decomposition group is triangular after any map to a local ring.

% before: sha-automorphy-realization-components-transfer-arc
The patching bounds are stated at arcs, which are points with values in a power series ring over a finite field. A prime of dimension one containing $\ell$ in a completion of the ring of linked points gives such an arc.

% before: sha-automorphy-realization-components-transfer-nice-separate
The conditions on an arc at a nice prime are checked at single elements of the decomposition groups. The characters of the arc have ratios that are one-units of a power series ring in characteristic $\ell$, which have no torsion, so one element can separate all of them at once.

% before: sha-automorphy-realization-components-transfer-nice-irreducible
Condition (N2) and the separable characteristic polynomials of (N3) rest on irreducibility over an algebraic closure. It is needed on every open subgroup of the kernel of the cyclotomic character, for the whole representation at an off-triangular prime and for its corner at an endoscopic one.

% before: sha-automorphy-realization-components-transfer-nice-type-e
At an endoscopic prime the entries of odd weight vanish. The arc then splits into the corner block and the middle character, which is what type E asks for.

% before: sha-automorphy-realization-components-transfer-nice
For the patching bounds to apply, the arc must be nice. The conditions defining nice primes are chosen so that the arcs they give satisfy the conditions of type I or E.

% before: sha-automorphy-realization-components-transfer-stable
The patching bound along chains needs stable primes of the universal ordinary determinant ring as starting points. An off-triangular prime of a completion gives one, as the linked point is irreducible there.

% before: sha-automorphy-realization-components-transfer-thick-irreducible
At a thick prime the representation is block triangular, with the second basis vector spanning a quotient. Its determinant cannot be that of an irreducible representation, because the middle diagonal entry would then define a character of a full matrix algebra.

% before: sha-automorphy-realization-components-transfer-thick-split
Semisimplicity fails for a similar reason. A semisimple quotient through which the block triangular representation factors would split off the quotient line, which gives a common eigenvector with second coordinate one.

% before: sha-automorphy-realization-components-transfer-thick
The chains of the patching bound pass through thick primes. A thick prime of a completion gives one of the universal ordinary determinant ring, as its extension class prevents both irreducibility and semisimplicity.

% before: sha-automorphy-realization-components-transfer-field-point
The seed of the Hecke ideal is a point over an algebraically closed field. To place it in the ring of linked points, a point whose traces and characters are integral over a complete local ring must come from that ring.

% before: sha-automorphy-realization-components-transfer-untwist-trace
The untwisted representation is known only through its traces, which lie in a complete local ring, and through the characters of the pair. To become a point of the ring of linked points it must be conjugated into the shape of a linked point with integral entries.

% before: sha-automorphy-realization-components-transfer-untwist
The seed is given as a determinant over a component of the Hecke ideal, together with an irreducible linked point over another field with that determinant. Untwisting by the character of the realization produces an irreducible linked point whose traces lie in the component.

% before: sha-automorphy-realization-components-transfer-local-point
To produce points of the Hecke ideal from linked points over discrete valuation rings, these linked points must come from the ring of linked points. The map exists already when only the diagonal entries and the characters reduce to the residual ones.

% before: sha-automorphy-realization-components-transfer
The component arguments take place in the ring of linked points, while the Hecke ideal lives in the universal ordinary determinant ring. A compatible surjection onto the subalgebra of weight zero carries the properties of the Hecke ideal across: finiteness over the Iwasawa algebra, the patching bounds at nice primes, and the seed. Conversely, an ideal that lies in every off-triangular prime gives the points of the Hecke ideal.

% before: sha-automorphy-realization-components-propagation-global
The ideal is spread over the spectrum by a connectedness argument. If the core of the space is connected, local agreement at closed points lets a property of one irreducible component pass to all the components that meet the off-triangular locus.

% before: sha-automorphy-realization-components-propagation-local
At a closed point, the argument is carried out in the completion. Components of the completion that meet in large dimension pass the property to each other through nice primes, by the two patching bounds.

% before: sha-automorphy-realization-components-propagation-charts
The local and global arguments are joined by the flat map to the completion. Its minimal primes give the irreducible components through the point, by going down.

% before: sha-automorphy-realization-components-propagation-geometry-lift
The seed gives a prime of the subalgebra of weight zero below an off-triangular prime. To reach a closed point of the core through it, the prime is first lifted to a homogeneous off-triangular prime of transcendence degree one over it.

% before: sha-automorphy-realization-components-propagation-geometry-dimension
The lifted prime must then be seen in a completion at the closed point, with the right dimension. Dimensions in a completion are bounded by those in the subalgebra of weight zero plus a transcendence degree, and the bound is attained above every prime through the point.

% before: sha-automorphy-realization-components-propagation-geometry
The global argument needs a connected core, enough closed points in it, closed points on every component, and a closed point through which the seed passes. These are properties of the spectrum of the ring of linked points.

% before: sha-automorphy-realization-components-propagation-completions
The local argument needs completions at the closed points and a comparison of their entries with those of the universal linked point. It also needs chains of minimal primes meeting in large dimension, and nice primes above primes of large dimension.

% before: sha-automorphy-realization-components-propagation-endoscopic-corner
Modulo its entries of odd weight, the universal linked point of rank three is the sum of its corner block and its middle entry. The corner is a linked point of rank two, so it comes from the ring of linked points of rank two.

% before: sha-automorphy-realization-components-propagation-endoscopic-lambda
The ideal of rank three is pulled back from the ideal of rank two through the corner. Its finiteness over the Iwasawa algebra must be checked again, as the middle entry is not seen by the corner.

% before: sha-automorphy-realization-components-propagation-endoscopic
In rank three, the thick components of a completion are reached from the endoscopic primes. An ideal that lies in every off-triangular prime in rank two gives one that lies in every endoscopic prime in rank three.

% before: sha-automorphy-realization-components-propagation-diagonal
For the characteristic polynomials in characteristic $\ell$, the diagonal of a triangular linked point must be identified with the residual characters. The characters are trivial on enough inertia above $\ell$ to determine them.

% before: sha-automorphy-realization-components-propagation
Once transferred, the properties spread the ideal over the off-triangular primes. Rank two is done first, and rank three uses it, together with the chains of thick primes through the endoscopic locus.

% before: sha-automorphy-realization-components
The component arguments spread the zero locus of that ideal over the whole space of linked points. They give points of the ideal for all irreducible linked points over complete discrete valuation rings with finite residue field, and the residual characteristic polynomials over those of characteristic $\ell$.

% before: sha-automorphy-realization-away
Classicality then makes the irreducible linked points automorphic after the twist. The windowed condition is obtained away from the auxiliary place of the realization.

% before: sha-automorphy-unramified-multiplicity
The multiplicity bound compares the dimension of a generalized eigenspace of a Hecke module at an arithmetic point with that at the points of distinct type of the same weight. A free module over the Iwasawa algebra and an idempotent cutting out part of it are the data.

% before: sha-automorphy-unramified-hecke-local-kirillov-forms
At the auxiliary place the forms are fixed by the principal congruence subgroup $K_1$, and on such vectors the integral upper triangular unipotent matrices act through a finite group. The vectors on which they act by a generic character $\psi_t$ are counted with a Kirillov form, a linear form that transforms by $\psi_t$ and detects a vector through its translates by the mirabolic subgroup.

% before: sha-automorphy-unramified-hecke-local-sums
The operator at the auxiliary place is applied to forms with coefficients in the Iwasawa algebra, so it must be defined with integer coefficients. When finitely many characters of a finite group are pairwise distinct, the sum of their inverses is a function with integer values, and the corresponding combination of translations is the order of the group times the projection onto the sum of the isotypic parts.

% before: sha-automorphy-unramified-hecke-local-projector
Over the residue field, the characters of the upper triangular unipotent group that are nontrivial on every entry next to the diagonal are pairwise distinct, and the sum of their inverses has integer values. Lifted to integral matrices, they give an operator whose eigenspace on the vectors fixed by $K_1$ is the sum of the spaces $W(\psi,t)$, for an additive character $\psi$ of level one and finitely many $t$ with unit entries.

% before: sha-automorphy-unramified-hecke-local-move
The Kirillov model of a local component is given for one additive character and the character $\psi_{(1,\dots,1)}$ of the unipotent group. Conjugation by a diagonal matrix, together with a rescaling of the additive character, moves it to every $\psi_t$ whose entries are nonzero.

% before: sha-automorphy-unramified-hecke-local-nonzero
A Kirillov form is nonzero on some translate of any nonzero vector fixed by $K_1$. Averaging this translate over $K_1$, and then over the integral unipotent matrices against $\psi_t$, keeps the value of the form nonzero when the additive character has level one and the entries of $t$ are units.

% before: sha-automorphy-unramified-hecke-local-injective
A Kirillov form for such $\psi$ and $t$ vanishes on no nonzero vector of $W(\psi,t)$. If it vanishes on a vector of $W(\psi,t)$, it vanishes on all its translates by the mirabolic subgroup, by clearing the rows of the translating matrix one at a time modulo the maximal ideal.

% before: sha-automorphy-unramified-hecke-local
The idempotent of the multiplicity bound comes from the auxiliary place $v_0$, where the forms are invariant under the principal congruence subgroup. A combination of translations by integral matrices, with integer coefficients, cuts out of the invariants of every generic local representation a space whose dimension does not depend on the representation.

% before: sha-automorphy-unramified-hecke-generic
The local representations at $v_0$ to which the previous lemma applies are local components of cuspidal representations of $GL_n$. These are generic, and so embed in their Kirillov models.

% before: sha-automorphy-unramified-hecke-projector
On Hida's forms, the combination of translations at $v_0$, divided by the power $q^j$, is an idempotent. It preserves the localized module, since translations at $v_0$ normalize the level and commute with the Hecke operators.

% before: sha-automorphy-unramified-hecke-weights
The multiplicity bound is stated for weights of points of the universal ordinary determinant ring, while Hida's control theorem uses weights of points of the Iwasawa algebra. The two notions correspond through the character $\theta$ and the cyclotomic character on inertia.

% before: sha-automorphy-unramified-hecke-eigenspace-multiplicity
A multiplicity is the dimension of a joint generalized eigenspace of a module over the Iwasawa algebra, after base change to a field. When the module embeds into a space of classical forms by a map that keeps bases independent, this dimension can be computed on the classical side.

% before: sha-automorphy-unramified-hecke-eigenspace-specialization
The specialization of Hida's forms at an arithmetic point is determined by its control map. Translations away from $\ell$ that normalize the level preserve classical ordinary forms and commute with the control map, so they commute with the specialization.

% before: sha-automorphy-unramified-hecke-eigenspace-localization-eigenvector
A classical form that is a common eigenvector of the Hecke operators gives a character of the Hecke algebra of the finite free module it comes from. When its eigenvalues reduce to given residual eigenvalues, this character factors through the localization at them, which is therefore nonzero.

% before: sha-automorphy-unramified-hecke-eigenspace-localization
The multiplicity is defined on the localized module, while the classical forms come from all of Hida's ordinary forms. At a point whose Hecke eigenvalues reduce to the fixed residual ones, the complement of the localized module contributes nothing to the generalized eigenspace.

% before: sha-automorphy-unramified-hecke-eigenspace-galois
An automorphic part that meets this generalized eigenspace has the same Hecke eigenvalues at the split places. Its Galois representation then has the same characteristic polynomials as the one attached to the point, and so is conjugate to it when that one is irreducible.

% before: sha-automorphy-unramified-hecke-eigenspace
At an arithmetic point, Hida's control theorem turns the multiplicity of the module into the dimension of a space of classical ordinary forms with the same Hecke eigenvalues. The idempotent becomes the same combination of translations at $v_0$ on the classical forms.

% before: sha-automorphy-unramified-hecke-count-inputs
The count of classical forms runs through the automorphic representations of the definite unitary group whose forms meet the space $Z$. What it needs about them is listed here: how the space decomposes along them, which of them meet $Z$, their base change to $GL_n$, the shape of their part of $Z$ as a local eigenspace above $\ell$ times a space of invariants, and how many of them share a base change.

% before: sha-automorphy-unramified-hecke-count
The dimension of that space of classical forms is computed through the automorphic representations of the definite unitary group that contribute to it. All of them have the same base change, and their number depends only on the weight.

% before: sha-automorphy-unramified-hecke-invariants
The local factor of a contribution is a space of invariants in a tensor product of local representations, one at the places of $S_{\mathrm a}$ and one at $v_0$. Its dimension is the product of the dimensions of the invariants in the two factors.

% before: sha-automorphy-unramified-hecke-types
At the places of $S_{\mathrm a}$ the type fixes the restriction of the Galois representation to inertia. A generic representation with such a parameter is a principal series, and its restriction to the integral matrices depends only on the characters on the units.

% before: sha-automorphy-unramified-hecke-dimension
Combining the count with the local computations, the dimension cut out by the idempotent at a point is a product of three factors. None of them depends on the point, only on the weight, and none vanishes.

% before: sha-automorphy-unramified-hecke-bound
The Hecke ideal acts faithfully on the localized module of Hida's forms, which is free over the Iwasawa algebra. With the idempotent at $v_0$, the dimensions above give the multiplicity bound when the Hecke ideal is radical.

% before: sha-automorphy-unramified-hecke
In rank three the Hecke ideal can be taken to satisfy the multiplicity bound. Its density makes it radical, and the multiplicities come from the automorphic representations of the definite unitary group in three variables with a given base change.

% before: sha-automorphy-realization-presentation
Conjugation by the diagonal torus acts on the universal ring of linked points, as it does on the local rings of the Siegel case, and multiplies the entry $(i,j)$ by a power of weight $w_{n,i}-w_{n,j}$. This gives the ring a weighted presentation. The subalgebra $B$ is its part of weight zero, which makes $B$ a complete local noetherian ring and the subalgebra generated by the entries a noetherian ring.

% before: sha-automorphy-unramified-realization-irreducible
At a point of the ring of linked points whose kernel is off-triangular, the linked point is irreducible. The base element $g_0$ puts the diagonal idempotents in the span of the representation, and the nonzero entries off the diagonal then connect every coordinate line to every other.

% before: sha-automorphy-unramified-realization
Proposition (d) needs a realization of rank three with more structure: a power series Iwasawa algebra, a multiplicity bound, and a surjection onto the subalgebra of weight zero of the universal ring of linked points. These come from realizations of ranks two and three together.

% before: sha-automorphy-unramified-point
A linked point gives a point of the universal ring of linked points. Over the realization it gives an arithmetic point of the universal ordinary determinant ring.

% before: sha-automorphy-unramified-algebra-finite
The absence of tangent vectors at a point of a Hecke module is proved by a sequence of lemmas of commutative algebra. The first ones are elementary: the ring acts on the module through a finite algebra over the base, each of its elements is a root of its characteristic polynomial modulo the ideal, and a multiplicity is bounded by a root multiplicity.

% before: sha-automorphy-unramified-algebra-primitive
Over a base of characteristic $0$, the quotient by the ideal has a primitive element up to denominators. It can also be chosen so that, at a given point, its value and the base determine the point among the points of the support.

% before: sha-automorphy-unramified-algebra-factors
The characteristic polynomial of the primitive element factors over a normal base into monic irreducible polynomials. Each factor defines a prime ideal of the ring, a component of the support, and points of the base and roots of the factor lift to points on that component.

% before: sha-automorphy-unramified-algebra-separating
At a point that the primitive element separates, the root multiplicity of its value in the characteristic polynomial is a lower bound for the multiplicity of the point. Together with the opposite bound of Lemma~\ref{sha-automorphy-unramified-algebra-finite}, this ties multiplicities to root multiplicities.

% before: sha-automorphy-unramified-algebra-conductor
A tangent vector is determined by its values on a component and on polynomials in the primitive element. The derivative of the factor of a component lies in the conductor of the polynomials in the primitive element, by Euler's formula for the trace dual.

% before: sha-automorphy-unramified-algebra-derivation
Suppose that the value of the primitive element at a point is a simple root of the product of the factors. Then a tangent vector that kills the ideal and the base kills the primitive element, then the component through the point, and then everything.

% before: sha-automorphy-unramified-algebra-rigid
The multiplicity bound forces the simple root. A point of a dense set at which the multiplicity is at least that at the given point lies on a single component, with multiplicity at most the exponent of its factor, and this caps the root multiplicity at the given point.

% before: sha-automorphy-unramified-algebra-image
Multiplicities cut out by an idempotent are the multiplicities of another Hecke module, the image of the idempotent. This reduces the multiplicity bound with an idempotent to the case without one.

% before: sha-automorphy-unramified-algebra
The absence of tangent vectors is a statement of commutative algebra about a Hecke module free over a normal base. If the multiplicity at a point is at most that at a dense set of points, the support is unramified over the base there.

% before: sha-automorphy-unramified-tangent
At an irreducible arithmetic point on the zero locus of the ideal, there are no tangent vectors that kill the ideal and the inertial characters. This comes from purity and the multiplicity bound.

% before: sha-automorphy-unramified-near-dense
Points that are dense in the zero locus of an ideal stay dense on the components of largest dimension of its image under a surjection. A separating element, which lies in every minimal prime but one, places a dense point on a chosen component.

% before: sha-automorphy-unramified-near-additive
A continuous additive function on a compact group is small on every element that the finite $\ell$-quotients cannot separate from a given finite set on which it vanishes, because multiplying by $\ell^k$ shrinks its values. So it vanishes on such elements.

% before: sha-automorphy-unramified-near-tangent-points
The bound on tangent vectors uses the continuity of a tangent vector along the local characters. Points of a complete local ring with values in $\bar{\mathbf Q}_\ell$ take integral values, small on the maximal ideal, and their images are lattices; these facts make the continuity checkable.

% before: sha-automorphy-unramified-near-tangent-small
A tangent vector at such a point is determined modulo the square of the kernel, and the ring modulo $\varpi$ and that square is finite. So a high power of the maximal ideal lies in $\varpi^n$ times the ring plus the square, where the tangent vector and the point are both small.

% before: sha-automorphy-unramified-near-tangent-filtration
The local characters take values in the subring of weight $0$, but their continuity is stated for the maximal ideal of the whole ring. A monomial of a given weight and large degree contains many monomials of weight $0$ as factors, which transfers the topology of the whole ring to the parts.

% before: sha-automorphy-unramified-near-tangent-adic
For the part of weight $0$, the comparison of topologies says that the topology of $A_n$ induces on $B$ the topology of its own maximal ideal. So each local character of the universal linked point, which is continuous as a map to $A_n$, is continuous as a map to $B$.

% before: sha-automorphy-unramified-near-tangent
The local characters of the universal linked point are free to move along the inertia groups above $\ell$ only in $3d$ directions. A tangent vector that vanishes on the values of these characters at $d_v$ chosen elements of each inertia group vanishes on the whole inertia group, by the residual condition on these elements. This bounds the tangent space by $3d$ plus the tangent vectors that fix the characters on inertia.

% before: sha-automorphy-unramified-near-scale
The argument needs points at which the entries off the diagonal are small. Conjugation by the diagonal torus multiplies the parts of $A_n^\circ$ by powers of a scalar, and a scalar of suitable absolute value makes these entries small without changing the point on $B$ or which homogeneous elements vanish.

% before: sha-automorphy-unramified-near-generic
Each component of $A_n^\circ$ through the point must contain a point of distinct type at which the representation is irreducible. Points of $B$ extend to $A_n^\circ$, which is of finite type over $B$, away from a nonzero element, and the rescaling then makes them small; minimal primes are homogeneous, so the rescaled point still vanishes on the component.

% before: sha-automorphy-unramified-near-fibre
The dimension of a component of $A_n^\circ$ exceeds that of its image in $B$ by the transcendence degree of the fibre. Where some entry below and some entry above the diagonal are nonzero, products of weight $0$ make every entry algebraic over one of them, so the fibre has transcendence degree at most $1$.

% before: sha-automorphy-unramified-near-completion-flag
The deformation theory at a point, and the unit step for the tangent vectors, need the linked point in triangular form on each decomposition group above $\ell$, over any local ring that receives $A_n^\circ$. An element of the decomposition group at which the residual characters differ separates the diagonal characters, and the ordered products of the linked point then split the space into lines, as in Lemma~\ref{sha-rigidity-lagrange}.

% before: sha-automorphy-unramified-near-completion-points-subring
A linked point over a complete local ring $C$ whose residue field is larger than $k_1$ cannot be compared directly with the universal ring $A_n$, which classifies linked points over rings with residue field $k_1$. The elements of $C$ that are constants of $\mathcal O_1$ modulo the maximal ideal form a smaller complete local ring with residue field $k_1$, over which $C$ is finite. The diagonal entries and the local characters of the point take their values there.

% before: sha-automorphy-unramified-near-completion-points-triangular
The entries off the diagonal of such a point need not lie in the maximal ideal. The multiplicativity of the diagonal entries, together with the polarization, still forces the reduction to be triangular, one way or the other. So a product of entries of total weight zero with a factor off the diagonal always meets a factor that vanishes modulo the maximal ideal.

% before: sha-automorphy-unramified-near-completion-points-laurent
To bring the entries off the diagonal into a maximal ideal, each entry is multiplied by a power of a Laurent variable given by its weight. The ring generated in this way has a maximal ideal defined by the constant coefficient, and the conditions on products of weight zero make it an ideal with the right residue field.

% before: sha-automorphy-unramified-near-completion-points-cone
In the Laurent model the twisted point is a linked point in the sense of the universal ring at every finite level. The maps from $A_n$ at these levels fit together on $A_n^\circ$, the elements of $B$ having limits in the coefficient ring by completeness, and setting the Laurent variable to $1$ returns the original point.

% before: sha-automorphy-unramified-near-completion-points
The complete local ring at a point of $A_n^\circ$ is to be universal for linked points over complete local rings with the reduction of that point. Such a reduction need not be diagonal, as entries off the diagonal may be units. Its diagonal entries and local characters still reduce to the residual ones, and this is enough to produce a map from $A_n^\circ$.

% before: sha-automorphy-unramified-near-completion-unique
A map from $A_n^\circ$ to a local ring is determined by the linked point it produces, once it sends the maximal ideal of $B$ into the maximal ideal. The entries and the characters generate a dense subalgebra of $B$, and such a map is continuous on $B$ for its maximal ideal.

% before: sha-automorphy-unramified-near-completion-tangent-triangulation
Deformations of a linked point are built by lifting its matrices one small step at a time. A linked point of rank $n$ extends to the whole group $\Gamma$ through the semidirect product of $GL_n$ with the transpose-inverse involution, and its triangular forms above $\ell$ are recorded by the matrices that conjugate it into upper triangular form. These data are the objects that the lifting arguments move.

% before: sha-automorphy-unramified-near-completion-tangent-correct
When the matrices of a triangulation are lifted along a small extension, they fail to be multiplicative and triangular by terms of square zero. If these failures are the coboundary of a cochain with values in the adjoint representation, multiplying the lifts by that cochain removes them, and the result is a triangulation over the larger ring.

% before: sha-automorphy-unramified-near-completion-tangent-normal
A triangulation produced in this way need not have the normal form of a linked point: the complex conjugation need not act by the identity matrix, and the element $g_0$ need not give a diagonal matrix. Conjugating by a matrix congruent to $1$ restores both, since the residual point separates the diagonal entries at $g_0$ and $2$ is invertible. The diagonal of the triangular forms then gives the local characters.

% before: sha-automorphy-unramified-near-completion-tangent
The tangent space of a universal ring of linked points is compared with the cohomology of the adjoint representation with the triangular conditions above $\ell$. Every cocycle with these conditions comes from a tangent vector, and the diagonal matrix of the weights gives one more tangent vector, whose cocycle is a coboundary. These are the two inputs of the count of the tangent space.

% before: sha-automorphy-unramified-near-completion-obstructions
A presentation of a complete local ring by power series comes from control of the lifts along small extensions. An obstruction space receives, linearly in the extension, the classes that decide whether a lift exists.

% before: sha-automorphy-unramified-near-completion-obstruction
For a universal ring of linked points the obstructions lie in the second cohomology of the adjoint representation with the triangular conditions above $\ell$. The Euler characteristic condition among the residual conditions bounds the dimension of this group.

% before: sha-automorphy-unramified-near-completion-mazur
A finite-dimensional obstruction space bounds the number of relations in a presentation by power series in as many variables as the dimension of the tangent space. Mazur's argument for deformation rings gives this for any obstruction space.

% before: sha-automorphy-unramified-near-completion-deformation
The complete local ring of $A_n^\circ$ at a point is a deformation ring: it carries the linked point at that point, universally. Its tangent space and its obstructions are measured by the cohomology of the adjoint representation with the triangular conditions above $\ell$, which is the cohomology that the Euler characteristic condition among the residual conditions controls.

% before: sha-automorphy-unramified-near-completion
The heights of primes of $A_n^\circ$ are computed in a complete local ring at each point. Its presentation by power series with $nd+1$ more variables than relations comes from the count of cocycles of the linked point and their obstructions.

% before: sha-automorphy-unramified-near-heights
The regularity argument needs both bounds on the height at the point: from below by the presentation of the complete local ring, and from above by the dimension bound for the components of the zero locus of $J_3$. Together they make that ring a complete intersection of the expected dimension, which is unmixed, and every component through the point has the full dimension.

% before: sha-automorphy-unramified-near-determined-cofinal
A tangent vector is continuous on $B$ for the topology of its own maximal ideal, while the density of a subalgebra of $B$ is measured by the powers of the maximal ideal of $A_n$. In a complete local noetherian ring, every decreasing sequence of ideals with intersection $0$ eventually lies in each power of the maximal ideal, and this compares the two topologies.

% before: sha-automorphy-unramified-near-determined
A tangent vector of $A_n^\circ$ is determined by its values on the entries and on the local characters. These generate a dense subalgebra of $B$, and a tangent vector is continuous on $B$ for its maximal ideal.

% before: sha-automorphy-unramified-near-torus
The last step of the regularity argument bounds the tangent vectors that fix the characters on inertia. If each acts on the entries through the adjoint action of a matrix, the polarization and the irreducibility of the point force that matrix to come from the diagonal torus up to a scalar, and the ordinary relations at $\ell$ make the tangent vector vanish on the local characters. Such tangent vectors then form a space of dimension at most $1$.

% before: sha-automorphy-unramified-near-weight
The tangent vectors at the small points are controlled through a vanishing theorem that needs the local characters to be generic: some element of the decomposition group with trivial cyclotomic character must separate them, in every power. A weight that is twice as large at one embedding as at the others forces this at every point of that weight, through a unit whose conjugates have product $1$.

% before: sha-automorphy-unramified-near-locally-generic
The vanishing theorem applies to points of the universal ordinary determinant ring that come from automorphic representations and whose Galois representations are irreducible and locally generic. Its conclusion concerns cocycles of the adjoint representation that are ordinary at the places above $\ell$ and unramified at the auxiliary places.

% before: sha-automorphy-unramified-near-vanishing
At an automorphic locally generic point the ordinary deformations of the Galois representation are trivial to first order. This is the input that bounds the tangent spaces at the small points.

% before: sha-automorphy-unramified-near-small-continuity
The Galois representation at a small point must be continuous, and so must the cocycles attached to tangent vectors at such a point. The entries of the universal linked point vary continuously modulo powers of the maximal ideal of $B$ times their span, by the Artin-Rees lemma, and a point or a tangent vector is small on high powers of that ideal.

% before: sha-automorphy-unramified-near-small
The small points of $A_3^\circ$ with off-triangular kernel give such points of $B_3$, after the twist of the realization. Their automorphy is the last part of Lemma~\ref{sha-automorphy-unramified-realization}, and their irreducibility comes from the off-triangular kernel.

% before: sha-automorphy-unramified-near-matrix-adjoint
An ordinary cocycle of the vanishing theorem is a cocycle of the whole group over $L^+$, acting on matrices through an adjoint extension. Such an extension exists for every polarization of a representation of the subgroup fixing $L$ whose values span the matrices, since the polarization then gives an involution compatible with the adjoint action.

% before: sha-automorphy-unramified-near-matrix-unipotent
At the auxiliary places the twisted lift is trivial on the inertia elements that kill the residual characters. There it is a matrix congruent to $1$ whose order is prime to $\ell$, and such a matrix is $1$.

% before: sha-automorphy-unramified-near-matrix
A tangent vector of $A_3^\circ$ at such a point is a first-order deformation of the linked point, and so gives a cocycle of the adjoint representation. The cocycle is ordinary because the tangent vector kills the characters on inertia, so it is a coboundary, and its matrix is the one Lemma~\ref{sha-automorphy-unramified-near-torus} needs.

% before: sha-automorphy-unramified-near
The ideal vanishes on a neighbourhood of the point of the subalgebra $A_3^\circ$ of the universal ring of linked points. So tangent vectors at that point kill it.

% before: sha-automorphy-unramified-tangents-rank
Tangent vectors with values in $\mathcal O$ are found by counting tangent vectors modulo $\ell^k$. The tangent vectors at a point form the dual of a finitely generated module over $\mathcal O$, and if that module had free rank at most one, their number modulo $\ell^k$ would grow at most like $|\mathcal O/\ell^k|$. A growth like $|\mathcal O/\ell^k|^2$ therefore forces two independent tangent vectors.

% before: sha-automorphy-unramified-tangents-unit-duality
The unit step needs an unramified class whose cup product with a crystalline extension class vanishes. This vanishing comes from the local duality of Bloch and Kato between the crystalline classes of a representation and those of its Tate dual.

% before: sha-automorphy-unramified-tangents-unit-crystalline-periods
The facts on crystalline representations below are proved from a few properties of Fontaine's period rings. The ring of crystalline periods sits inside the field of de Rham periods, whose Galois invariants are the base field, and it contains the Witt vectors of the tilt and a period of the cyclotomic character that is a unit.

% before: sha-automorphy-unramified-tangents-unit-crystalline-tilt
A representation that is trivial on inertia will be trivialized over the Witt vectors of the tilt, by a matrix built from the finite fields inside the tilt. That construction needs the Galois group to fix each element of these finite fields on an open subgroup, and an element fixing all of them to lie in the inertia group.

% before: sha-automorphy-unramified-tangents-unit-crystalline-lang
The matrix that trivializes a representation trivial on inertia is found by a form of Lang's theorem. Hilbert's Theorem~90 over a finite field gives it modulo $\ell$, and averaging over finite quotients of the group lifts it one power of $\ell$ at a time.

% before: sha-automorphy-unramified-tangents-unit-crystalline
The unit step passes from a potentially crystalline point to a crystalline extension class through a chain of operations on representations: restriction to a smaller group, change of coefficients, subquotients, twists by characters and extensions of the trivial representation. Each keeps the crystalline property, measured by the dimension of the invariants in the tensor product with $B_{\mathrm{cris}}$. The bound on that dimension, which rests on the regularity of $B_{\mathrm{cris}}$, gives most of them at once.

% before: sha-automorphy-unramified-tangents-unit-cup
For the cyclotomic character the Tate dual is the trivial representation, whose crystalline classes are the unramified homomorphisms. So the cup product of a crystalline class with cyclotomic coefficients and the homomorphism that measures the Frobenius is a coboundary, and the cochain that trivializes it is the $\beta$ of the unit step.

% before: sha-automorphy-unramified-tangents-unit
At a place above $\ell$ where two diagonal characters of a crystalline linked point differ by the cyclotomic character, the extension between them is constrained by $p$-adic Hodge theory. The weights force the two characters to be adjacent on the diagonal, and the crystalline condition gives an unramified class whose cup product with the extension class is a coboundary.

% before: sha-automorphy-unramified-tangents-count-counts
Cocycles modulo $\ell^k$ are counted through two formulas for the link, one at each place above $\ell$ and one for the whole group. The local one is the Euler characteristic formula, with the dual term written as equivariant maps to the cyclotomic character; the global one is the inequality of Greenberg and Wiles for cocycles with local conditions.

% before: sha-automorphy-unramified-tangents-count-inputs-unramified
The global count runs over a Galois group with restricted ramification, which has places that the link does not see. At such a place, not above $\ell$, the cocycles that are coboundaries on inertia are at least as many as the elements of the module, because the decomposition group modulo inertia is procyclic.

% before: sha-automorphy-unramified-tangents-count-inputs-level
The group of the link is a quotient of the Galois group of $L^+$ by tame inertia at the places outside $\ell$, cut down by an open subgroup. Only finitely many of these inertia groups leave that open subgroup, so the quotient factors through a Galois group with restricted ramification, and its kernel there is generated by inertia at finitely many places.

% before: sha-automorphy-unramified-tangents-count-inputs-global
Over a Galois group with restricted ramification, the formula of Greenberg and Wiles bounds the classes with local conditions below. Written with cocycles in place of classes, and with the dual Selmer group bounded by $1$, it becomes an inequality between numbers of cocycles.

% before: sha-automorphy-unramified-tangents-count-inputs
Both formulas come from Galois cohomology over $L^+$ and its completions, through the quotient $\Gamma$ of the Galois group that defines the link. The complex conjugations at the infinite places of $L^+$ enter the global formula.

% before: sha-automorphy-unramified-tangents-count-cocycles
A triangulation of a linked point over $\mathcal O_1$ gives local conditions for cocycles of the adjoint representation modulo $\ell^k$: at each place above $\ell$, cocycles that are upper triangular after the change of basis $g_v$ and have no diagonal part on inertia, up to coboundaries. The cocycles counted below are the global cocycles with these local conditions, normalized at the conjugation and at $g_0$.

% before: sha-automorphy-unramified-tangents-count-triangulation
The linked point given by $\xi$ has triangular forms at every place above $\ell$ by the flag lemma. At the place of the unit step the triangular form can be taken to be the one that the unit step provides.

% before: sha-automorphy-unramified-tangents-count-local-conj
At a place above $\ell$ the triangular form $t_v$ acts on matrices by conjugation, and the local conditions are read off from the entries. Conjugation by an upper triangular matrix keeps upper triangular matrices upper triangular, fixes their diagonal, and changes the entries just above the diagonal by a simple formula; when the diagonal entries of $t_v$ differ by units, a matrix whose conjugates differ from it by upper triangular matrices is itself upper triangular.

% before: sha-automorphy-unramified-tangents-count-local-lift
The unit step gives a map that is a cocycle only modulo the corner entry $(1,3)$. The local count, applied to the matrices, to the corner and to the quotient, corrects such a map to a true cocycle after multiplying it by a factorial that depends only on a bound for the equivariant maps from the corner to the cyclotomic character.

% before: sha-automorphy-unramified-tangents-count-local
At a place above $\ell$, the local Euler characteristic formula gives many upper triangular cocycles, and the vanishing of their diagonal on inertia costs a bounded power of $|\mathcal O_1/\ell^k|$. At the place of the unit step, the class $\beta$ produces a cocycle whose diagonal moves, and this gives one more factor $|\mathcal O_1/\ell^k|$, up to a constant.

% before: sha-automorphy-unramified-tangents-count-adjoint
The global count is applied to the adjoint representation modulo $\ell^k$, and four of its terms need bounds. The action has open kernel, and there is no equivariant map to the cyclotomic character. Each complex conjugation fixes at most $|\mathcal O_1/\ell^k|^3$ matrices, and the normalizations at $c_\Gamma$ and $g_0$ cost at most a factor $|\mathcal O_1/\ell^k|^8$.

% before: sha-automorphy-unramified-tangents-count-selmer
Putting the local bounds into the inequality of Greenberg and Wiles, the factors for the places and for the infinite places cancel. The extra factor at the place of the unit step remains, and gives the growth like $|\mathcal O_1/\ell^k|^2$.

% before: sha-automorphy-unramified-tangents-count-tangent
A cocycle with these local conditions deforms the triangulation to first order modulo $\ell^k$, and the deformation is a linked point over the ring of dual numbers with coefficients modulo $\ell^k$. The universal property of $A_3$ turns it into a tangent vector that kills the characters on inertia.

% before: sha-automorphy-unramified-tangents-count
The triangular forms at the places above $\ell$, with the unit step at $v$, define local conditions for cocycles modulo $\ell^k$. The Greenberg-Wiles formula and the local Euler characteristic formulas count these cocycles, the count grows like $|\mathcal O_1/\ell^k|^2$, and each cocycle gives a tangent vector modulo $\ell^k$.

% before: sha-automorphy-unramified-tangents
On the other side, a linked point that is potentially crystalline with two diagonal characters related by the cyclotomic character has two independent tangent vectors. The contradiction with the previous lemmas proves (d).

% before: sha-residually-reducible-automorphy
To bound the group over $\mathbf Q$ by congruences between automorphic forms, every ordinary deformation of the reducible representation must be automorphic. The stable points come from cuspidal representations of $GL_3$ over the imaginary quadratic field, and the endoscopic points from newforms.

% before: sha-host-datum
The bounds for the character pieces depend on the prime $\ell$ and the twist $k$ in different ways. A finite set of bad primes, determined by $\psi$ and by an auxiliary imaginary quadratic field, is treated separately, as are the twists with $\ell-1\mid2k$.

% before: sha-cases
When the auxiliary primes $q$ and $d$ are large, every character piece of a host field falls into one of the cases. Outside the bad primes, $\ell$ also does not divide the order of the characters, so the decomposition of the plus part is exact.

% before: sha-siegel-datum
Cases II and III are bounded through the points of a deformation problem for $GSp_4$: four-dimensional representations of $\operatorname{Gal}(\bar{\mathbf Q}/\mathbf Q)$, symplectic up to a twist, that are congruent to a sum of four characters. The characters are built from the character piece, the twist and an auxiliary odd character of $2$-power conductor. The datum fixes these characters, the places and Galois elements that normalize the points, and the local conditions: ordinary at $\ell$, and trivial on part of the inertia at the other primes where the character piece has $\ell$-power ramification.

% before: sha-flagged-residues
Several arguments below separate the two blocks of a point of a Siegel datum by evaluating its diagonal characters at a single Galois element. The conditions of the datum make the residual characters of the two blocks distinct at the Frobenius elements $\sigma_p$, and also at $\phi$, whose residual image is that of $\sigma_\ell$. Complex conjugation acts on the diagonal representation by a fixed pattern of signs.

% before: sha-flagged-block-idempotent
A block diagonal matrix whose diagonal residues are distinct across the blocks generates the idempotent of the first block. The proof applies the Cayley-Hamilton theorem to each block in turn. The lemma is used twice: for the span of a Siegel point, and for the flag at $\ell$.

% before: sha-flagged-idempotents
The shape of a point is read off from its span, once that span contains the diagonal matrix units. The normalization at $\phi$ gives the idempotent of the first block, the normalization at $g_0$ splits it into two diagonal matrix units, and the symplectic relation gives the other two.

% before: sha-flagged-relation
When the diagonal matrix units lie in the span, the span is described by the entries that are somewhere nonzero. This relation on the indices is a preorder compatible with the symplectic involution, and it detects irreducibility and the span of the first blocks.

% before: sha-flagged-cocycles
An off-diagonal entry of a point, divided by a diagonal entry, gives a cocycle with values in a residue field for one of the characters $\alpha^{\pm1}$, $\beta^{\pm1}$ and $\theta^{\pm1}$. The last condition of a Siegel datum makes such a cocycle a coboundary once it is one locally at the primes of $T$. At the primes of $S_b$ the local condition holds automatically, by local duality.

% before: sha-flagged-common-line
The flag at $\ell$ gives a common eigenvector for a two-by-two corner of a triangular representation, provided the corner is cut off from the other entries. In the proof of Proposition~\ref{sha-siegel-points} this turns the ordinary condition into the local condition at $\ell$ for the cocycles above.

% before: sha-flagged-extension-cocycle
A block diagonal point becomes a non-split block extension when it is twisted by a nonzero cocycle with values in the unipotent part, ordinary at $\ell$ and trivial on inertia elsewhere. Such cocycles exist by the Greenberg-Wiles formula: complex conjugation acts by $-1$, so the local condition at $\ell$ forces the Selmer groups modulo $\mathfrak m_A^n$ to grow with $n$. A compactness argument then gives a cocycle with values in the lattice.

% before: sha-counts-vd-model-torus-characters
A two-dimensional torus acts on deformation rings by conjugation, and the weights of this action are read off by evaluating at a universal torus element. Over power series in one variable per basis vector of the weight lattice, the characters of distinct weights are linearly independent over every coefficient ring. This independence separates the components of each weight.

% before: sha-counts-vd-model-torus-action
When the kernel of a presentation is homogeneous for the weights of its variables, the torus acts on the ring itself. Each element of the part of weight $w$ is multiplied by the character of weight $w$. The action is built on the power series ring and passes to the quotient because the kernel is homogeneous.

% before: sha-counts-vd-model-torus-grading
Conversely, a torus action on a ring, given on generators of its maximal ideal, makes a presentation by these generators homogeneous. The action is tested at the universal torus element modulo each power of the maximal ideal, where the characters of distinct weights can be separated. For a universal lift this applies to the entries of the lift, so that each entry lies in the part of its weight.

% before: sha-flagged-universal
Finiteness of a Siegel datum is a statement about all of its points at once, so it is tested on the universal point. Conjugation by a two-dimensional torus of diagonal similitudes acts on the universal ring through the weights of the entries, and the products of entries of weight zero lie in its part of weight zero.

% before: sha-flagged-weight-zero
If the part of weight zero is not finite over $\mathcal O$, it has a point of characteristic $\ell$ that does not vanish on its maximal ideal. Rescaling by the torus extends this point to the whole universal ring. The result is a point of characteristic $\ell$ at which some product of entries of weight zero is a nonzero element of the maximal ideal.

% before: sha-flagged-relation-cases
The points of a Siegel datum are sorted by the pattern of entries that vanish identically, recorded as a relation on the four indices. The symplectic form and a cocycle argument at $\ell$ restrict the patterns that can occur. A finite check then leaves three: irreducible, diagonal, and block triangular.

% before: sha-flagged-relation-cases-argument
The finite check can also be done by hand. The argument sorts the indices into classes of indices related in both directions, which the symmetry $i\mapsto5-i$ permutes.

% before: sha-siegel-tower-ambient
The points of a Siegel datum are studied over valuation rings of algebraically closed fields, of characteristic $0$ or $\ell$, through flagged representations: representations congruent to $r_0$ with a flag that makes them upper triangular at $\ell$. Purity of the Frobenius eigenvalues is the property that rules out the endoscopic points.

% before: sha-siegel-tower-towers
Automorphy of a flagged representation is proved after restriction to a totally real abelian field $F$, inside the CM field $L=FK$. The conditions on such a tower control the residual characters on $\operatorname{Gal}(\bar{\mathbf Q}/L)$ and the local behaviour at $2$ and at $\ell$. Layers of the cyclotomic $\mathbf Z_\ell$-extension make the degree of $F$ large while keeping these conditions.

% before: sha-siegel-tower-pairs
Over $F$, a representation with flags at the places above $\ell$ is a flagged pair. The residual ones are lower triangular with a prescribed diagonal, and they are joined to each other by families over a polynomial ring; chains of such links carry automorphy from one point to another.

% before: sha-siegel-tower-space
The flagged pairs over $F$ have a universal family over a ring $A$ of finite type over a complete local ring $B$, which is generated by characteristic polynomials and inertial characters. The modularity of a point is expressed by an ideal of $B$ that lies in the minimal primes of $A$ through which the point factors.

% before: sha-siegel-tower-ideal
A modularity ideal plays the part of the kernel of a map from $B$ to a Hecke algebra: the points that kill it are pure and automorphic. Two further properties, at seeds and along curves in characteristic $\ell$, come from patching. The cohomology bounds control the dimensions that enter the propagation of the ideal.

% before: sha-siegel-step-towers-input
In case II the first tower is cut out by characters of $G$. Three quadratic characters, ramified at three auxiliary primes, and a character of $\ell$-power order adjust the residues of $\Psi$ and the field $K$ so that the fixed field is totally real, abelian and of even degree. The data are bundled once, since several lemmas use them together.

% before: sha-siegel-step-towers-quadratic
The quadratic characters come from Legendre symbols at a prime chosen in an arithmetic progression. The congruences make the primes of a given finite set split, and fix the sign at complex conjugation.

% before: sha-siegel-step-towers-kappa-dirichlet
The auxiliary character and the elements of the tower input are read off from Dirichlet characters, through the action of $G$ on roots of unity. A Dirichlet character splits into components at the primes dividing its modulus, and the inertia group at a prime sees only the component at that prime.

% before: sha-siegel-step-towers-kappa-units
Values of characters are compared through their residues. Roots of unity of order prime to $\ell$ are determined by their residues, and the characters trivial on the kernel of $\omega$ are the powers of $\omega$.

% before: sha-siegel-step-towers-kappa-dyadic
The element where $\xi^4$ is not $1$ is found by bounding the conductor of $\xi$. A primitive character of $2$-power conductor that kills a power of an odd number bounds that conductor by the $2$-adic valuation of the power.

% before: sha-siegel-step-towers-kappa-exponent
For that bound, a power of $\psi$ with a controlled $2$-part must be trivial. The components of $\psi$ at odd primes have order dividing $2o_{\psi_0}$, as $\psi$ is trivial on the subgroup $H_h$ of the host datum.

% before: sha-siegel-step-towers-kappa
The remaining data come from the Dirichlet characters of the datum. The auxiliary character removes the part of $\psi$ of $\ell$-power order, and the elements give witnesses for the conditions of a fit tower at $\ell$ and at $2$.

% before: sha-siegel-step-towers-subgroup
The subgroup of a tower input is the Galois group of the totally real field of the first tower. Since the three quadratic characters are ramified at distinct primes, their values can be adjusted one at a time without changing the other characters.

% before: sha-siegel-step-towers-elements
Three elements of the subgroup witness most of the conditions of a fit tower. They are found by adjusting a complex conjugation, an inertia element at $2$ and the element of the tower input with $\bar\Psi=-1$.

% before: sha-siegel-step-towers-make
A tower needs, besides its fields, one place above $\ell$ in each orbit of $G_F$, chosen to induce on $K$ the place of $v_\ell$. An element of $G_F$ that moves $K$ makes this choice possible.

% before: sha-siegel-step-towers-three-host
In case III the character $\psi$ has odd order, and the prime $\ell$ is outside the bad primes of $\psi_0$. This makes $\ell$ prime to every modulus in sight, so the Dirichlet characters of the datum are unramified at $\ell$.

% before: sha-siegel-step-towers-three-fields
In case III the field $F$ is cut out by a subgroup on which the residues of the characters are tied together, cut down by two quadratic characters. The quadratic characters supply the imaginary quadratic field $K$ and the even degree of $F$.

% before: sha-siegel-step-towers-three-elements
The witnesses for the conditions of a fit tower in case III come from elements on which $\xi$ and $\psi$ take prescribed values. They are found by the Chinese remainder theorem applied to the moduli of $\xi$ and $\psi$, which are coprime.

% before: sha-siegel-step-towers
Towers with these conditions exist for every relabelling that satisfies (a), the conditions pass to larger towers, and there are deep layers of any degree. Depth makes $\psi$ of order prime to $\ell$ over $F$.

% before: sha-siegel-step-labels-flag
The flag of a flagged representation is seen through the Frobenius element at $\ell$, at which the residual characters of the two blocks differ. The identity between the flag and the residual diagonal representation forces two entries of the first column of the flag to vanish residually and decides which relabelling fits.

% before: sha-siegel-step-labels-hensel
A matrix over a henselian local ring that is residually diagonal with distinct diagonal entries has eigenvalues lifting those entries. Its eigenvectors are read off from the adjugate of the matrix minus an eigenvalue, and they lift the basis vectors.

% before: sha-siegel-step-labels-diagonal
The normal form also needs $r(g_0)$ to be diagonal. A symplectic matrix that is residually diagonal with distinct diagonal entries is diagonalized, over a complete local ring, by a symplectic matrix congruent to $1$.

% before: sha-siegel-step-labels
The first step puts a flagged representation in a normal form over $\mathbf Q$, a rational flagged pair close to the diagonal representation. A relabelling of the characters may be needed to place the flag.

% before: sha-siegel-step-lattice-shapes
A change of lattice is a conjugation by a diagonal matrix, which scales each entry by a character of its weight. Which roots survive in the reduction is recorded by the tight set, and only six configurations of tight sets can occur after a suitable conjugation.

% before: sha-siegel-step-lattice-vertex-program
The diagonal matrix of the next lemma is $(p^x,p^y)$ for a square root $p$ of a uniformizer, and the integers $x$ and $y$ solve an integer program in the valuations of the entries. Taking $y$ largest, and then $x$ largest, keeps the conjugated entries integral. The entries that become units then have weights of one of the six shapes.

% before: sha-siegel-step-lattice-vertex
The diagonal matrix is chosen from the valuations of the entries of the pair. For this choice the reduction keeps enough entries for its tight set to take one of the six shapes.

% before: sha-siegel-step-lattice-conjugate
Conjugation by a diagonal matrix keeps the conditions of a flagged pair, since it preserves products, the symplectic involution, diagonal entries and vanishing entries. Integrality of the conjugated pair is part of the choice of the matrix.

% before: sha-siegel-step-lattice-vanishing
An entry of the reduction whose weight is an odd root, and which is not fed by other entries, is a cocycle for the character of its weight. The cocycle vanishing for the Siegel datum kills it.

% before: sha-siegel-step-lattice-exclusion
The vanishing of such entries rules out tight sets that contain $(-1,-1)$ without $(-1,0)$, and $(-1,0)$ without $(-1,1)$ or $(0,1)$. Together these exclude five of the six shapes.

% before: sha-siegel-step-lattice-reading
Once the tight set has the first of the six shapes, the reduction itself has the shape asked for by the chains. When it has no entry of weight $(0,-1)$, the vanishing applied to the entry $(2,3)$ gives the remaining conditions of lattice shape.

% before: sha-siegel-step-lattice
A change of lattice gives the rational flagged pair a residual image of a fixed lower triangular shape. The chains of Lemma~\ref{sha-siegel-step-chains} start from residual images of this shape.

% before: sha-siegel-step-restriction-matrices
A flag at a place $v$ above $\ell$ other than $v_\ell$ comes from the flag at $v_\ell$ moved by an element of $G$. Finding it is a factorization of a symplectic matrix as a lower unipotent one times an upper triangular one, whose pivots are units when the residual representation has lattice shape.

% before: sha-siegel-step-restriction
Restriction to $\operatorname{Gal}(\bar{\mathbf Q}/F)$ turns a rational flagged pair into a flagged pair for the tower, with flags at all places above $\ell$. Its residual image lies over the residual image over $\mathbf Q$.

% before: sha-siegel-step-seeds-final
A seed is first built over $\bar{\mathbf Q}_\ell$, with integral entries and congruences to the residual characters. Its entries then lie in the ring of integers of a finite extension of $\mathbf Q_\ell$, over which it is an integral seed, and the congruences become identities between residues.

% before: sha-siegel-step-seeds-induced-datum
The seed over $\bar{\mathbf Q}_\ell$ is the four-dimensional representation of $G$ induced from a two-dimensional representation of $G_{k_0}$, for the real quadratic field $k_0$ of condition (p), twisted on $G_F$ by a character close to $1$. Several lemmas use the same data, so they are collected in one definition, with the induced representation written in an explicit basis through the matrices $M(\gamma,A,B)$.

% before: sha-siegel-step-seeds-induced-characters
Once $k_0$ is fixed, the residues of the relabelling characters on $G_F$ are determined by those of $\xi$ and of the quadratic character of $k_0$. These residues fix the diagonal of the induced representation modulo the maximal ideal.

% before: sha-siegel-step-seeds-induced-hilbert-form-ordinary
The Hilbert modular form is produced together with its Galois representation and the local data that make it ordinary at the primes above $\ell$. These data are an eigenvector of $U_w$ at each prime where $\xi$ ramifies or that lies above $\ell$, an unramified vector at the other primes, and weights at the infinite places that induce the primes above $\ell$.

% before: sha-siegel-step-seeds-induced-hilbert-form-exists-holomorphic
The form is built adelically, as a function on $GL_2(\mathbf A_E)$ with a level at each prime and a weight at each infinite place. The weight is read off from the action of the rotations at that place, and holomorphy is the vanishing of the lowering operator there. The Hecke operators are the operators $U_w$ of the ordinary form datum, scaled by $\operatorname N(w)^{k/2-1}$ for the parallel weight $k$.

% before: sha-siegel-step-seeds-induced-hilbert-form-exists-eisenstein-expansion
The eigenform is found through congruences between Fourier coefficients. A holomorphic form is recorded by its Fourier expansion along the totally positive elements of $E$, with coefficients indexed by integral ideals, and the operators $U_w$ act on these coefficients by an explicit formula.

% before: sha-siegel-step-seeds-induced-hilbert-form-exists-eisenstein-series
The starting point is an Eisenstein series of parallel weight one. It exists because $\omega$ is odd, and its coefficients are divisor sums of the character $\omega$ on ideals prime to its conductor.

% before: sha-siegel-step-seeds-induced-hilbert-form-exists-eisenstein-congruent
To move from weight one to a large weight without changing coefficients modulo the maximal ideal, the Eisenstein series is multiplied by a form whose expansion is congruent to $1$. Such forms exist in every parallel weight divisible by $(\ell-1)\ell^N$ for a suitable $N$.

% before: sha-siegel-step-seeds-induced-hilbert-form-exists-eisenstein-raise
At the primes above $\ell$ the Eisenstein series has to be made an eigenvector of $U_w$ with a prescribed eigenvalue. This uses the translate of a form by a uniformizer at one prime, which raises the level there by one and shifts the coefficients by that prime.

% before: sha-siegel-step-seeds-induced-hilbert-form-exists-eisenstein-stabilized
Subtracting a multiple of the translated form removes one of the two roots of the Hecke polynomial at a prime above $\ell$. Doing this at every prime above $\ell$ gives the eigenvalues $e(w)$ there, and the constant term vanishes once one of these primes lies in $P$.

% before: sha-siegel-step-seeds-induced-hilbert-form-exists-eisenstein-product-holomorphic
The weight is raised by multiplying by a form of level $0$ and trivial central character. The product of two holomorphic forms is again one, with the weights added: the conditions at infinity follow from the Leibniz rule, and the finiteness under the Capelli operators from their eigenvalues on the product.

% before: sha-siegel-step-seeds-induced-hilbert-form-exists-eisenstein-product-finite
The Fourier coefficients of a product are convolutions over decompositions $\xi=\eta+(\xi-\eta)$ into totally positive elements. Only finitely many such decompositions carry nonzero coefficients, which makes the congruence of the convolution an ultrametric estimate over a finite sum.

% before: sha-siegel-step-seeds-induced-hilbert-form-exists-eisenstein-product
Multiplying by the form congruent to $1$ raises the weight from $1$ to $k+1$. The coefficients of the product are congruent to those of the weight one form.

% before: sha-siegel-step-seeds-induced-hilbert-form-exists-eisenstein-lift-finite-rotation
The weight condition fixes how a holomorphic form changes under the rotations at infinity. With the signs at the infinite places, the rotations make up the whole maximal compact subgroup at infinity, so the form has one type under that subgroup.

% before: sha-siegel-step-seeds-induced-hilbert-form-exists-eisenstein-lift-finite-harish-chandra
Cusp forms with a fixed level, a fixed type under the maximal compact subgroup at infinity and fixed eigenvalues of the Capelli operators form a finite-dimensional space. This is a form of Harish-Chandra's finiteness theorem, stated for $GL_a$ over a totally real field.

% before: sha-siegel-step-seeds-induced-hilbert-form-exists-eisenstein-lift-finite-cusp
Finite dimensionality is proved first for cusp forms. A holomorphic form has a fixed type under the rotations at infinity and fixed eigenvalues of the Capelli operators, and the cusp forms with these data and a fixed level form a finite-dimensional space.

% before: sha-siegel-step-seeds-induced-hilbert-form-exists-eisenstein-lift-finite-iwasawa
Every element of $GL_2(\mathbf A_E)$ is a scalar times an upper triangular matrix of the form $a(y,x)$ times an element of the maximal compact subgroup. The constant term behaves simply under each of these factors, which is how a finite set of points comes to control it.

% before: sha-siegel-step-seeds-induced-hilbert-form-exists-eisenstein-lift-finite-points
A holomorphic form is a cusp form once its constant term vanishes. For forms of a fixed level this is decided at finitely many points, since the set where the constant term vanishes is stable under translations that multiply the constant term by nonzero factors.

% before: sha-siegel-step-seeds-induced-hilbert-form-exists-eisenstein-lift-finite
The lift works inside the space of holomorphic forms of the given weight, level and central character. Linear algebra over $\bar{\mathbf Q}_\ell$ needs this space to be finite-dimensional.

% before: sha-siegel-step-seeds-induced-hilbert-form-exists-eisenstein-lift-hecke-cosets
The operator $U_w$ is computed through the left cosets of its double coset. There are $\operatorname N(w)$ upper triangular cosets, and one more when the level at $w$ is zero.

% before: sha-siegel-step-seeds-induced-hilbert-form-exists-eisenstein-lift-hecke-translation
The upper triangular cosets move the variable $x$ of a Fourier expansion. A holomorphic form is invariant under integral unipotent translations, and a rational unipotent factor can be moved to the left, where the form is invariant under $GL_2(E)$.

% before: sha-siegel-step-seeds-induced-hilbert-form-exists-eisenstein-lift-hecke-characters
The translations of $x$ by elements of $E$ multiply the Fourier terms by additive characters, which are summed over the residues at $w$. These sums are $\operatorname N(w)$ or $0$ according to the position of the element relative to the inverse different.

% before: sha-siegel-step-seeds-induced-hilbert-form-exists-eisenstein-lift-hecke-values
With these facts, the value of $h_w\phi$ at a point of the Borel subgroup is a sum of values of $\phi$ at such points. The representatives of the residues are scaled by an element of $E$ of the right valuation at $w$, so that each coset becomes a translation of $x$ by an element of $E$.

% before: sha-siegel-step-seeds-induced-hilbert-form-exists-eisenstein-lift-hecke-coefficients
The character sum turns the Fourier coefficients of these translates into the coefficients $a(w\mathfrak m)$. The sum vanishes when the idele does not generate an ideal of the right form, which removes the other terms.

% before: sha-siegel-step-seeds-induced-hilbert-form-exists-eisenstein-lift-hecke
The congruences are stated on Fourier coefficients, so the operators $U_w$ have to be read on coefficients. On a form with a Fourier expansion, the normalized operator $U_w$ acts on the coefficients by the formula $T^k_w$, and these operators commute.

% before: sha-siegel-step-seeds-induced-hilbert-form-exists-eisenstein-lift-unique-exponentials
A Fourier expansion is an absolutely convergent sum of exponentials in the variable $x$, with frequencies the totally positive elements. Averaging over translations of $x$ removes the terms one frequency at a time, so a sum that is constant in $x$ has no nonzero frequency.

% before: sha-siegel-step-seeds-induced-hilbert-form-exists-eisenstein-lift-unique-rotation
The weight condition says that the rotations at infinity act on a holomorphic form by a character. In particular they do not change its absolute value.

% before: sha-siegel-step-seeds-induced-hilbert-form-exists-eisenstein-lift-unique-approximation
Up to $GL_2(E)$ on the left and an open subgroup on the right, every element of $GL_2(\mathbf A_E)$ is an element at infinity times a diagonal matrix $\operatorname{diag}(y,1)$ at the finite places. This is strong approximation for $SL_2$, applied once the determinant is moved into $\operatorname{diag}(y,1)$.

% before: sha-siegel-step-seeds-induced-hilbert-form-exists-eisenstein-lift-unique
To pass from coefficients back to forms, a form must be determined by its coefficients. In positive weight the constant term carries no further information, and a form whose coefficients vanish is zero.

% before: sha-siegel-step-seeds-induced-hilbert-form-exists-eisenstein-lift-deligne-serre
The lift itself is linear algebra over $\bar{\mathbf Q}_\ell$. A vector that is an eigenvector modulo the maximal ideal for a commuting family of operators preserving integrality is replaced by a true common eigenvector with congruent eigenvalues.

% before: sha-siegel-step-seeds-induced-hilbert-form-exists-eisenstein-lift
A form whose coefficients are congruent to those of an eigenform need not be an eigenform. The lemma of Deligne and Serre, applied to the finite-dimensional space of forms of the given weight and level, gives a true eigenform with congruent eigenvalues.

% before: sha-siegel-step-seeds-induced-hilbert-form-exists-eisenstein-noncuspidal-constant-term-periodic
The constant term is an integral over a fundamental domain for $E$ in $\mathbf A_E$ of a function that is periodic under $E$. Such integrals converge for continuous functions, since $\mathbf A_E/E$ is compact, and they do not depend on a translation or on an automorphism preserving $E$. They can also be differentiated under the integral sign along a parameter.

% before: sha-siegel-step-seeds-induced-hilbert-form-exists-eisenstein-noncuspidal-constant-term-matrix
The derivative of the constant term along the diagonal at infinity is computed one infinite place at a time. At a place, the weight, the lowering condition and the invariance under the unipotent radical give four linear relations among the derivatives in the directions of the matrix units, and these fix the derivative along the diagonal.

% before: sha-siegel-step-seeds-induced-hilbert-form-exists-eisenstein-noncuspidal-constant-term
An eigenform that is not a cusp form has a nonzero constant term along the unipotent radical. The constant term is invariant under the unipotent radical and under $E^\times$ in the torus, and the diagonal matrices at infinity act on it through the weight.

% before: sha-siegel-step-seeds-induced-hilbert-form-exists-eisenstein-noncuspidal-characters
The constant term is a function on the ideles in the first diagonal entry, invariant under a subgroup of finite index by the finiteness of a ray class group. Projecting to one character of the finite quotient gives a function on which the Borel subgroup acts through a pair of Hecke characters.

% before: sha-siegel-step-seeds-induced-hilbert-form-exists-eisenstein-noncuspidal-local
At each prime the projected constant term restricts to a vector in a principal series of $GL_2(E_w)$ that is fixed by $K_1(w^c)$ and is an eigenvector of $U_w$. The possible eigenvalues are read off from the values of the vector at $1$ and at the Weyl element.

% before: sha-siegel-step-seeds-induced-hilbert-form-exists-eisenstein-noncuspidal
The eigenform obtained in this way could be an Eisenstein series. An eigenform that is not a cusp form has the eigenvalues of an Eisenstein series attached to two characters of finite order, which the primes $w_1$ and $w_2$ rule out.

% before: sha-siegel-step-seeds-induced-hilbert-form-exists-eisenstein-chebotarev
The congruences at almost all primes are turned into congruences at every idele. This uses the Galois characters of the two characters and the density of Frobenius elements.

% before: sha-siegel-step-seeds-induced-hilbert-form-exists-eisenstein
The construction starts from the Hecke character $\omega$ that class field theory attaches to $\xi$. A cuspidal eigenform of parallel weight $k+1$ and central character $\omega$ is found whose eigenvalues are congruent to those of an Eisenstein series of $\omega$. At the primes above $\ell$ these eigenvalues are $\omega_w(\varpi_w)$ or $1$, according as the prime lies in $P$ or not.

% before: sha-siegel-step-seeds-induced-hilbert-form-exists-weight-model
The weight is raised by comparing traces of Hecke operators in two weights. Traces are computed on a classical model of the holomorphic forms, as functions of a point in a product of half-planes and of a finite adelic matrix, on which the Hecke operators are sums of right translations by representatives.

% before: sha-siegel-step-seeds-induced-hilbert-form-exists-weight-model-facts
The comparison needs the model to match the adelic holomorphic forms in both directions. It also needs the spaces in the model to be finite-dimensional and stable under the Hecke operators, so that their traces are defined.

% before: sha-siegel-step-seeds-induced-hilbert-form-exists-weight-traces
The arithmetic input is a congruence between traces. Once the operators at the primes above $\ell$ are normalized by $\ell^t$ in the raised weight, the traces of the Hecke operators in weights $k_0$ and $k_0+2t$ agree modulo a power of $\ell$ that grows with $t$ and with the number of factors at these primes.

% before: sha-siegel-step-seeds-induced-hilbert-form-exists-weight-linear-systems
Congruent traces are to be turned into congruent eigenvalues, so the traces are written in terms of eigenvalues. For a commuting family of operators, the trace of every product is a sum over the common eigenvalue systems, with positive integer multiplicities that add up to the dimension.

% before: sha-siegel-step-seeds-induced-hilbert-form-exists-weight-linear-powers
Raising a number $1+x$ with $\nu(x)<1$ to the power $\ell$ multiplies its distance from $1$ by at most $\max(\nu(\ell),\nu(x))$. Iterated, this makes $(1+x)^{\ell^a}$ as close to $1$ as needed, which gives the small error terms of the argument that follows.

% before: sha-siegel-step-seeds-induced-hilbert-form-exists-weight-linear-factors
The eigenvalue systems are separated by multiplying the products of eigenvalues by polynomials in them. When the coefficients of these polynomials are integral, a bound on the sums of products of eigenvalues passes to the sums of the new products.

% before: sha-siegel-step-seeds-induced-hilbert-form-exists-weight-linear-integral
Bounded traces force integral eigenvalues. A polynomial that vanishes on every other eigenvalue system isolates one system, and its powers then stay bounded only if its eigenvalues are integral.

% before: sha-siegel-step-seeds-induced-hilbert-form-exists-weight-linear
Congruent traces give congruent eigenvalues. This is a statement about two finite-dimensional spaces with commuting operators, and it needs the congruence to be finer than the absolute values of the integers up to the dimension.

% before: sha-siegel-step-seeds-induced-hilbert-form-exists-weight
The ordinary form datum asks for the weights $(t,-k-t)$ at the infinite places that induce the primes of $P$. Raising the weight by $2t$ at these places, with $t$ divisible by $\ell-1$, keeps the eigenvalues modulo the maximal ideal, up to a factor $\ell^t$ at the primes above $\ell$ that they induce.

% before: sha-siegel-step-seeds-induced-hilbert-form-exists-representation-casimir
The infinitesimal character of the representation generated by a holomorphic form is read off from the form itself. The Capelli operator at an infinite place is a polynomial in the derivatives of right translation, and on a holomorphic form the weight condition and the vanishing of the lowering operator determine its value. The rotations at that place give the sign of the central character there.

% before: sha-siegel-step-seeds-induced-hilbert-form-exists-representation-constituent
A cusp form need not lie in one cuspidal automorphic representation, but it lies in a finite sum of them. Projecting it to one summand keeps every eigenvalue relation that it satisfies for operators preserving the stable spaces, such as translations, the operators $U_w$ and the Capelli operators.

% before: sha-siegel-step-seeds-induced-hilbert-form-exists-representation
A cuspidal eigenform of this kind gives a cuspidal automorphic representation, regular algebraic and polarized after a twist by a power of $|\det|$. Its local components carry the eigenvalues of the operators $U_w$, and at the primes of level zero its Satake parameters are the roots of the Hecke polynomial of the form.

% before: sha-siegel-step-seeds-induced-hilbert-form-exists
The form is first built with determinant $\chi_{\mathrm{cyc}}^{-k}\xi$. It comes from an eigenform of parallel weight congruent to an Eisenstein series built from $\xi$, whose weight is then raised by $t$ at the infinite places that induce the primes of $P$; its eigenvalues at the primes above $\ell$ give the constants $a(w)$.

% before: sha-siegel-step-seeds-induced-hilbert-form-galois-jacquet
At a prime of positive level, the local components of an ordinary form datum are controlled by a principal series. An eigenvector of $U_w$ with nonzero eigenvalue has nonzero image in the Jacquet module, and a linear form on that module that transforms by a character of the Borel subgroup places the representation in a principal series by Frobenius reciprocity.

% before: sha-siegel-step-seeds-induced-hilbert-form-galois-component
The conditions of an ordinary form datum concern the action of $GL_2(E_w)$ on the whole space of $\pi$. That action is a sum of copies of one local component, so fixed vectors, eigenvectors of $U_w$ and the central character pass to every local component.

% before: sha-siegel-step-seeds-induced-hilbert-form-galois-central
The two characters of the principal series are tied together by the central character of $\pi$. Comparing the determinant of the attached Galois representation with the Satake parameters at almost all primes identifies the central character with the Hecke character of $\xi$ times a power of the norm.

% before: sha-siegel-step-seeds-induced-hilbert-form-galois
The Galois representation of an ordinary form datum has the local shape that the form predicts. Away from $\ell$ inertia acts through $1$ and $\xi$; above $\ell$ the representation is triangular, with the weights of the form on inertia; and the gap between the two weights above $\ell$ makes it irreducible. The same gap, and the ramification of $\xi$ at the other primes of positive level, make the local components of the form plain at those primes.

% before: sha-siegel-step-seeds-induced-hilbert-form-twist
The determinant $\chi_{\mathrm{cyc}}^{-k}\xi$ becomes $\chi_{\mathrm{cyc}}^k\xi$ after twisting by $\chi_{\mathrm{cyc}}^k$. On the automorphic side this is the twist by $|\det|^k$, which shifts the weights at the infinite places by $k$ and keeps the local components plain or unramified.

% before: sha-siegel-step-seeds-induced-hilbert-form-ordinary-twist
After the twist, the triangular form at a place above $\ell$ is compared with the characters $1$ and $\xi$ modulo the maximal ideal. Powers of $\chi_{\mathrm{cyc}}$ with exponent divisible by $\ell-1$ are congruent to $1$, and a continuous character of the decomposition group is determined modulo the maximal ideal by its values on inertia and at a Frobenius element.

% before: sha-siegel-step-seeds-induced-hilbert-form-irreducible-normal
Irreducibility on open subgroups reduces to normal subgroups. If an irreducible representation of dimension $2$ becomes reducible on a normal subgroup that contains an element with distinct eigenvalues, every element that moves the stable line has trace $0$, and the congruences for the characteristic polynomials then constrain $\Phi$.

% before: sha-siegel-step-seeds-induced-hilbert-form-irreducible
Inside every open subgroup of $G_E$, the inertia group at a place above $\ell$ has an element on which the ordinary representation has distinct eigenvalues. So a reducible restriction to an open normal subgroup would force $\xi^2$ to be congruent to $1$, which the hypothesis on $\xi$ excludes.

% before: sha-siegel-step-seeds-induced-hilbert-form
Over a totally real field of even degree in which $\ell$ splits completely, an odd character $\xi$ unramified above $\ell$ is the residual partner of $1$ in the Galois representation of a cuspidal Hilbert modular form. The set $P$ chooses, prime by prime above $\ell$, which of the two ordinary weights the form has there.

% before: sha-siegel-step-seeds-induced-hilbert
The two-dimensional representation comes from a Hilbert modular form over $k_0$ whose residual representation has constituents $1$ and $\xi$. The form is ordinary at the two places of $k_0$ above $\ell$, with different weights at the two.

% before: sha-siegel-step-seeds-induced-ribet-lattice
An irreducible two-dimensional representation whose characteristic polynomials are congruent to those of two characters can be conjugated to be upper triangular modulo the maximal ideal, in the spirit of Ribet's lemma. An element at which the two characters have distinct residues can be made to act diagonally, and compactness gives an entry $(1,2)$ of absolute value $1$.

% before: sha-siegel-step-seeds-induced-ribet
The residual representation of the Hilbert modular form is reducible, with constituents $1$ and $\xi$. A suitable stable lattice makes its reduction a non-diagonal upper triangular representation, which is the source of the entries of weight $(-1,0)$ or $(-1,1)$ and $(-1,-1)$ asked for in the seed.

% before: sha-siegel-step-seeds-induced-conjugate
Changing the basis of the two-dimensional representation does not change its determinant, its characteristic polynomials or its conjugacy classes on subgroups. So the properties of the Hilbert modular form survive the choice of the lattice.

% before: sha-siegel-step-seeds-induced-mix
Inducing from a subgroup of index $2$ gives a four-dimensional representation. When the square of an element outside the subgroup acts diagonally, it can be written through the matrices $M(\gamma,A,B)$, and it is then symplectic, with multiplier the determinant.

% before: sha-siegel-step-seeds-induced-square-root
The multiplier of the induced representation is $\mu$, which is congruent to $\theta$ on $G_F$. For $\ell$ odd, a character congruent to $1$ has a continuous square root congruent to $1$, and twisting by it corrects the multiplier.

% before: sha-siegel-step-seeds-induced-exists
With these inputs the data of an induction datum can be built: the field $k_0$, the Hilbert modular form, the lattice, square roots for the diagonal of $\rho((g_0^T)^2)$, the induced representation and the twist. The Hilbert modular form also supplies the automorphy of $\rho$ used later.

% before: sha-siegel-step-seeds-induced-irreducible-twist
The group $G_L\cap\ker\chi_{\mathrm{cyc}}$ contains the commutators of $G_L$, so irreducibility on it can be tested through characters of the abelian quotient. If an irreducible representation becomes reducible on a normal subgroup with abelian quotient, it is isomorphic to its twist by a nontrivial character of the quotient, by Clifford's argument.

% before: sha-siegel-step-seeds-induced-irreducible-commutators
For a representation that is irreducible on every open subgroup, such a twisting character has an open kernel. On that kernel the intertwining matrix commutes with an irreducible representation, so it is a scalar and the character is trivial.

% before: sha-siegel-step-seeds-induced-irreducible-mackey
An induced representation from a subgroup of index $2$ is a sum of two planes exchanged by an element outside the subgroup. It is irreducible when the planes are irreducible and not isomorphic, by Mackey's criterion in the case of index $2$.

% before: sha-siegel-step-seeds-induced-irreducible-planes
In the basis of the matrices $M(\gamma,A,B)$, the two planes of the induced representation are explicit, with $\rho$ acting on one of them and its conjugate by $\tau$ on the other. Mackey's criterion then reduces the irreducibility of $R$ to two statements about $\rho$.

% before: sha-siegel-step-seeds-induced-irreducible-sign
To exclude an isomorphism between $\rho$ and its conjugate, it is enough to compare them at the square of a suitable element. An isomorphism on a subgroup normalized by $s$ extends to $s^2$ when the determinants agree at $s$, as the defect at $s$ is a scalar of square $1$.

% before: sha-siegel-step-seeds-induced-irreducible-residue
The subgroup $G_L\cap\ker\chi_{\mathrm{cyc}}$ must not lie in $G_{k_0}$, whose elements of $G_F$ are those where $\bar\chi^{(0,1)}$ is trivial. If it did, $\bar\chi^{(0,1)}$ would be a quadratic character of $G_L$ permuting the residual characters and trivial where $\omega$ is, which conditions (f) and (g) of a fit tower exclude.

% before: sha-siegel-step-seeds-induced-irreducible
A seed must have irreducible restriction to $G_L\cap\ker\chi_{\mathrm{cyc}}$. For an induced representation this comes down to the irreducibility of $\rho$ on a subgroup and to the distinctness of $\rho$ from its conjugate there.

% before: sha-siegel-step-seeds-induced-flags-line
An ordinary filtration of $\rho$ at a place above $\ell$ is a line stable under the decomposition group. When $\rho$ is upper triangular modulo the maximal ideal, with diagonal residues $1$ and $x$ that differ somewhere, the residue of the eigenvalue on the line fixes its position: it is spanned by $(u,1)$ with $u$ integral, or by $(1,y)$ with $y$ in the maximal ideal.

% before: sha-siegel-step-seeds-induced-flags-matrix
An eigenvector $(u,1)$ of $A$ and an eigenvector $(1,y)$ of $B$, or the same with $A$ and $B$ exchanged, give a flag stable under $M(\gamma,A,B)$. An explicit lower unipotent matrix carries the standard flag to it, and its entries $(3,2)$ and $(3,1)$ are those asked for in condition (b) of Lemma~\ref{sha-siegel-step-seeds-final}.

% before: sha-siegel-step-seeds-induced-flags
The ordinary filtrations of $\rho$ at the places of $k_0$ above $\ell$ give flags for the induced representation at the places of $F$ above $\ell$. The lower unipotent matrices of these flags have the shape asked for in condition (b) of a seed.

% before: sha-siegel-step-seeds-induced-deletion
The twisting character $\delta$ multiplies the diagonal characters of the flags by $\delta$. Its square is, up to a character of finite order, a power of the cyclotomic character, and on the inertia groups above $\ell$ this keeps the inertial weights integral.

% before: sha-siegel-step-seeds-induced-local
The flags of $R$ and the inertial behaviour of $\delta$ combine into the local conditions of a seed for the twisted restriction. The weights come out as $2m_i+b$, with gaps of at least $4$.

% before: sha-siegel-step-seeds-induced-automorphic-published
The automorphy of the induced representation over $L$ rests on three operations on automorphic representations of $GL_2$: base change from the real quadratic field $k_0$ to a CM field, base change up a solvable extension, and automorphic induction from a quadratic extension. They are published results of Arthur and Clozel, stated here in the form used.

% before: sha-siegel-step-seeds-induced-automorphic-induction
Automorphic induction gives an automorphic representation with prescribed Satake parameters, but a Galois representation is attached to it only after a comparison of Frobenius eigenvalues. The characteristic polynomials of an induced representation are determined by those of the representation it is induced from, on the subgroup and off it, and this is all the comparison uses.

% before: sha-siegel-step-seeds-induced-automorphic-windows-plain
A bound on the absolute values of the characters in a supercuspidal support confines them to an interval whose endpoints have ratio $q_w$, so no two of them are linked. The local component is then the product of its characters, and the local Langlands correspondence sends it to a sum of characters with no monodromy; if the characters are unramified, the local component is spherical.

% before: sha-siegel-step-seeds-induced-automorphic-windows-twist
The bound on supercuspidal supports survives the operations that the passage between fields uses. Twisting by a Hecke character shifts the bound by a constant, since the absolute value of a Hecke character is a power of the norm, and strong multiplicity one passes the bound from one representation attached to a Galois representation to all of them.

% before: sha-siegel-step-seeds-induced-automorphic-windows-up
Along a cyclic extension of prime degree, the local components of a base change have as supercuspidal supports the characters of the original support composed with the norm. The bound therefore goes up such an extension, and up a solvable one by induction on the degree.

% before: sha-siegel-step-seeds-induced-automorphic-windows-real
The bound is first established over the CM quadratic extension of a totally real field. It holds at every place at which the base change is plain, by the bound for unitary generic representations, so the remaining point is plainness at every place.

% before: sha-siegel-step-seeds-induced-automorphic-windows-induction-local
An automorphic induction from a quadratic extension has at each place a supercuspidal support related, through the norm, to the supports of the induced-from representation above it. When those consist of characters that come from the smaller field up to unramified twists, the bound descends.

% before: sha-siegel-step-seeds-induced-automorphic-windows
A seed asks for a bound on the supercuspidal supports at every finite place, which is more than attachment gives. Alpöge obtains it over the real field, carries it through base change and automorphic induction, and uses it to show that the base-changed representation has Satake parameters at the places where the induction needs them.

% before: sha-siegel-step-seeds-induced-automorphic
On $G_L$ the induced representation is again induced, from the compositum of $k_0$ and $L$, where the Hilbert modular form can be base changed. Automorphic induction then makes it automorphic over $L$.

% before: sha-siegel-step-seeds-induced-residues
On elements of $G_{k_0}$ the induced representation is $M(\gamma,A,B)$ for two matrices of the Hilbert modular representation that are upper triangular modulo the maximal ideal. The integrality of a seed then becomes a computation with these matrices.

% before: sha-siegel-step-seeds-induced-kill
A seed must be trivial on the inertia elements of the problem of $T$. On such elements every character in sight is trivial, and the Hilbert modular representation is trivial by its inertial type away from $\ell$.

% before: sha-siegel-step-seeds-induced-global
The global conditions of a seed are integrality, the residues of the diagonal and of the entries outside the weights of type V, and an entry of absolute value $1$. Elements of $G_{k_0}$ give them through the explicit matrices, and the other elements of $G_F$ through the image of $g_0^T$, which is diagonal.

% before: sha-siegel-step-seeds-induced
The seed over $\bar{\mathbf Q}_\ell$ comes from a Hilbert modular form over the real quadratic field $k_0$ of condition (p). Its Galois representation, induced to $G$ and twisted on $G_F$, has the shape of a residual pair of type V, and automorphic induction makes it automorphic over $L$.

% before: sha-siegel-step-seeds
Seeds are points of the universal family over $F$ that are attached to automorphic representations. Over a fit deep tower they exist, with a residual image of type V.

% before: sha-siegel-step-chains-notions
The chains are built by allowing ramification at auxiliary primes and then enlarging the tower until that ramification disappears. Pairs are compared through their entries of weights $(0,-1)$ and $(-1,0)$, which decide whether two residual pairs can be joined by a line.

% before: sha-siegel-step-chains-lines-parabolic
The lines that link residual pairs are affine in the entries where the two pairs differ. They are homomorphisms because these entries lie in the unipotent radical of a parabolic subgroup, whose square is zero.

% before: sha-siegel-step-chains-lines
Two residual pairs of the same lower shape that agree in one weight can be joined by a line of residual pairs over $k[X]$. These lines are the links of the chains.

% before: sha-siegel-step-chains-levi
A residual pair supported in $W_{\mathrm m}$ can be replaced by one that keeps only its entries of weight $(-1,0)$. The entries $(3,1)$ of its flags can then be set freely, so as to match those of a seed.

% before: sha-siegel-step-chains-rational-flags
Restricting a residual pair to a larger field needs new flags at the new places. Over a field they are found by solving four linear equations, one for each entry of the first column and one for the entry $(3,2)$; conjugating by the moving element then keeps the flags upper triangular with the same diagonal.

% before: sha-siegel-step-chains-rational
A pair of lattice shape over $\mathbf Q$ restricts to residual pairs for the problem of the tower. The restriction keeps the shape, and it is ample.

% before: sha-siegel-step-chains-gysin-cochains
The pairs of the next lemmas are built from cocycles for the characters of two weights and a cochain linking them. The local conditions at the places of $\Sigma_T$ are written with explicit constants, which become the entries of the flags.

% before: sha-siegel-step-chains-gysin-data
The part of weight $(0,-1)$ of a residual pair supported in $W_{\mathrm r}$ is a cocycle. Its local constants are the entries $(3,1)$ of the flags.

% before: sha-siegel-step-chains-gysin-raw-module
A cocycle of weight $w$ and cochains linking it to given cocycles $y_j$ are the coordinates of one cocycle with values in a finite module. The module is an extension of the character of weight $w$ by copies of the character of weight $w+w'$, glued by the $y_j$.

% before: sha-siegel-step-chains-gysin-raw-local
At an auxiliary place where the module and the cyclotomic character are trivial, the local condition allows ramification along one line. Local duality supplies both an unramified and a ramified character of order $\ell$, and they give at least $2|M|$ classes.

% before: sha-siegel-step-chains-gysin-raw-split
The auxiliary places are taken where the module is trivial on the whole decomposition group. The Chebotarev density theorem gives infinitely many of them, above primes that split in the tower.

% before: sha-siegel-step-chains-gysin-raw-count
The cocycle is found by counting. Allowing ramification along a line at many auxiliary places makes a Selmer group larger, by the Greenberg-Wiles formula, than the finite group of cocycles unramified outside the original places, so one of its cocycles is ramified somewhere along that line.

% before: sha-siegel-step-chains-gysin-raw-cocycles
For the problem of a tower the original places are those above the primes of $S$ and the auxiliary primes. The places added by the count lie above primes that split in the tower, which become new auxiliary primes.

% before: sha-siegel-step-chains-gysin-raw
Adding auxiliary primes that split in the tower produces a locally trivial cocycle of a given weight that is not zero, together with cochains linking it to given cocycles. The cocycle may ramify at the added primes, so it belongs to the problem with the larger set of auxiliary primes.

% before: sha-siegel-step-chains-gysin-cup
The cocycles of the previous lemma can be moved by coboundaries so that they vanish at the element $g_0^T$. This is needed for the pairs built from them to be diagonal there.

% before: sha-siegel-step-chains-gysin-matrix
A cocycle of weight $(-1,0)$, one of weight $(0,-1)$ and a cochain linking them fill the entries below the diagonal of a matrix in a Heisenberg group. Multiplied by the diagonal characters, these matrices form a residual pair whose flags come from the local constants.

% before: sha-siegel-step-chains-gysin
To reach a pair with an entry of weight $(-1,0)$, ramification is allowed at auxiliary primes that split in the tower. Two given pairs are moved, keeping their parts of weight $(0,-1)$, to pairs that share their parts of weight $(-1,0)$.

% before: sha-siegel-step-chains-dies-cyclotomic
The characters of order $\ell$ that are ramified only at an auxiliary prime $p$ must be shown to die on an explicit open subgroup. An abelian extension of $\mathbf Q$ unramified outside the primes dividing $M$ lies in a field of $M$-power roots of unity, by the Kronecker-Weber theorem and the inertia groups of cyclotomic fields.

% before: sha-siegel-step-chains-dies
The auxiliary ramification is removed by enlarging the tower. Over a larger field the characters of order $\ell$ ramified only at the auxiliary primes become trivial.

% before: sha-siegel-step-chains-restrict-kill-roots
Two facts about elements close to $1$ at a place are used on inertia groups. A root of unity of order prime to the residue characteristic that is congruent to $1$ is $1$, and an automorphism that moves every element by a factor congruent to $1$ is determined on a field by a power of order prime to the residue characteristic.

% before: sha-siegel-step-chains-restrict-kill-tame
An inertia group modulo its wild elements is generated topologically by one element. The generator is built from the action on roots of an element of the valuation of the residue characteristic, and characters of exponent prime to the residue characteristic do not see the wild elements.

% before: sha-siegel-step-chains-restrict-kill
Restricting a residual pair to a larger tower must make it trivial on the inertia at the auxiliary primes. At such a prime the pair has exponent $\ell$ on inertia, so it factors through tame inertia, where a character of order $\ell$ ramified only at that prime detects it.

% before: sha-siegel-step-chains-restrict-flags
A pair over another is determined by its homomorphism once its flags are lower unipotent, since two such flags that make the same matrix upper triangular agree. The condition of being over can be tested with any element that moves the places.

% before: sha-siegel-step-chains-restrict-powers
Restricting to a larger tower keeps only the $\ell$-power elements of the smaller Galois group. An entry that behaves like a cocycle stays nonzero on such powers when the two diagonal characters differ somewhere.

% before: sha-siegel-step-chains-restrict
Residual pairs for a tower restrict to residual pairs for a larger tower once the auxiliary primes are killed. Restriction keeps supports, entries and the comparisons of weights, and a pair over a rational pair is determined by it.

% before: sha-siegel-step-chains
The residual image of a point and that of a seed are joined by a chain of linked ample residual pairs. The tower may first have to be enlarged.

% before: sha-siegel-step-bounds-iwasawa
Along the cyclotomic $\mathbf Z_\ell$-extension of an abelian field, the unramified characters of order $\ell$ are dual to the class groups modulo $\ell$-th powers. The theorem of Ferrero and Washington bounds these uniformly in the layer.

% before: sha-siegel-step-bounds-places
The local terms of the Euler characteristic formulas are indexed by the places above the bad primes. In the cyclotomic $\mathbf Z_\ell$-extension every prime has an open decomposition group, so their number stays bounded along the layers.

% before: sha-siegel-step-bounds-strict
A cocycle that vanishes on inertia restricts to an unramified homomorphism on a subgroup of index prime to $\ell$ that acts trivially. The restriction and an average over the cosets determine it, so the unramified characters of the subgroup bound the number of such cocycles.

% before: sha-siegel-step-bounds-global
The Greenberg-Wiles formula, with every class allowed at the bad places, compares the global cocycles with a dual group of classes that vanish locally. Those dual classes vanish on inertia on a subgroup of index prime to $\ell$, so the previous lemma bounds them.

% before: sha-siegel-step-bounds-base
The cohomology bounds are asked along every homomorphism from $\mathcal O$ to a field, while the counts are made over the residue field of $\mathcal O$. Linear forms over the residue field transport the bounds.

% before: sha-siegel-step-bounds
The dimension estimates of the propagation need the residual Galois cohomology over the layers of the cyclotomic $\mathbf Z_\ell$-extension to grow at most like the degree. The theorem of Ferrero and Washington gives this.

% before: sha-siegel-step-ideal-links
The modularity ideal is built from two links of the tower, whose auxiliary places lie above distinct primes. Each link gives the automorphy of points away from its own auxiliary place, and two such places together cover every place.

% before: sha-siegel-step-ideal-finite
The finiteness property of the modularity ideal comes from the finiteness of the Hecke algebra over the Iwasawa algebra. Modulo a prime containing the ideal, the generators of the deformation ring become integral over the coefficients, hence constant.

% before: sha-siegel-step-ideal-setting
Each link contributes the ideal generated by the image of its Hecke ideal. The properties of a modularity ideal are checked one link at a time, on the universal ordinary determinant ring of the link.

% before: sha-siegel-step-ideal-patching-arc
The patching bound applies to maps from the universal ordinary determinant ring that are attached to nice arcs. A point of the deformation space over a power series ring over a finite field gives such an arc once the twist by $\kappa$ is pushed down to $\Gamma_0$, and its niceness comes from the hypotheses on the point and from the conditions on the tower and the link.

% before: sha-siegel-step-ideal-patching
A point of the deformation space over a power series ring that is regular at the places above $\ell$ and spans $M_4$ gives a nice arc of type I for the link. The patching bound then places every prime below it on the Hecke locus.

% before: sha-siegel-step-ideal-seeds-phi
The weights of a seed need not have the parity that the character $\kappa$ of the link matches. Twisting by a power of the second character $\beta_L$ of the link adjusts the parity, and the construction of $\varphi$ goes through for every such twist.

% before: sha-siegel-step-ideal-seeds-weights
The inertial weights of a point of the deformation space become, after the twist, the weights of a point of the universal ordinary determinant ring. Halving them is possible once the parities agree.

% before: sha-siegel-step-ideal-seeds-point
The twisted representation at a point of the deformation space inherits the properties that the Hecke ideal asks of a point: its determinant, its triangular forms above $\ell$, its triviality at the auxiliary place and a polarization. These come from the flagged pair and from the conditions on the link.

% before: sha-siegel-step-ideal-seeds-automorphy
A seed is automorphic after a twist by its own character. The twist by the characters of the link differs from it by an algebraic character, so it is automorphic too, and the weights of the seed make it windowed.

% before: sha-siegel-step-ideal-seeds-inputs-twisting
The maps $\varphi_n$ for different $n$ differ by a twist by a power of $\beta_L$. Such twists make sense on the universal ordinary determinant ring for characters of $\Gamma_0$ that are residually trivial, continuous, inverted by conjugation by $c$ and trivial where the ordinary pairs are constrained.

% before: sha-siegel-step-ideal-seeds-inputs-twist-map
The universal property turns a twist of an ordinary pair into a twist of the map that classifies it. Uniqueness comes from the fact that a determinant is determined by its characteristic polynomials at group elements.

% before: sha-siegel-step-ideal-seeds-inputs-twist-quotient
The twist property of the Hecke ideal is stated for characters of finite $\ell$-power order with values in $\bar{\mathbf Q}_\ell$. Modulo a power of the maximal ideal, a twisting character factors through a finite $\ell$-group, and the group ring of that group carries the universal twist.

% before: sha-siegel-step-ideal-seeds-inputs-twist-stable
The twist by a twisting character preserves the Hecke ideal. Modulo each power of the maximal ideal, it is read off from the universal twist over a group ring, whose coefficients lie in the ideal by the density and twist properties and by Fourier analysis on the finite group.

% before: sha-siegel-step-ideal-seeds-inputs-twist-descend
The character $\beta_L$ of the link is defined on $G_L$. It factors through $\Gamma_0$, since it is trivial on the kernel of $t_L$, and its conditions make it a twisting character there.

% before: sha-siegel-step-ideal-seeds-inputs-regular
The second fact is proved through tangent vectors. A point of a ring finite and torsion-free over a power series ring, at which every tangent vector trivial on the power series ring vanishes, has a regular local ring, so only one component of the Hecke locus passes through it.

% before: sha-siegel-step-ideal-seeds-inputs-tangent-lift
A tangent vector at a point whose representation is absolutely irreducible lifts the representation to the dual numbers. The lift inherits the ordinary relations and the polarization of the universal pair, and its first-order part extends to a cocycle of $\Gamma$ in the polarized adjoint representation.

% before: sha-siegel-step-ideal-seeds-inputs-tangent-v0-local-model
The local statement at $v_0$ is proved for representations with values in the integers of a finite extension of $\mathbf Q_\ell$. A continuous representation of a compact group over $\bar{\mathbf Q}_\ell$ has such a model, by a stable lattice, and its entries are then locally constant modulo every power of the maximal ideal.

% before: sha-siegel-step-ideal-seeds-inputs-tangent-v0-local-tame
At a prime $p$ the image of inertia in a finite $\ell$-group with $\ell\ne p$ is reached through the tame inertia, on which a Frobenius element acts through the cyclotomic character. In the form needed here, characters of inertia of order prime to $p$ transform under conjugation by the Galois action on their values.

% before: sha-siegel-step-ideal-seeds-inputs-tangent-v0-local-integral
An integral representation at $p\ne\ell$ whose reduction has a single eigenvalue, unramified, sends inertia into an $\ell$-group modulo each power of the maximal ideal. An element whose cyclotomic character is not $1$ modulo $\ell$ acts by a scalar modulo that power after raising to an $\ell$-power, and conjugation by it would then fix a nontrivial $\ell$-power root of unity, which it cannot.

% before: sha-siegel-step-ideal-seeds-inputs-tangent-v0-local
At the auxiliary place $v_0$ the residual conditions of the link force the triviality of inertia. The local statement concerns representations of a decomposition group at $p\ne\ell$ whose eigenvalues stay close to an unramified residual character.

% before: sha-siegel-step-ideal-seeds-inputs-tangent-v0
The cocycle of a tangent vector must be unramified at $v_0$ for the Bloch-Kato conditions. Doubling the representation by the cocycle reduces this to the local statement at $v_0$.

% before: sha-siegel-step-ideal-seeds-inputs-tangent-crystalline-descent
The Bloch-Kato condition above $\ell$ is first obtained on an open subgroup of the decomposition group. It descends to the whole group by averaging over the cosets of a normal open subgroup, since the coefficients have characteristic $0$.

% before: sha-siegel-step-ideal-seeds-inputs-tangent-crystalline-extension
An extension of a crystalline character by a crystalline representation is crystalline exactly when its extension class is a crystalline cocycle. The coefficients here are modules over a finite algebra, such as the dual numbers over a finite extension of $\mathbf Q_\ell$.

% before: sha-siegel-step-ideal-seeds-inputs-tangent-crystalline-positive
For representations whose Hodge-Tate weights all have the same sign, the dimension formula of Bloch and Kato and Tate's Euler characteristic formula show that every class is crystalline. The only obstruction is a quotient isomorphic to the cyclotomic character, which the residual conditions exclude.

% before: sha-siegel-step-ideal-seeds-inputs-tangent-crystalline-triangular
A triangular representation over the dual numbers, with crystalline diagonal characters spaced so that no ratio is cyclotomic, is built column by column as a chain of extensions. Each new column is a cocycle of a representation to which the dimension count applies.

% before: sha-siegel-step-ideal-seeds-inputs-tangent-crystalline-deformation
A first-order deformation of a representation is the extension of the representation by itself through its cocycle. It is crystalline exactly when the representation is crystalline and the cocycle is a crystalline cocycle of the adjoint representation.

% before: sha-siegel-step-ideal-seeds-inputs-tangent-crystalline
Above $\ell$ the Bloch-Kato condition is crystallinity. A first-order deformation that stays triangular, with weights spaced by at least two and no first-order part of the diagonal characters on inertia, is crystalline by the local theory of Bloch and Kato.

% before: sha-siegel-step-ideal-seeds-inputs-tangent-bloch-kato
The adjoint lifting theorem applies only to cocycles that satisfy the Bloch-Kato conditions. Away from $\ell$ these come from the auxiliary conditions of the ordinary pairs, and above $\ell$ from the triangular forms over the dual numbers.

% before: sha-siegel-step-ideal-seeds-inputs-tangent-constant-separating
Triangular forms over the dual numbers are compared through one element of the group algebra at which the diagonal characters differ by units. Over a local ring with enough residues, distinct residual characters provide such an element, since they are linearly independent.

% before: sha-siegel-step-ideal-seeds-inputs-tangent-constant-flags
A tangent vector whose cocycle is a coboundary has two triangular forms on each decomposition group above $\ell$, one with the universal diagonal characters and one with constant ones. When the residual diagonal characters are distinct, the diagonal characters of a triangular form over a local ring are determined by the representation, so the two agree.

% before: sha-siegel-step-ideal-seeds-inputs-tangent-constant
Once the cocycle is a coboundary, the lift to the dual numbers is a conjugate of the constant lift. Its characteristic polynomials and its diagonal characters then have no first-order part, and these topologically generate the universal ordinary determinant ring.

% before: sha-siegel-step-ideal-seeds-inputs
Two facts about the universal ordinary determinant ring finish the argument. The Hecke ideal is stable under the twists by the second character of the link, and at an automorphic point the vanishing of adjoint Bloch-Kato classes leaves no room for a component outside the Hecke locus.

% before: sha-siegel-step-ideal-seeds
At a seed the twisted representation is automorphic, so the corresponding point lies on the Hecke locus. The adjoint lifting theorem shows that no other component passes through it.

% before: sha-siegel-step-ideal-purity-roots
The eigenvalues of a point of the universal ordinary determinant ring remember the residual characters. Each eigenvalue of an element of the group algebra lies within distance less than $1$ of a lift of the value of some residual character.

% before: sha-siegel-step-ideal-purity-multiplicity
The purity statement of the Hecke ideal needs every multiplicity in its factorization to be $1$. A repeated factor would make a residual character agree with another one times the cyclotomic character, which a separating element of the group algebra rules out.

% before: sha-siegel-step-ideal-purity-frobenius
The bound on Frobenius eigenvalues comes in terms of geometric Frobenius elements at a place. Purity asks for it at every element that acts on a residue field as a power of the absolute Frobenius, which differs from a power of a Frobenius element by an element of inertia.

% before: sha-siegel-step-ideal-purity
On the Hecke locus, the points of regular weight have pure Galois representations. The inertial weights $2m$ rule out the repeated twists that the purity statement allows.

% before: sha-siegel-step-ideal-away
On the Hecke locus, an irreducible arithmetic point is automorphic over $L$ and windowed away from the auxiliary place of the link. Twisting by $\kappa$ turns the points of the deformation space into such points.

% before: sha-siegel-step-ideal-descent-third
Automorphy is carried from one imaginary quadratic field to another through their compositum, which needs irreducibility on the intersection of their Galois groups. A third imaginary quadratic field provides this whenever the representation is irreducible over the first two.

% before: sha-siegel-step-ideal-descent-split
The windowed condition is carried between imaginary quadratic fields through a third one, as automorphy is. For the windows the third field must also be split above $\ell$, so that the decomposition groups above $\ell$ do not change in the compositum; fields $\mathbf Q(\sqrt{-p})$ with $p\equiv-1$ modulo $4\ell$ provide this.

% before: sha-siegel-step-ideal-descent-conditions
The passage of the windowed condition between fields uses three local properties of the representation. Away from $\ell$ its characteristic polynomials on inertia come from characters of the decomposition group, and above $\ell$ it is triangular, with a condition on the ratios of the diagonal characters that may or may not survive in an extension.

% before: sha-siegel-step-ideal-descent-updown-local-above
Above $\ell$, the characters of a supercuspidal support are compared with the diagonal characters of a triangular Galois representation through $p$-adic Hodge theory. Each diagonal character of a de Rham triangular representation is locally algebraic, and when it extends to the decomposition group of the smaller field its smooth part comes from that field through the norm, up to an unramified character.

% before: sha-siegel-step-ideal-descent-updown-local-down
Going down a cyclic extension of prime degree needs the converse of the description of the supports of a base change. The local components below must have supports of characters, and every character below must appear, composed with the norm, in a support above.

% before: sha-siegel-step-ideal-descent-updown-local
Going down a cyclic extension, the bound on supercuspidal supports over the larger field gives the bound over the smaller one only when each character of the support comes, up to an unramified character, from the smaller field through the norm. This is where condition (I) and potential $\ell$-distinctness enter. The automorphic representations involved must also be polarized, which follows from the polarization of the Galois representation.

% before: sha-siegel-step-ideal-descent-updown-auxiliary
The descent of the windowed condition is proved one extension of prime degree at a time. When the extension has degree one, the argument still passes through an extension of prime degree, larger than the dimension of the representation. Such an extension of a CM field is obtained from a subfield of a cyclotomic field.

% before: sha-siegel-step-ideal-descent-updown-prime-index
Base change to the auxiliary field needs the representation to stay irreducible there. A representation that becomes reducible on a normal subgroup of prime index is a nontrivial twist of itself, and the determinant then bounds the order of the twisting character by the dimension.

% before: sha-siegel-step-ideal-descent-updown
Automorphy descends along solvable extensions of CM fields by base change. The windowed condition does as well, in both directions, for polarized automorphic representations that are triangular above $\ell$.

% before: sha-siegel-step-ideal-descent-across
From one imaginary quadratic field over which $r$ is windowed, the condition reaches every other one over which $r$ is irreducible. The passage goes up to the compositum and down again, twice, through a third field split above $\ell$.

% before: sha-siegel-step-ideal-descent-laws
The local properties needed for the windows hold for every representation of $G$ that extends a point of the deformation space with inertial weights $2m$. Over $L$ the triangular condition above $\ell$ holds in its stronger form at every place.

% before: sha-siegel-step-ideal-descent
Two links whose auxiliary places lie above distinct primes leave no place without the windowed condition. Automorphy over $L$ of the twist then descends to the imaginary quadratic fields over which $r$ is irreducible.

% before: sha-siegel-step-ideal
Over a fit deep tower, patching produces a modularity ideal on every deformation space with large enough coefficients. This is the input from automorphic forms in the proof of Step 1.

% before: sha-siegel-step-coefficients-extension
A family over $V$ has its entries in a complete discrete valuation ring inside $V$, but a complete local ring over a given $\mathcal O'$ must also contain the image of $\mathcal O'$. Adjoining finitely many elements to the fraction field and intersecting with $V$ gives a larger complete discrete valuation ring of the same kind.

% before: sha-siegel-step-coefficients-subring
A complete local ring over $\mathcal O'$ may have a residue field larger than that of $\mathcal O'$. The elements congruent to constants form a subring with the residue field of $\mathcal O'$, and it keeps the other properties needed of a coefficient ring.

% before: sha-siegel-step-coefficients-standard
A deformation space is built from universal rings, and their existence needs finiteness of the tangent spaces. For the problem of a fit tower this finiteness comes from Hermite's theorem globally and from local class field theory at the places above $\ell$, together with the conditions of a fit tower on the residual characters.

% before: sha-siegel-step-coefficients-universal
Before the flags are imposed, the problem asks only for a residually diagonal representation and for characters of the decomposition groups. These are pro-represented separately, and the universal ring is their completed tensor product.

% before: sha-siegel-step-coefficients-space-core-weights
The ring $B$ of weight $0$ must be determined by traces and diagonal characters. A product of matrix entries of total weight $0$ can be shortened by merging two of its factors whose weights pair negatively, until only diagonal entries remain, and those are traces.

% before: sha-siegel-step-coefficients-space-core-traces
For a representation the merging needs a subring of matrices that contains the diagonal idempotents and is stable under the symplectic adjoint. A regular diagonal element supplies the idempotents, and the multipliers make the span of the image stable.

% before: sha-siegel-step-coefficients-space-core-laurent
A representation that is residually triangular in the order given by a linear form $\tau$ on the weights becomes residually diagonal after conjugation by a diagonal matrix of powers of a variable. Over an Artinian ring this conjugation can be carried out inside a complete local ring of Laurent series.

% before: sha-siegel-step-coefficients-space-core
Conjugation by the diagonal torus grades the universal ring by the weights of the entries, and its part of weight $0$ will be the ring $B$. Three properties of this graded ring carry the construction: the topologies agree, the traces and diagonal characters determine homomorphisms on $B$, and representations that are residually triangular in the order given by a linear form on the weights still receive a homomorphism from the graded subring.

% before: sha-siegel-step-coefficients-space
The deformation space is cut out of the universal near representation by the torus that conjugates by diagonal matrices. Its invariants form the ring $B$ of traces and diagonal characters, and the entries of the representation generate $A$ over $B$.

% before: sha-siegel-step-coefficients
The argument needs deformation spaces over the rings $\mathcal O'$, and it needs chains and points over $V$ to come from residue fields and complete local rings over some $\mathcal O'$. Enlarging $\mathcal O'$ achieves both.

% before: sha-siegel-step-linked
Linked points share a minimal prime of $A$, so containing $J$ passes from one to the other. Along a chain it passes from a seed to the point, provided that at each ample residual point it passes from one minimal prime to all of them.

% before: sha-siegel-step-completion-facts-bockstein
The condition at a set $\mu_v$ is imposed by two equations only if the elements of order $\ell$ modulo commutators in the decomposition group form a cyclic group. Two independent characters of order $\ell$ would give, through the Bockstein construction, two independent classes in the second cohomology with coefficients in $\mathbf Z/\ell$, and local duality allows only one.

% before: sha-siegel-step-completion-facts-cyclic
In a residually finite commutative group, elements of $\ell$-power order have bounded order once those of order $\ell$ form a cyclic group. An element of the largest order then generates all of them.

% before: sha-siegel-step-completion-facts
The presentation of a universal lift ring is computed with the Galois cohomology of the adjoint representation of the residual pair. That representation factors through the quotient of $G_F$ by tame inertia, and the multiplier $\theta$ is trivial on the groups the problem kills and is $-1$ at complex conjugations. Local duality and the Bockstein argument make each set $\mu_v$ cyclic modulo commutators.

% before: sha-siegel-step-completion-adjoint
The adjoint representation on $\mathfrak{sp}_4$ carries a Borel subalgebra at each place above $\ell$, with a complement of dimension $4$. A complex conjugation has invariants of dimension at most $4$, and condition (e) of a fit tower leaves no equivariant linear form twisted by $\omega$.

% before: sha-siegel-step-completion-euler
The Euler characteristic formulas bound the obstructions by the tangent space. With Borel conditions of dimension $m$ at the places above $\ell$ and invariants of dimension at most $e$ at the real places, the second cohomology exceeds the first by at most $[F:\mathbf Q](e-m)$.

% before: sha-siegel-step-completion-tangent
A cohomology class with Borel conditions, normalized at $g_0^T$ by the torus and at the places above $\ell$ by the complement of the Borel subalgebra, is a lift of the residual pair to the dual numbers. Universality turns these lifts into tangent vectors of the universal ring, and the torus of $\mathfrak{sp}_4$ accounts for the term $2$.

% before: sha-siegel-step-completion-mu
The condition at the sets $\mu_v$ asks that the diagonal characters vanish there. As each $\mu_v$ is cyclic modulo commutators and the characters come in pairs with product $\theta$, two equations at each place impose it.

% before: sha-siegel-step-completion-obstruction-lift
The obstruction is computed from lifts through a small extension of the representation, the flags and the triangular forms. A continuous lift of the entries can be corrected to be symplectic with the right multiplier, since the error squares to zero.

% before: sha-siegel-step-completion-obstruction-diagonal
A lift found through a small extension must be normalized before the universal property applies. A symplectic matrix whose diagonal entries are distinct modulo the maximal ideal and that is diagonal modulo the kernel can be conjugated to a diagonal matrix by a symplectic matrix that is $1$ modulo the kernel.

% before: sha-siegel-step-completion-obstruction-flags
The flags of a lift must also be lower triangular with diagonal entries $1$. Multiplying a lifted flag on the right by an upper triangular symplectic matrix that is $1$ modulo the kernel achieves this, and it leaves the triangular forms upper triangular.

% before: sha-siegel-step-completion-obstruction-generators
The obstruction classes of the lift ring are attached to the generators of the kernel of a small extension, one class to each generator. A combination of the generators whose combined class vanishes can then be killed while the lift survives, which is what a count of relations needs.

% before: sha-siegel-step-completion-obstruction
Lifting the residual pair through a small extension meets an obstruction in the second cohomology with Borel conditions. When it vanishes, the lifted pair can be normalized and gives a section.

% before: sha-siegel-step-completion-obstruction-presentation
Mazur's comparison of a minimal presentation with the obstructions bounds the number of relations by the dimension of an obstruction space. In the form needed here the obstruction space is one for minimal generators, and the number of variables is read off from the maps to the dual numbers.

% before: sha-siegel-step-completion
At a residual point, the universal lift ring is flat over $A$ and plays the part of its completion. A count of generators and relations from Galois cohomology bounds the dimension of that ring from below.

% before: sha-siegel-step-dimensions-pinning
The third dimension bound counts parameters for the inertial characters at the places above $\ell$. By local class field theory, a character of a decomposition group with values in a local domain of characteristic $\ell$, residually trivial on inertia, is trivial on inertia once it is trivial at as many elements as the local degree.

% before: sha-siegel-step-dimensions-tangent-pairs
The first dimension bound reduces to a count of tangent vectors of the universal lift where it is supported in a set of weights. These tangent vectors are pairs of a cocycle with values in $\mathfrak{sp}_4$ and of elements that move the flags, and they make sense for any representation and flags.

% before: sha-siegel-step-dimensions-tangent-flow
A residual pair degenerates to a simpler pair when its entries are graded by integers attached to the indices and only the part of degree $0$ is kept. Taking the part of lowest degree of a tangent pair gives a tangent pair of the degenerate pair, so the degeneration can only increase the dimension of the space of tangent pairs.

% before: sha-siegel-step-dimensions-tangent-split
For a diagonal residual representation the entries of a tangent cocycle are cocycles for the characters $\chi^w$, so the cohomology bounds apply to them one at a time. If all entries but one lie in the smaller spaces $S_w$, the bound has the form $d+20C$.

% before: sha-siegel-step-dimensions-tangent-cases
Two degenerate pairs occur: the diagonal representation with trivial flags, and the diagonal representation with the flags of a pair of type V. In each, the conditions at the places above $\ell$ make all the relevant entries but one locally coboundaries.

% before: sha-siegel-step-dimensions-tangent
Where the universal lift is supported in one of the nine sets of weights, its tangent vectors are cocycles whose values are supported in the same set. Filtering them by weight reduces their number to the cohomology bounds.

% before: sha-siegel-step-dimensions-torus
Where a symplectic representation spans $M_4$ and has a regular diagonal element, its traces determine its entries up to the conjugation by the diagonal torus, which two entries fix. A lower unitriangular symplectic flag is determined by four of its entries.

% before: sha-siegel-step-dimensions-generators
The first dimension bound passes from tangent vectors to the dimension of a complete local ring. Derivations that vanish on enough elements of the maximal ideal vanish on all of it, so those elements and a uniformizer generate the maximal ideal, and Krull's height theorem bounds the dimension.

% before: sha-siegel-step-dimensions
Three dimension bounds enter the propagation: for the loci where the universal lift is reducible, for the fibres over $B$ where it is irreducible, and for the locus in $B$ where two inertial characters agree on $U_v$. The last uses the modularity ideal.

% before: sha-siegel-step-irreducible
A symplectic representation with a regular diagonal element is reducible only if it is supported in one of nine sets of weights. This turns irreducibility into the avoidance of nine ideals.

% before: sha-siegel-step-commutative
The propagation uses four facts of commutative algebra: going down for flat maps, Grothendieck's connectedness theorem, the existence of curves that avoid small closed subsets, and the parametrization of curves in characteristic $\ell$ by power series. The connectedness theorem joins any two minimal primes of the universal lift ring.

% before: sha-siegel-step-propagation
At an ample residual point, the minimal primes of the universal lift ring are joined through intersections of large dimension. A curve in such an intersection that avoids the reducible loci and the agreements of inertial characters is one to which the patching property applies.

% before: sha-siegel-step-last
The last step of Step 1 combines a seed, a chain from its residual image to that of the point, and the propagation. It produces the deformation space, the modularity ideal and the homomorphism that make the point promodular.

% before: sha-siegel-automorphy-clifford
The rigidity of an irreducible point is proved over a CM field $L$ that is abelian over an imaginary quadratic field, and the point must stay irreducible over $L$. An irreducible representation that becomes reducible on a normal subgroup with abelian quotient is isomorphic to a nontrivial twist of itself, and such a twist is incompatible with the distinct residual characters of a Siegel point.

% before: sha-siegel-automorphy-base-change
The definition of an automorphic point gives automorphy over imaginary quadratic fields, and the adjoint vanishing theorem needs it over $L$. Base change along the solvable extension $L/K$ supplies it, once the point is polarized by an algebraic character of $\mathbf Q$, which its symplectic multiplier provides.

% before: sha-siegel-automorphy-polarized
Galois representations, and with them local-global compatibility above $\ell$, are available for polarized cuspidal representations. A representation of $\operatorname{Gal}(\bar{\mathbf Q}/\mathbf Q)$ that is self-dual up to a power of the cyclotomic character times a character of finite order polarizes every cuspidal representation attached to its restriction to an imaginary quadratic field.

% before: sha-siegel-automorphy-free
The adjoint vanishing theorem needs, above $\ell$, an element of trivial cyclotomic character whose eigenvalues have no ratio that is a root of unity. Over an imaginary quadratic field it comes from the windowed condition and local-global compatibility, as in Lemma~\ref{sha-rigidity-free}, now for a polarized cuspidal representation.

% before: sha-siegel-automorphy-adjoint
The vanishing over $L^+$ is a vanishing of adjoint Bloch-Kato Selmer groups at an automorphic point of a Hida family over $L$ whose residual representation is a sum of distinct characters. Its hypotheses at the places above $\ell$, at the auxiliary place $v_0$ and on the image of $\Gamma'$ are those that a twist of an irreducible point of a Siegel datum satisfies.

% before: sha-siegel-automorphy-descent
The Hida setting sees only a quotient of $\operatorname{Gal}(\bar{\mathbf Q}/L^+)$ and a twist of the restriction of the point to $L$. Restriction, inflation from that quotient and averaging over the cosets of $\operatorname{Gal}(\bar{\mathbf Q}/L^+)$ bring the vanishing down to every polarized adjoint representation over the imaginary quadratic field.

% before: sha-siegel-automorphy-link
The vanishing over $L^+$ applies Theorem~\ref{sha-siegel-automorphy-adjoint} to a Hida family over $L$ whose residual representation is a sum of characters. A link of a tower supplies this family: a Hida setting over $L$, a character $\kappa$ whose twist turns the symplectic multiplier $\theta$ of a point into the pairing with $\chi_{\mathrm{cyc}}^{-3}$ that the setting requires, and residual characters that are the reductions of the twisted relabelling characters. Its conditions on $\kappa$ above $\ell$ and away from $\ell$ give the hypotheses of the theorem at the twisted point. It also carries a character $\beta_L$, congruent to $1$ and a square root of an odd power of the cyclotomic character on inertia above $\ell$ up to finite order, with its logarithm $t_L$; $\kappa$ is built from $\beta_L$ when the integer $k$ of the datum is odd.

% before: sha-siegel-automorphy-phi
A point of a deformation space of the tower must become a point of the Hida family over $L$. The twist by $\kappa$ of the family of representations of the deformation space has characteristic polynomials and triangular forms defined over its ring $B$, and it is unramified where the Hida setting requires. So it is an ordinary pair over $B$, and the universal property gives a homomorphism from the universal ordinary determinant ring to $B$.

% before: sha-siegel-automorphy-stabilizer
Hypothesis (5) of Theorem~\ref{sha-siegel-automorphy-adjoint} asks for irreducibility on the subgroup $\Gamma'$ on which the cyclotomic character is trivial. If an irreducible representation becomes reducible on a normal subgroup containing an element with distinct eigenvalues, the translates of a minimal invariant subspace are spanned by eigenvectors of that element. Counting eigenvectors then bounds the index of the stabilizer of such a subspace by the dimension.

% before: sha-siegel-automorphy-sign
A reducible restriction to $\Gamma'$ will produce a quadratic character, trivial on the kernel of $\varepsilon$, that twists the representation into itself. The residual conditions of a link exclude such characters only when they are trivial on the larger kernel of $\varepsilon$ modulo $\ell$. Since $\ell$ is odd, a quadratic character with open kernel is trivial on this larger kernel as well.

% before: sha-siegel-automorphy-sign-twist
A self-twist by a sign character is visible on characteristic polynomials, and these reduce to products of the residual characters. So the reduction of the sign character permutes the residual characters, which is the form in which the residual conditions of a link apply.

% before: sha-siegel-automorphy-monomial
The second half of hypothesis (5) asks for elements of $\Gamma'$ with separable characteristic polynomial in every open set, and the argument for it needs normal subgroups whose images do not commute. If the images of a normal subgroup commute and include an element with distinct eigenvalues, the whole group permutes its eigenlines, so the representation is monomial. Reducing the eigenvalues then shows that the residual characters agree on the kernel of the permutation action.

% before: sha-siegel-automorphy-noncommuting
In rank four the permutation action of a monomial representation takes values in the symmetric group on four letters, whose elements have order at most four. So a pair of residual characters whose ratio has order greater than four rules out commuting images on every normal subgroup of finite index.

% before: sha-siegel-automorphy-zariski
To find a separable characteristic polynomial in a coset of a subgroup of finite index, the argument works with the Zariski topology that the representation induces on the group. Its irreducible component through the identity is a normal subgroup of finite index. Since this component is irreducible, finitely many polynomial conditions that can each be met on a coset of it can be met together, and inside any subgroup of finite index.

% before: sha-siegel-automorphy-weight
The separability is first proved after replacing the powers of a diagonal element by powers of an indeterminate. This replacement needs integer weights on the eigenvalues that respect every multiplicative relation between them and still separate them.

% before: sha-siegel-automorphy-torus
The discriminant of a characteristic polynomial is a polynomial in the entries of a diagonal twist. If it vanished at all powers of the eigenvalues, the relations among them, read through the weights of Lemma~\ref{sha-siegel-automorphy-weight}, would make it vanish at the powers of an indeterminate. So separability over the field of rational functions gives separability at some power.

% before: sha-siegel-automorphy-quartic
Over the field of rational functions, separability is decided by the degrees of the coefficients of the characteristic polynomial. Two shapes of matrices occur in rank four: those whose diagonal entries and principal minors of size three are nonzero, and block antidiagonal ones made of two invertible blocks of size two.

% before: sha-siegel-automorphy-coset-four
In rank four, Lemma~\ref{sha-rigidity-coset} is replaced by an argument on the supports of the matrices in a coset of the identity component. A relation between indices records which entries are somewhere nonzero on such a coset, and the noncommuting elements force one of the two shapes of Lemma~\ref{sha-siegel-automorphy-quartic}. Lemma~\ref{sha-siegel-automorphy-torus} then passes from an indeterminate back to a power of a diagonal element.

% before: sha-siegel-automorphy-link-self-twist
If an irreducible representation becomes reducible on a normal subgroup that contains the commutators, it is isomorphic to its twist by a nontrivial character trivial on that subgroup. The argument needs only that the index of the subgroup be invertible in the field of coefficients: some element moves the projection onto an invariant summand, and a common eigenvector of the group acting on the endomorphisms fixed by the subgroup gives the twist.

% before: sha-siegel-automorphy-rank-four
These lemmas give hypothesis (5) of Theorem~\ref{sha-siegel-automorphy-adjoint} for a representation of rank four. The residual conditions they use are those that the link and the tower supply: no quadratic character that is trivial where the cyclotomic character is $1$ modulo $\ell$ permutes the residual characters, and two of them have a ratio of order greater than four.

% before: sha-siegel-automorphy-link-root
The character $\kappa$ of a link is built from characters with values in $1+\ell\mathbf Z_\ell$ and from their roots. Such a character has a logarithm with values in $\mathbf Z_\ell$, and its $n$-th roots close to $1$ are the powers of a single $n$-th root of $1+\ell$.

% before: sha-siegel-automorphy-link-half
In the relation (a) of a link, $\kappa(g)\kappa(\tau g\tau^{-1})$ must be $\theta(g)^{-1}\chi_{\mathrm{cyc}}(g)^{-3}$, a power of $\chi_{\mathrm{cyc}}\omega^{-1}$ times a character of finite order. A character of all of $G$ takes the same value at $g$ and at $\tau g\tau^{-1}$, so a square root of $\chi_{\mathrm{cyc}}\omega^{-1}$ on $G$ supplies that power.

% before: sha-siegel-automorphy-link-hecke
When $k$ is odd, the power of that square root is a half-integral power of the cyclotomic character on inertia above $\ell$, while (e) of a link asks for an integral one up to finite order. An $\ell$-adic character of the imaginary quadratic field, unramified outside one prime above $\ell$, corrects this. Class field theory gives its reductions modulo the powers of $\ell$, and the density of Frobenius elements glues them.

% before: sha-siegel-automorphy-link-beta
A root of a twist of that character is congruent to $1$, satisfies $\beta(s)\beta(csc^{-1})=1$, and on inertia above $\ell$ is the square root of an odd power of the cyclotomic character, up to finite order. Its product with the inverse of the square root of $\chi_{\mathrm{cyc}}\omega^{-1}$ is then an integral power of the cyclotomic character up to finite order on these inertia groups.

% before: sha-siegel-automorphy-link-dirichlet
What remains of $\theta(g)^{-1}\chi_{\mathrm{cyc}}(g)^{-3}$ is the character $\omega^{-(k+3)}\xi^{-1}\psi^{-1}$, of finite order. On the field of a deep tower only its part of order prime to $\ell$ matters, and that part is an even Dirichlet character, unramified outside $S$.

% before: sha-siegel-automorphy-link-kappa-zero
On $G_K$, an even Dirichlet character whose ramified primes split in $K$ has the form $s\mapsto\kappa_0(s)\kappa_0(csc^{-1})$ for a character $\kappa_0$ of finite order. Class field theory gives $\kappa_0$ from a character of the ideals of $K$ prime to the conductor, defined through a homomorphism from the ring of integers of $K$ to the integers modulo the conductor.

% before: sha-siegel-automorphy-link-characters
The character $\kappa$ of a link is the product, on $G_L$, of $\kappa_0$, a power of the inverse of the square root of $\chi_{\mathrm{cyc}}\omega^{-1}$ and, when $k$ is odd, the root $\beta$; the characters $\beta_L$ and $t_L$ are the restrictions of $\beta$ and of its logarithm. Their conditions follow from those of the pieces and, for the relation (a), from the depth of the tower, which makes the values of $\psi$ on $G_L$ of order prime to $\ell$.

% before: sha-siegel-automorphy-link-places
A link also needs the places of Situation~\ref{sha-hecke-situation} over $L$: a finite set of places split in $L$ and an auxiliary place $v_0$ at which the residual characters coincide and the cyclotomic character modulo $\ell$ is nontrivial. Condition (e) of a tower gives an element of $G_L$ with these two properties, and the Chebotarev density theorem gives a place $v_0$ whose decomposition group is generated by it modulo a small open subgroup.

% before: sha-siegel-automorphy-link-exists
With these pieces a link exists over every large enough ring of coefficients. The residual characters are the reductions of the twisted relabelling characters; they factor through the Galois group of the maximal extension unramified outside the chosen places, are distinct by (f) and are polarized by the relation (a).

% before: sha-siegel-points-three-types
In case III the promodularity of a flagged representation follows the same steps as in case II, with three changes. The first tower carries an imaginary quadratic field whose quadratic character is the residual character of weight $(1,0)$ on the absolute Galois group of the field of the tower, the seeds are of a new type A, and the pairs joining a pair over the lattice to a seed may pass through pairs of type A that are not ample.

% before: sha-siegel-points-three-tower-quadratic
The tensor field of the first tower is an imaginary quadratic field ramified at a single prime, chosen larger than a given bound and split at the primes of $S$. Its quadratic character is obtained from a Legendre symbol, and a Gauss sum identifies its field.

% before: sha-siegel-points-three-tower-elements
The element $g_0^T$ of the first tower is built from three elements of $G$, one of which moves $\psi$ while fixing $\xi$ and $\omega$. The prime of the tensor field is chosen large enough for $\xi$, $\psi$ and $\omega$ to be unramified there.

% before: sha-siegel-points-three-tower
The first tower in case III satisfies conditions (a) to (m) and has a tensor field. Being of tensor type also records the conditions $\ell\ge5$, $S_b=\emptyset$ and $\psi^m=1$ of case III, so the tower is fit and the lemmas on fit towers apply to it.

% before: sha-siegel-points-three-seeds
The seeds of case III are integral seeds of type A. They lie over a fixed residual pair that is defined at one auxiliary prime, and they exist over every larger tower that kills that prime.

% before: sha-siegel-points-three-gysin
The pair of the previous lemma must still be joined to a pair over the lattice. The Gysin construction of Lemma~\ref{sha-siegel-step-chains-gysin}, with the two weights exchanged, gives a pair with the same entries of weight $(-1,0)$ and with an entry of weight $(0,-1)$, through which the two can be linked.

% before: sha-siegel-points-three-chains
With these pairs, a pair over the lattice is joined to a seed of type A in four links, alternately keeping the entries $(3,1)$ and the entries $(2,1)$. The last pair before the seed is of type A, so the sequence is an A-chain.

% before: sha-siegel-points-three-descent
The last step of case II uses three formal properties of chains: they descend to a smaller ring of coefficients, they end at residual pairs, and they propagate the modularity ideal from one end to the other. Each holds for A-chains with the same proof.

% before: sha-siegel-points-three-bounds-rank
At a pair of type A the flags have a nonzero entry $(2,1)$, and the tangent pairs of the split pair with these flags are controlled once three of their entries vanish: one entry of an odd root and the two diagonal entries $(1,1)$ and $(2,2)$. The count is that of Lemma~\ref{sha-siegel-step-dimensions-tangent-split}, with two more spaces of cocycles of weight $0$.

% before: sha-siegel-points-three-bounds-cases
The split pair of type A satisfies the hypothesis of the previous lemma for five of the nine sets of weights. The flags relate the entries $(1,2)$ and $(2,1)$ of a tangent pair at the places of $\Sigma_T$, and make the entries $(3,2)$ and $(4,1)$ coboundaries there.

% before: sha-siegel-points-three-bounds-tangent
A residual pair of type A degenerates, along the grading $m=(0,0,1,1)$, to the split pair with flags $1+b_v(E_{21}-E_{43})$. On the decomposition groups of $\Sigma_T$ the first two diagonal characters of such a pair agree, which is the input of the previous lemma.

% before: sha-siegel-points-three-bounds-torus
On the locus of type $\{\pm(-1,1),\pm(1,1)\}$ the representation does not span $M_4$, so Lemma~\ref{sha-siegel-step-dimensions-torus}(1) does not apply to it. When all four entries are present, two of them, with the traces and a regular diagonal element, still generate the field of all the entries.

% before: sha-siegel-points-three-bounds
Propagating the modularity ideal at a point of type A needs bounds on the dimensions of the loci where the representation degenerates. The bound of the first part covers eight of the nine types of Lemma~\ref{sha-siegel-step-dimensions}; for the type $\{\pm(-1,1),\pm(1,1)\}$ the second part bounds the fibres over $B$ instead, and the third shows that on it the diagonal characters coincide at every place where the flag has a nonzero entry $(2,1)$.

% before: sha-siegel-points-three-propagation
At a point of type A the argument of Lemma~\ref{sha-siegel-step-propagation} goes through with the bounds of the previous lemma. The one new case is a curve on which the representation is supported in $\{\pm(-1,1),\pm(1,1)\}$, and the coincidence of diagonal characters excludes it.

% before: sha-siegel-points-three-last
The last step of case III is that of case II, with A-chains in place of chains. The descent of A-chains and the propagation at points of type A replace the corresponding lemmas of case II.

% before: sha-siegel-points
The bound for cases II and III needs four properties of the points of a Siegel datum: the functions of weight zero on them form a finite algebra, none of them is endoscopic, and the irreducible ones are automorphic and rigid. Finiteness and automorphy come from comparing the points with modularity ideals over towers of abelian fields; rigidity comes from the vanishing of an adjoint Bloch-Kato Selmer group.

% before: sha-eigensystem-level
The proof of case I compares Frobenius traces of Galois representations of bounded ramification with an Eisenstein value. Their ramification is controlled at every prime by characters that factor through the $M$-th cyclotomic field, and at the primes above $\ell$ by a triangular shape whose diagonal characters are powers of the cyclotomic character up to such characters.

% before: sha-eigensystem-triple
The representations compared with the Eisenstein value come in two kinds. The stable ones are three-dimensional and automorphic for $GL_3$ over $K$; the endoscopic ones place a two-dimensional representation of $\operatorname{Gal}(\bar{\mathbf Q}/\mathbf Q)$, attached to a newform, in the corners of a $3\times3$ matrix.

% before: sha-eigensystem-data
At a fixed prime, a set of automorphic representations of bounded level and weight gives a finite family of Frobenius eigenvalues or traces. The family is stable under the automorphisms of $\mathbf C$, its members become algebraic integers after multiplication by a fixed power of the norm, and their absolute values are bounded; these are the properties that the product of differences in Lemma~\ref{sha-eigensystem-norm} needs.

% before: sha-eigensystem-cuspidal-count
The size of these families is bounded by the number of cuspidal representations of given infinitesimal character and level. Such representations contribute to the cohomology of a locally symmetric space, so their number is at most a dimension of cohomology, which grows polynomially in the level and in the weight.

% before: sha-satake-slopes
The eigenvalue data below need the Satake parameters of the counted representations to become algebraic integers after multiplication by powers of the norm. Integrality is tested one prime $\ell$ at a time, through every isomorphism $\iota:\bar{\mathbf Q}_\ell\cong\mathbf C$, and at each $\ell$ a Satake parameter is bounded above and below by products over the embeddings that induce the place. The exponents in these products are read off from the infinitesimal character, and the bounds are required for every range of integers that contains them.

% before: sha-satake-integer-criterion
A complex number that is integral at every prime, through every isomorphism $\iota$, is an algebraic integer. Transcendental numbers fail the test, because some $\iota$ moves them out of the unit ball.

% before: sha-satake-place-embeddings
Through $\iota$, an embedding of a number field into $\mathbf C$ becomes $\ell$-adic, so it induces a finite place over $\ell$ and extends to the completion there. Over a CM field, complex conjugation exchanges the two places over a split place of the totally real subfield, which decides which embeddings induce a given split place.

% before: sha-satake-norm-product
The bounds of Definition~\ref{sha-satake-slopes} are products of absolute values of a uniformizer under the embeddings that induce $w$. Their product over all these embeddings is the absolute value of the norm of $w$, which turns a bound with constant exponents into a power of $q$.

% before: sha-satake-character-value
A twist by an algebraic Hecke character multiplies the Satake parameters by its value at a uniformizer and shifts the infinitesimal character by its infinity type. The two changes cancel in the slope bounds, by the triviality of the character on a generator of a power of $w$.

% before: sha-satake-newton-polygon
The Hecke operators at a split place bound the elementary symmetric functions of the Satake parameters. An ultrametric argument on the characteristic polynomial turns these into bounds for the parameters themselves.

% before: sha-satake-weight-matrices
On the definite unitary group, the action of the archimedean weights is transported to places over $\ell$, where it is given by an algebraic representation of $GL_n$ evaluated on $\ell$-adic matrices. The matrix entries of this representation on integral matrices are controlled by the determinant, through the Cartan decomposition and the weights of an integral basis.

% before: sha-satake-eigenvalue-bound
A Hecke eigenvalue on a vector-valued form is bounded by the growth of the words in the operator, because the form, moved back by the representation, is bounded on the finitely many double cosets. Only these properties are used, so the bound is stated for an abstract group.

% before: sha-satake-vector-forms
The bound of Lemma~\ref{sha-satake-eigenvalue-bound} needs forms with values in a representation of the finite adelic points. A copy of the tensor product of the weights at infinity inside the unramified vectors gives such forms, with the rational points acting through the weights.

% before: sha-satake-hecke-operators
At a place split over the totally real subfield, the unitary group is $GL_n$, and its spherical Hecke operators act on the unramified forms through the Satake parameters of the local component. Multiplied by $q^{j(j-1)/2}$, their eigenvalues become the elementary symmetric functions of the parameters scaled by $q^{(n-1)/2}$.

% before: sha-satake-operator-bound
The Hecke operators of Lemma~\ref{sha-satake-hecke-operators} consist of integral matrices whose determinants have a fixed valuation. Transporting the weights to the places over $\ell$, Lemma~\ref{sha-satake-weight-matrices} bounds the words in these matrices, and Lemma~\ref{sha-satake-eigenvalue-bound} bounds the eigenvalues.

% before: sha-satake-unitary-slopes
The eigenvalues of Lemma~\ref{sha-satake-hecke-operators}, their bounds in Lemma~\ref{sha-satake-operator-bound} and the Newton polygon give the slope bounds for automorphic representations of the definite unitary group. The exponents are the weights shifted as in the infinitesimal character.

% before: sha-satake-descent
Conjugate self-dual cuspidal representations reach the definite unitary group through Labesse's descent. It keeps the local components at the split places and matches the weights with the infinitesimal characters, which are the two properties that the slope bounds use.

% before: sha-satake-translates
The twist below must be unramified wherever the good scope gives Satake parameters, and its proof shows that a character is unramified where both $\pi$ and the twist of $\pi$ by that character are spherical. The local input is that a vector on which $GL_a(\mathcal O_w)$ acts through a ramified character of the determinant has all its translates in the span of the vectors $\rho(k)z-z$. The Cartan decomposition reduces this to diagonal elements, which commute with a diagonal unit of any given determinant.

% before: sha-satake-extension
The twist itself is a Hecke character of $E$ whose restriction to $F$ is prescribed by the polarization. It has to be algebraic, and unramified at the given split place and at every place that does not split, where the self-dual scope asks for Satake parameters. Alpöge builds it on a congruence subgroup of the ideles of $E$ whose only global elements with absolute value one at every embedding are trivial, which an auxiliary split place ensures.

% before: sha-satake-twist
Descent applies only in the self-dual scope, where the polarizing character is a power of the quadratic character. A twist by an algebraic Hecke character, unramified at the given split place, moves a pair from the good scope into the self-dual scope.

% before: sha-satake-twist-support
The ramification step of the twist argues with the vectors of $\pi$ and the irreducibility of $\pi$. The local theory of $GL_a$ gives a shorter route.

% before: sha-satake-quadratic-extension
A general place need not split over the totally real subfield, and the good scope need not hold. A quadratic base change to $E(\sqrt{-p})$, for suitable primes $p\equiv-1$ modulo a fixed integer, makes the given place split, keeps its Satake parameters, and puts the base change in good scope.

% before: sha-satake-unramified-places
The quadratic base change in the last step of the next proof comes from the weak form of cyclic base change, which controls the Satake parameters only at almost all places. The other places are reached through the local properties of the base change of Arthur and Clozel: it takes unramified representations to unramified ones at unramified places, and principal series to principal series at a place with only one place above it.

% before: sha-satake-unramified-principal
At a ramified place the Satake parameters of the base change are read off from principal series. A spherical representation is a subquotient of an unramified principal series whose characters give its Satake parameters, and an unramified principal series with no two characters differing by the absolute value is irreducible and spherical.

% before: sha-satake-ramified-places
At a place ramified in a cyclic extension of prime degree the residue field does not grow, so the norm takes unramified characters to unramified characters with the same value at a uniformizer. The Jacquet-Shalika bound keeps two Satake parameters of a cuspidal representation from differing by the factor $q$, so the principal series of the base change is irreducible and spherical, with the same Satake parameters.

% before: sha-satake-integrality
The integrality of the Satake parameters reduces to the definite unitary group through a quadratic base change, a twist and descent. The slope bounds obtained there, for all primes and all $\iota$, give the two integrality conditions of Definition~\ref{sha-eigensystem-data}.

% before: sha-satake-integrality-base-change
The quadratic base change in the last step goes through the weak form of cyclic base change. The strong form gives it at once.

% before: sha-eigensystem-distinctness
Local-global compatibility at a prime above $\ell$ uses that the diagonal characters of the representation there differ from their cyclotomic twists. If no two weights differ by one, this follows from the weights; otherwise the counts below assume it at $\ell$, together with $\ell\mid M$. Over an imaginary quadratic field the assumption has to survive restriction to subgroups of finite index, so it is also needed for all powers of the characters.

% before: sha-eigensystem-unitary
Over a CM field, the cuspidal representations counted above give an eigenvalue datum through their Satake parameters at a prime away from the level. Clozel's theorem makes the family stable under $\operatorname{Aut}(\mathbf C)$, the Jacquet-Shalika bound gives the window, and local-global compatibility shows that a representation of the given level and weights is determined by the Satake parameters of the cuspidal representation it is attached to.

% before: sha-eigensystem-twists
The count over $\mathbf Q$ restricts a representation to an imaginary quadratic field on which it stays irreducible. Two facts about twists make this possible: an irreducible representation of dimension $n$ is isomorphic to its twist by at most $n^2$ characters, and its restriction to a subgroup of index two is irreducible unless it is isomorphic to its twist by the quadratic character. An irreducible restriction then determines the representation up to that twist.

% before: sha-eigensystem-split-field
To carry $\ell$-distinctness from $\mathbf Q$ to an imaginary quadratic field unchanged, the field should not shrink the decomposition groups at $\ell$, that is, $\ell$ should split in it. The fields $\mathbf Q(\sqrt{-P})$ with $P\equiv-1$ modulo $4\ell$ have this property, and among $n^2+1$ of them one keeps the representation irreducible.

% before: sha-eigensystem-flags
Potential $\ell$-distinctness behaves well under a change of field. It passes to a subgroup of finite index, since a fixed power of every element lies in the subgroup, and it descends from such a subgroup to the whole decomposition group, since the trace of a triangular form determines its diagonal characters.

% before: sha-eigensystem-potential
For a representation of $\operatorname{Gal}(\bar{\mathbf Q}/\mathbf Q)$ that is automorphic and windowed over imaginary quadratic fields, $\ell$-distinctness implies potential $\ell$-distinctness. The windowed condition gives it over a field in which $\ell$ splits, and it descends to $\mathbf Q$ and then passes to every number field. The count over $\mathbf Q$ needs this to apply the count over an imaginary quadratic field.

% before: sha-eigensystem-rational
The endoscopic triples involve representations of $\operatorname{Gal}(\bar{\mathbf Q}/\mathbf Q)$, and so do the representations of Siegel type below; the count over a CM field does not apply to them directly. Restricting to one of several imaginary quadratic fields, chosen so that the restriction stays irreducible, reduces the count over $\mathbf Q$ to counts over these fields, at the cost of taking square roots of Frobenius eigenvalues at the inert primes.

% before: sha-eigensystem-traces
The counting proposition only uses traces of Frobenius. The eigenvalue data of Lemmas~\ref{sha-eigensystem-unitary} and~\ref{sha-eigensystem-rational} give trace data of the same size for the stable and the endoscopic triples, and their window conditions bound the traces by $\sqrt q+2$.

% before: sha-eigensystem-eisenstein
The Eisenstein value has absolute value about $q^k$, which exceeds the bound $\sqrt q+2$ on the traces of a trace datum. So a trace differs from an Eisenstein value, and the difference times $q^k$ is an algebraic integer of controlled size.

% before: sha-eigensystem-norm
Multiplying these differences over a trace datum and over a Galois-stable family of roots of unity gives a nonzero integer. Its archimedean size is controlled by the size of the datum, and its $\ell$-adic size by the differences that come from the representations being counted.

% before: sha-algebra-graded
The one-character proposition bounds a Selmer group by the length of a cotangent module of a graded model of a deformation ring: the ring universal for the lifts of $r_0(\eta,\Psi)$ whose restriction at $\ell$ is cut out by an ordinary ideal. The ring is graded by the action of the diagonal matrices $\operatorname{diag}(t,1,t^{-1})$, and the bound comes from its points at which functions of positive and of negative weight do not all vanish.

% before: sha-cotangent-finite
The proof of the next commutative algebra statement needs a finite set of two-signed points that meets every orbit. Such points are counted by slices on which one function of weight $-2$ equals $1$ and one of weight $2$ is invertible. Each slice is a finitely generated algebra over the fraction field with no tangent vectors at its points, so it has finitely many points.

% before: sha-cotangent-implicit
The chart of a graded algebra may have variables of weight $0$, while the cotangent bound needs a chart without them. The functions of weight $0$ that span the cotangent lattice of weight $0$ can be solved for these variables, by the implicit function theorem for formal power series in a form that respects weights.

% before: sha-cotangent-euler
The cotangent bound below is proved by comparing two Euler characteristics on a Rees algebra graded by $\mathbf Z^3$. They are defined one element at a time, by the quotient and the kernel of the multiplication, as for the multiplicity symbol of a local ring, and are only used in large degrees.

% before: sha-cotangent-calculus
The Euler characteristics obey the rules of the multiplicity symbol: additivity on exact sequences, independence of the order, positivity, the bound by the quotient, and invariance under changes that kill the module. Here they hold in large degrees, once the list of elements is admitted.

% before: sha-cotangent-nullstellensatz
The proof of the cotangent bound produces points of the graded algebra from primes of its Rees algebra. Each comes from a ring homomorphism on the part of degree $0$, extended by the Nullstellensatz so that it kills a given ideal and not a given element.

% before: sha-cotangent-herbrand
The upper bound for the Euler characteristic of a degree-zero element reduces to modules of finite length over the localizations at the points. On such a module the length of the cokernel minus that of the kernel of multiplication by the element is a sum over the points of the local lengths times the values of the element.

% before: sha-algebra-cotangent
The first commutative algebra statement compares the cotangent module of weight $-1$ of a graded algebra with the values of a function of weight $0$ at the two-signed points. Its hypotheses say that the algebra is smooth of the expected dimensions at the base point, finite in weight $0$, and rigid at its other points.

% before: sha-algebra-derivations
The second commutative algebra statement bounds the length of a module by the same cotangent length, when the elements of the module give derivations of the right weight modulo a constant. Such a module embeds in the dual of the cotangent module, which has the same length.

% before: sha-kato-selmer
At the endoscopic points, the vanishing of the adjoint Selmer group uses the vanishing of the Selmer groups of the representation attached to the newform and of its twist by the quadratic character of $K$. These are Bloch-Kato Selmer groups of the Galois representation of a modular form, which Kato's Euler system controls.

% before: sha-rigidity-normalize
An ordinary point is normalized relative to one element $g_1$, while a link comes with its own base element. To pass from one to the other, the point is conjugated by a matrix congruent to $1$, found by solving quadratic equations with Hensel's lemma. Two further quadratic equations first make the matrix part at the complex conjugation equal to $1$.

% before: sha-rigidity-free
The adjoint lifting theorem needs, at each place above $\ell$, an element of the decomposition group with trivial cyclotomic character whose eigenvalues have no ratio that is a root of unity. Such an element comes from local-global compatibility at $\ell$: the windowed condition keeps the ratios of the Frobenius eigenvalues of the local Weil-Deligne representation below $q$ in absolute value, while distinct Hodge-Tate weights would force ratios that are powers of $q$.

% before: sha-rigidity-adjoint
The vanishing of adjoint Selmer groups at the stable points here, and at the lifts of Definition~\ref{sha-lifts-problem}, comes from the adjoint lifting theorem over the abelian CM field of a link. It uses only four properties of an ordinary point: the parity of $\Psi$ at $c$ with the ramification of $\eta$, a link with a seed datum, irreducibility with automorphy, and triangular forms at $\ell$ whose characters are not related by the cyclotomic character.

% before: sha-rigidity-two-by-two
At a block shaped point the corner block is a representation $W$ of dimension two, and its level and irreducibility are read off from the triangular and diagonal forms of the point on inertia. These forms are given on the kernel of the quadratic character, so passing to $W$ and to the whole inertia group needs a few facts about $2\times2$ matrices.

% before: sha-rigidity-corner
At a point whose matrix parts on the kernel of $\eta$ have the block shape of Lemma~\ref{sha-rigidity-points}(2), the corner block inherits triangular forms from those of the point at the places above $\ell$. The first and last vectors of a flag of the point lie in the corner as soon as the outer diagonal characters are nontrivial somewhere, and they give a flag of the corner block.

% before: sha-rigidity-lagrange
The ordinary relations of a linked point say that a product of the matrices $r(g_i)-\psi_i(g_i)$ vanishes. When one element $g_0$ separates the characters modulo the maximal ideal, the Lagrange idempotents of $r(g_0)$ split the module into lines, and the relations make $r$ triangular in a basis of such lines, over the coefficient ring itself.

% before: sha-rigidity-clifford
Hypothesis (5) of Theorem~\ref{sha-siegel-automorphy-adjoint} asks for irreducibility on the kernel of the cyclotomic character, a normal subgroup. In prime dimension, Clifford's theorem leaves three possibilities for the restriction of an irreducible representation to a normal subgroup.

% before: sha-rigidity-characters
The irreducibility on a normal subgroup is proved by computing the eigenvalues of the same matrices in two ways. Two families of characters with the same characteristic polynomials everywhere agree up to a permutation.

% before: sha-rigidity-normal
The monomial case of Clifford's theorem is excluded by the residual characters, since the permutation of the lines would identify two residual characters that differ. The scalar case is excluded by an element of the subgroup that separates the characters at a place above $\ell$.

% before: sha-rigidity-small
For an open subgroup, the intersection of its conjugates with the kernel of the cyclotomic character is normal, and a power of a free element lies in it. So irreducibility on these subgroups follows from a free element and from two residual characters that differ.

% before: sha-rigidity-coset
The second half of hypothesis (5) asks for elements with separable characteristic polynomial in every open set that meets the kernel of the cyclotomic character. In dimension at most three a translate $y\,r(h)$ can be made separable, using powers of an element with free eigenvalue ratios and an element whose matrix and inverse have no zero entries.

% before: sha-rigidity-separable
A neighbourhood of an element $s$ contains a coset $sV$ of an open subgroup, and the representation stays irreducible on the subgroups of finite index of $V\cap G'$. So Lemma~\ref{sha-rigidity-coset}, applied with $y=r(s)$, gives an element with separable characteristic polynomial in the neighbourhood.

% before: sha-rigidity-windowed-flags
Passing the windowed condition up a field extension compares triangular forms at $\ell$ over the two fields. A triangular form restricts to the decomposition groups of the larger field, and $\ell$-distinctness over the larger field descends to the smaller one where the representation is triangular. Lemma~\ref{sha-eigensystem-flags} gives the corresponding two facts for potential $\ell$-distinctness.

% before: sha-rigidity-windowed-local
Condition ($\beta$) of the windowed condition bounds the absolute values of the characters in the supercuspidal supports of the local components. Under a cyclic base change of prime degree these characters are composed with the local norm, and the order of the residue field is raised to the power of the residue degree. Both sides of the bound are raised to the same power, so the bound passes to the base change with the same exponent $t$.

% before: sha-rigidity-windowed-up
In the proof of Lemma~\ref{sha-rigidity-points}(2), the windowed condition for a two-dimensional representation over an imaginary quadratic field must pass to an abelian CM field containing it. The passage works for any solvable Galois extension of CM fields, by induction along a tower of extensions of prime degree, with the two preceding lemmas at each step. The polarization restricts to the intermediate fields, and solvable base change gives automorphy over them.

% before: sha-rigidity-points
Condition (4) of Lemma~\ref{sha-algebra-cotangent} asks that the deformation ring be rigid at its two-signed points. Rigidity at a point follows from the vanishing of the Bloch-Kato Selmer group of the adjoint representation there, which the next statement gives for the two kinds of points.

% before: sha-admissibility-stable
The stable points produced by the automorphy theorem must be automorphic triples of the level that the counting proposition counts. Their ramification away from $\ell$ and their triangular, potentially crystalline shape at $\ell$ come from the conditions on the $\mathfrak a_B$-lifts, and the condition on their characters at $\ell$ from the linked points of a link.

% before: sha-siegel-count
For cases II and III, each congruence between an irreducible point of a Siegel datum and the sum of four characters is measured by the difference of their traces at a Frobenius element. The product of these differences over all eigensystems of bounded level and weight is a nonzero integer whose logarithm is polynomial in the level and the weight, and it bounds the congruences at every prime $\ell$ at once.

% before: sha-lifts-problem
In case V the prime $\ell$ is a bad prime of $\psi$ and the twist is $1$, so the Eisenstein data used in case I are not available. The bound comes instead from deformations of $r_0(\eta,\Psi)$ over $\mathbf Q$, with $\eta$ the character of $\mathbf Q(i)$, that are ordinary at $\ell$ and diagonal at the primes where the $\ell$-power part of $\psi$ ramifies. The dyadic and Bernoulli conditions are what the argument uses of $\psi$; a large $2$-conductor separates the residual characters at these primes.

% before: sha-lifts-graded
Conjugation by the diagonal matrices $\tau(t)$ fixes $r_0$ and acts on unitary lifts, so it grades the ring that classifies them. The points of the graded subring are the lifts in characteristic zero up to this action, and a point is rigid when its first-order deformations with the same local conditions are trivial.

% before: sha-lifts-selmer
The size of the universal ring is governed by the residual cohomology of the graded pieces of the adjoint representation. The Selmer groups give its tangent space in each weight, the obstruction groups bound its relations, and the groups $\text{Ш}_a$ are the part of weight $-1$ that carries the Tate-Shafarevich group being bounded.

% before: sha-lifts-host-data
The hypotheses of case V with small $2$-conductor give the dyadic and Bernoulli conditions. The primes where the $\ell$-power part of $\psi'$ ramifies divide $N_\psi$, so the inequality on the $2$-conductor bounds them as well as $\ell$.

% before: sha-lifts-local
The dyadic condition provides, at each prime of $T_b$, an element of the decomposition group on which $\Psi^2$ is not congruent to $1$. These elements give the local vanishing statements and make the triangular conditions closed, while the other primes enter only through characters of order prime to $\ell$.

% before: sha-lifts-universal
Mazur's representability theorem applies to $\Gamma$, and the local conditions and the normalization cut out a quotient of the universal ring of all lifts. Conjugation by $\tau(t)$ preserves the conditions, which gives the grading of the quotient.

% before: sha-lifts-residual
The count of points needs five facts about the residual cohomology: local vanishing at the primes of $T_b$, an Euler characteristic count in each weight, the class number one of $\mathbf Q(i)$, torsion of the integral classes of nonpositive weight, and the vanishing of the Selmer group of weight $-2$. The last two come from Theorem~\ref{sha-soule-bound} and from the residual nonvanishing of Bernoulli numbers.

% before: sha-lifts-presentation
The universal ring is presented as a quotient of a power series ring graded by the weights of the matrix entries. The Euler characteristic count compares the numbers of variables and of relations in each weight, and the Selmer classes and the groups $\text{Ш}_a$ give tangent vectors at the base point.

% before: sha-lifts-points
The points of the graded ring are sorted by the weights on which they vanish. Points that vanish on neither the positive nor the negative weights are irreducible, and at a rigid point the only tangent direction is that of the orbit under $\tau$.

% before: sha-lifts-coordinates
The two commutative algebra arguments of case I apply to a graded power series ring whose relations are counted by weight. In the coordinates of the presentation they bound every module that injects into the tangent vectors of weight $-1$ by twice the valuations of one element at finitely many two-signed points.

% before: sha-lifts-abstract-count
The presentation, the description of the points and the coordinate count together bound the groups $\text{Ш}_a$ by valuations of differences of traces at irreducible points. The bound needs only a property of points that gives rigidity and the nonvanishing of these differences, and the automorphy, rigidity and counting lemmas below provide one.

% before: sha-lifts-comparison
The group defining $V'_\ell(e,\psi)$ is the union of its torsion submodules, and each of them is a group $\text{Ш}_a$ up to the local conditions at the primes of $T_b$. The elements $x_p$ of Lemma~\ref{sha-lifts-local}(b) make these local conditions agree.

% before: sha-lifts-automorphy
The irreducible points must be automorphic for the count to apply. The argument of Theorem~\ref{sha-residually-reducible-automorphy} gives this for the characters of a host datum, with $K=\mathbf Q(i)$, and its finiteness statement makes the part of weight $0$ of the universal ring finite.

% before: sha-lifts-rigidity
The count needs the irreducible points to be rigid. For automorphic points rigidity follows from the vanishing of the Bloch-Kato Selmer group of the adjoint representation, as for the points of case I.

% before: sha-lifts-number
The differences of traces at a Frobenius element are bounded by the counting proposition of case I, applied over $\mathbf Q(i)$. A prime inert in $\mathbf Q(i)$ and small compared with $\ell N$ makes the bound polynomial in $\ell N$.

% before: sha-lifts-automorphy-conditions
The descent step in the proof of Theorem~\ref{sha-residually-reducible-automorphy} builds its abelian CM field, its link and its seed datum from a short list of properties of the characters. For the characters of Definition~\ref{sha-lifts-problem} these properties follow from the dyadic condition, with $\mathbf Q(i)$ as the quadratic field.

% before: sha-lifts-automorphy-form
With these properties the rest of the proof of Theorem~\ref{sha-residually-reducible-automorphy} goes through for the characters of Definition~\ref{sha-lifts-problem}. It gives the finiteness of characteristic polynomials in characteristic $\ell$ and the automorphy over $\mathbf Q(i)$ of the ordinary points that are irreducible on every open subgroup.

% before: sha-lifts-triangularity
In characteristic $\ell$ the finiteness statement fixes the characteristic polynomials only on a subgroup of finite index. Together with the normalization at $g_1$, this forces a lift to be triangular on $\ker\eta$, with the residual diagonal.

% before: sha-lifts-finiteness-diagonal
The diagonal of a triangular lift is a deformation of the residual characters. The local conditions and the absence of unramified abelian extensions of $\mathbf Q(i)$ leave it no room to move.

% before: sha-lifts-finiteness-criterion
The part of weight $0$ of the universal ring is controlled by its points in characteristic $\ell$. If every such point kills the generators of weight $0$ and those of one sign, that part is finite over $\mathcal O$.

% before: sha-lifts-level
The counting proposition of case I applies to representations of $\operatorname{Gal}(\bar{\mathbf Q}/\mathbf Q(i))$ with a given level and given local behaviour at $\ell$. For the irreducible points both come from the local conditions of a unitary lift, and the condition at $\ell$ also enters the proof of rigidity.

% before: sha-model-local-component
The local properties of the ordinary ideal in the next lemma all come from a few facts about the ordinary component of the universal local deformation ring at the base point. They concern its completion, its points with values in finite $\mathbf Q_\ell$-algebras, its stability under conjugation, and its tangent directions of each weight.

% before: sha-model-local
In case I the bound comes from a ring that is universal for certain lifts of $r_0(\eta,\Psi)$ over $\mathbf Q$, and the condition at $\ell$ is imposed through the universal ring of local lifts at a place above $\ell$. The ordinary ideal of that local ring cuts out the components of the ordinary locus through $r_0$; the global ring is universal for the lifts whose restriction at $\ell$ kills it, not for the ordinary points themselves. Its quotient is smooth at the base point after inverting $\ell$, with twelve coordinates whose weights under the torus are known, and its tangent vectors there satisfy the Bloch-Kato condition.

% before: sha-model-universal
The global ring is cut out of a universal ring of lifts by the local condition at $\ell$, the triviality on the inertia elements away from $\ell$ that are residually trivial, and a normalization at $c$ and $g_1$. Its points with values in complete local domains are ordinary points, so the automorphy theorem applies to them.

% before: sha-model-points
The commutative algebra of the bound needs the finiteness of the part of weight $0$ of the global ring and a description of its points by the shape of their matrices on $\ker\eta$. The finiteness comes from the characteristic polynomials in characteristic $\ell$. A point that is both upper and lower triangular is the base point, and a two-signed point is either irreducible or endoscopic.

% before: sha-model-tangent
The tangent cocycle of a first-order deformation lies in the Bloch-Kato Selmer group of the adjoint representation. At the base point, in weights at most $0$, this group vanishes by class field theory and the bound of Soulé type; at a two-signed point it vanishes by the rigidity of automorphic points and Kato's theorem.

% before: sha-model-exists
The model replaces the global ring by a graded algebra to which the cotangent bound applies. It is built from a presentation of the global ring over the local ring at $\ell$, whose numbers of variables and relations in each weight come from Galois cohomology, followed by the eight normalization equations; these numbers are the ones the cotangent bound requires.

% before: sha-counts-siegel-dyadic
The Siegel datum used in cases II and III has a character $\xi$ of large $2$-power conductor, and its conditions at the Frobenius elements $\sigma_p$ come from the size of that conductor. A primitive character of conductor $2^n$ cannot take a value of small order at an odd integer $v$ unless $2^n$ is small compared with $v^2$, because it is nontrivial at $1+2^{n-1}$.

% before: sha-counts-siegel-auxiliary
Besides its characters, a Siegel datum consists of a complex conjugation, places with Frobenius elements, and the elements $\phi$, $g_0$ and $s_0$. These can be chosen for any characters $\psi$ and $\xi$; only the conditions at the $\sigma_p$ and the condition on residual cocycles depend on how $\xi$ is chosen.

% before: sha-counts-siegel-host
The Siegel datum attached to a character of a host datum is built from the parts of that character. Since the character is trivial on $H_h$, the host datum controls the orders and conductors of these parts and the primes at which the part of $\ell$-power order ramifies.

% before: sha-counts-siegel-stickelberger
The condition on residual cocycles in a Siegel datum asks that certain classes, unramified up to a subgroup of index prime to $\ell$, be coboundaries. For the residue of a Dirichlet character this follows from Stickelberger's theorem once a Bernoulli number of the inverse character is an $\ell$-adic unit.

% before: sha-counts-siegel-pair
In a Siegel datum the residual characters $\bar\alpha$, $\bar\beta$ and $\bar\theta$ are residues of products of a character of odd conductor and order prime to $\ell$ with a primitive character of large $2$-power conductor. For such products the Bernoulli number can be controlled modulo $\ell$.

% before: sha-counts-siegel-bernoulli
For the product of a Bernoulli pair the Bernoulli number is an $\ell$-adic unit. The large $2$-power conductor allows the sum defining it to be regrouped by residues modulo $2^n$ and computed modulo $\ell$.

% before: sha-counts-siegel-residual
The two preceding lemmas give the vanishing of the unramified residual classes for the product of a Bernoulli pair and for its inverse. The remaining hypothesis of the Stickelberger lemma comes from the $2$-power part, which takes the value $-1$ at an element fixing the $\ell$-th roots of unity.

% before: sha-counts-siegel-cocycles
In a Siegel datum the residual characters are residues of products of five Dirichlet characters, and the cocycles in its condition vanish on inertia only on the kernel of $\bar r_0$, or are known to be coboundaries locally at the primes of $T$. Both situations reduce to a Bernoulli pair and to classes unramified up to a subgroup of index prime to $\ell$.

% before: sha-counts-siegel-datum
Part (b) of Proposition~\ref{sha-counts} starts from a Siegel datum attached to the character. It exists once $\mathcal O$ contains enough $2$-power roots of unity: $\xi$ is chosen with $2$-power conductor so large that the conditions at the $\sigma_p$ hold and the residual characters come from Bernoulli pairs.

% before: sha-counts-vd-model
Part (b) of Proposition~\ref{sha-counts} bounds $V_D$ by the commutative algebra of Lemmas~\ref{sha-algebra-cotangent} and~\ref{sha-algebra-derivations}. That algebra needs a $\mathbf Z$-graded ring with a chart at a smooth point, while the entries of a Siegel point carry two weights, those of the torus elements $\tau(t)$. A graded model presents the ring of Siegel points as a quotient of a ring graded by $\mathbf Z^2$ with a chart, and the $\mathbf Z$-graded ring is its part of first weight $0$.

% before: sha-counts-vd-presentations
The rings of a graded model are complete local rings presented by power series in weighted variables, and the arguments below pass between such presentations. They use that quotients by homogeneous elements and changes of variables keep the parts of each weight, and that the part of weight $0$ is again a complete local ring, controlled by finitely many monomials.

% before: sha-counts-vd-extension
A homomorphism defined on the finite sums of homogeneous elements extends to the whole complete ring once the images of the variables are small. A torus element makes them small, provided the monomials of weight $0$ have small images; points of the $\mathbf Z$-graded ring become Siegel points in this way.

% before: sha-counts-vd-structure
The $\mathbf Z$-graded ring $A$ of a graded model satisfies the hypotheses of Lemma~\ref{sha-algebra-cotangent} on the ring and its generators. It is related to the ring $S$ of the Siegel point $r^{\mathrm u}$ through the kernel $I$, and an element of weight $0$ records the trace of a Frobenius element minus that of $r_0$.

% before: sha-counts-vd-span
When the Siegel datum is finite, every part of $S$ is spanned by products of entries of $r^{\mathrm u}$. Homomorphisms of the graded ring are then determined by their values on entries, and the part of weight $0$ is finite over $\mathcal O$, which is condition (2) of Lemma~\ref{sha-algebra-cotangent}.

% before: sha-counts-vd-points
The homomorphisms from $A$ to $\bar{\mathbf Q}_\ell$ that vanish on $I$ are the points of Lemma~\ref{sha-algebra-cotangent}. Each extends to the graded part of $S$, the torus acts on these extensions, and a torus conjugate of the resulting representation is a $\bar{\mathbf Q}_\ell$-point of the Siegel datum. Tangent vectors with values in the dual numbers behave in the same way.

% before: sha-counts-vd-conjugate
Passing from points of $A$ to Siegel points loses only the action of the torus, and the count of the bound needs points that are not conjugate. The normalization of Siegel points at $\phi$ and at $g_0$ forces a matrix conjugating one irreducible point into another to be a torus element, and forces a tangent vector at a rigid point to come from the torus.

% before: sha-counts-vd-flag
An ordinary lift at $\ell$ is defined by conditions on the idempotents $e_i$ built from $\rho(s_0)$. When the eigenvalues $\delta_i(s_0)$ are distinct, these conditions make the representation of $D_\ell$ upper triangular in a suitable basis, with diagonal characters equal to the $\delta_i$ on inertia. The shapes of the points of a Siegel datum use the ordinary condition in this form.

% before: sha-counts-vd-lattice
An off-diagonal entry of a point gives a residual cocycle only after it is divided by a common factor of its values. Dividing by a value of least valuation keeps one value a unit, so the residual cocycle is not zero. When the idempotent $E_{ss}$ is an integral combination of the matrices of the point, the residual cocycle is not even a multiple of the difference of two diagonal entries.

% before: sha-counts-vd-linked
The shapes of the points rest on a property of the relation $\succeq$ of a point: two linked indices in different blocks are linked to a third index. Otherwise the entry between them would give a residual cocycle that is a coboundary at every place, by the ordinary condition at $\ell$ and the vanishing at the bad primes, hence a global coboundary by the definition of a Siegel datum. Lemma~\ref{sha-counts-vd-lattice} excludes such a cocycle.

% before: sha-counts-vd-soule
The finiteness of the groups $V_\ell(k,\psi)$ reduces, by restriction to an abelian field containing the $\ell$-th roots of unity, to the bound of Soulé type of Theorem~\ref{sha-soule-bound}. The finiteness in turn shows that a cocycle with values in $\bar{\mathbf Q}_\ell$ for a negative Tate twist, which is a coboundary on the inertia groups and at $\ell$, is a global coboundary. This vanishing excludes tangent vectors of weight $(0,-1)$ at $r_0$.

% before: sha-counts-vd-tangent
To first order in $\varepsilon$, the ordinary condition at $\ell$ on a tangent vector at $r_0$ is a condition on the idempotents $e_i$. In weight $(0,-1)$ it makes the entry $(3,1)$ a coboundary on $D_\ell$ up to a nonzero constant, and in weight $0$ it makes the tangent vector vanish on inertia. These local facts start the global argument for the vanishing of such tangent vectors.

% before: sha-counts-vd-diag-henselian
The characters on the diagonal of a point of diagonal type are compared with those of the ordinary flag at $\ell$ through their residues. That comparison takes place over the complete local ring of a local lift, and integrality over a henselian ring gives the bound on the residues.

% before: sha-counts-vd-diag-local
A $\bar{\mathbf Q}_\ell$-point is ordinary at $\ell$, so it has a flag on $D_\ell$ whose characters are known on inertia. Their residues are those of the characters of $r_0$, because the point comes from a lift over a complete local domain.

% before: sha-counts-vd-diag-characters
For a point of diagonal type the diagonal entries are characters of $G$, and they must be the flag characters at $\ell$ in some order. The order is constrained by the residues at $\sigma_\ell$ and by the symplectic form, and the characters are Tate twists of characters of finite order.

% before: sha-counts-vd-diag-order
The order of the characters decides which entries of a point of diagonal type can be nonzero. A nonzero entry against the order of the flag would give a cohomology class that is locally trivial everywhere, and the bound of Soulé type makes it a coboundary.

% before: sha-counts-vd-chart-generators
The chart of $A$ is obtained from the chart of the graded part of $R$, whose generators of each weight are needed explicitly. Each variable is the linear term of a homogeneous element, and these elements span the cotangent spaces up to bounded denominators.

% before: sha-counts-vd-chart-kernel
In the chart of $A$ the cotangent space at the origin must be spanned by the new generators. Since the variables of $R$ of nonzero first weight enter $A$ only through products, it suffices to control a few monomials of low degree.

% before: sha-counts-vd-inputs
The conditions of Lemma~\ref{sha-algebra-cotangent} also need the shapes of the $\bar{\mathbf Q}_\ell$-points that are not irreducible, the vanishing of tangent vectors at $r_0$ of weights $(0,-1)$ and $(0,0)$, the block-diagonal limit of a torus orbit, the finiteness of the groups $V_\ell(k,\psi)$, and a chart of $A$. The shapes rest on the published local theorems at $\ell$ and on a bound of Soulé type, which also gives the vanishing of these tangent vectors and the finiteness.

% before: sha-counts-vd-irreducible
A two-signed point of $A$ must give an irreducible representation, since the bound is turned into a count of automorphic representations. The other shapes are excluded: a representation of diagonal type kills all positive or all negative degrees of $A$, and a block-diagonal limit would be an endoscopic point.

% before: sha-counts-vd-conditions
With these facts the graded model satisfies the conditions of Lemma~\ref{sha-algebra-cotangent} except the last, which depends on a Frobenius element. Rigidity enters condition (4): a tangent vector at a two-signed point comes from the torus, so it is a multiple of the Euler derivation.

% before: sha-counts-vd-sha
The classes counted by $V_D$ give derivations of the graded model of weight $-1$ at $x$, through the last condition in the definition of a graded model. Lemma~\ref{sha-algebra-derivations} then bounds $V_D$ by the cotangent length of the family $F$.

% before: sha-counts-vd-model-framing
The graded model is cut out of a ring that classifies global lifts together with framings at the primes of $T$, over a ring that classifies the local lifts at these primes. The framings $\alpha_v$ compare the global lift on $D_v$ with the universal local lift; they are normalized away at the end.

% before: sha-counts-vd-model-six
The framings and the block shape at $\phi$ are imposed by equations on six off-diagonal entries. For a symplectic matrix close to the identity these six entries control the others, and a few more equations make it the identity.

% before: sha-counts-vd-model-chart
The chart of the local base is given on its graded part, while the cover and the quotient in the construction need a chart on the whole ring. An element of the ring is approximated by elements of the graded part modulo powers of the kernel of the origin, and the chart extends by continuity.

% before: sha-counts-vd-model-tangent
The equations on the six entries can be imposed one at a time only if they are independent to first order. Conjugating the residual point by an element of $\mathfrak{sp}_4$ over the dual numbers gives the derivations that separate them.

% before: sha-counts-vd-model-cocycle
Condition (iv) of a graded model asks for a derivation of weight $(0,-1)$ attached to each locally trivial cocycle. A cocycle twists $r_0$ into a framing over a ring with square-zero ideal, and the universal property of the framed ring turns it into a derivation.

% before: sha-counts-vd-model-charted
The construction of a graded model passes through rings that carry charts but are not graded algebras, such as the local base and the cover of Step 2. A chart needs only a family of submodules, the parts, and a point $x$. Its variables measure the cotangent space at $x$ weight by weight, up to bounded denominators.

% before: sha-counts-vd-model-solve
Dividing a ring with a chart by equations whose linear terms are independent should remove one variable of the chart for each equation. Over a field this is the implicit function theorem for power series: the equations can be solved for some of the variables. The variables solved for are chosen with the weights of the equations, so that the solution respects the weights.

% before: sha-counts-vd-model-cut
Step 4 divides the cover by the equations of Step 3, and the quotient must keep a chart with the right numbers of variables. A chart passes to the quotient by elements whose images in the cotangent space are independent, and it loses one variable of the weight of each element.

% before: sha-counts-vd-model-kernel
The quotient of a ring by elements with independent linear terms need not be a domain, while a graded model is a domain. Dividing further by the kernel of the chart gives a domain with the same chart, and the weights and the derivations at the origin pass to it.

% before: sha-counts-vd-model-graded
Condition (i) of a graded model asks for a chart of the graded part, while the cut and the kernel produce charts of the whole ring. Every element of the ring is congruent to an element of the graded part modulo the square of the kernel of the origin, so the restriction of a chart is again a chart.

% before: sha-counts-vd-model-quotient
Step 4 of the construction imposes the equations of Step 3 on the cover. The three lemmas above give a quotient that is a domain with a chart of its graded part, and the derivations that kill the equations descend to it.

% before: sha-counts-vd-model-sum
The local base is a completed tensor product of the universal rings at the primes of $T$, and the cover of Step 2 is a power series ring over the local base. In both cases a chart of the whole ring is assembled from charts of the factors, using the variables of all of them.

% before: sha-counts-vd-model-cover
Step 2 of the construction needs a ring with a chart that maps onto the framed ring. Adding to the local base the variables of the framed ring that are not already there gives such a ring, and the relations of the framed ring generate the kernel. Derivations of the framed ring lift to it, since the added variables are free.

% before: sha-counts-vd-model-coproduct
The local data at the primes of $T$ are classified one prime at a time, and the local base classifies them together. A family of homomorphisms from the local rings is the same as one homomorphism from their completed tensor product, which inherits a weighted presentation from the factors.

% before: sha-counts-vd-model-universal
The local base is built from universal rings for the lifts of $\bar r_0$ on the decomposition groups at the primes of $T$ that satisfy the local conditions. Conjugation by the diagonal torus acts on these rings and gives them weighted presentations. A description of the ring near the point $r_0$, by power series over the fraction field of $\mathcal O$, then gives a chart.

% before: sha-counts-vd-model-base
With a chart at each prime of $T$, the completed tensor product of the universal local rings is a local base. The numbers of variables of its chart are the sums of those at the primes of $T$.

% before: sha-counts-vd-model-lift-chart-factor
Points of a universal ring with values in finite algebras over the fraction field are not points with values in complete local rings with residue field $\mathbf F$. A compactness argument puts the values of such a point in a complete local subring, where the universal property applies.

% before: sha-counts-vd-model-lift-chart-continuity
The lifting hypothesis is stated for homomorphisms with continuous entries for the $\ell$-adic topology. Points of the universal ring give such homomorphisms, because a homomorphism to a finite $\mathbf Q_\ell$-algebra is continuous for the adic topology of the universal ring.

% before: sha-counts-vd-model-lift-chart-generation
The universal lift generates its ring in the sense of complete local rings. The same holds for homomorphisms to finite $\mathbf Q_\ell$-algebras, which lets two points over pointed artinian algebras be compared through their lifts.

% before: sha-counts-vd-model-lift-chart-complete
The charts of the model are homomorphisms to power series rings over the fraction field that are dense and injective up to nonzero constants. Here the ring need not be graded.

% before: sha-counts-vd-model-lift-chart-smooth-parameters
The chart is built from first-order data at the point $x$. Derivations at $x$ dual to elements that span the cotangent space up to constants give the linear terms of the chart.

% before: sha-counts-vd-model-lift-chart-smooth-homogeneous
To compare the powers of $\ker x$ with the powers of the maximal ideal of the power series ring, an element of the $N$-th power of $\ker x$ is written, up to a constant, as a homogeneous polynomial of degree $N$ in the parameters. The error lies in the next power.

% before: sha-counts-vd-model-lift-chart-smooth-vanishing
The images of the parameters in the chart are the variables up to a constant and terms of higher order. So a homogeneous polynomial in them of degree $N$ has order $N$ unless it vanishes.

% before: sha-counts-vd-model-lift-chart-smooth-approximation
The chart has an image that is dense up to constants. Its linear terms are the variables up to a constant, so every series is approximated, after multiplication by a constant, by images of elements of the ring.

% before: sha-counts-vd-model-lift-chart-smooth-prime
The kernel of the chart differs from the ideal $I$ only away from $x$. An element of the kernel that lies, after multiplication by a unit at $x$, in every power of $\ker x$ modulo $I$ is killed into $I$ by one element outside $\ker x$.

% before: sha-counts-vd-model-lift-chart-smooth-height
The dimension of the quotient by the kernel of the chart is computed at the section $x$. The parameters bound the height of $\ker x$ from above, the chart bounds it from below, and the dimension of the ring exceeds this height by one.

% before: sha-counts-vd-model-lift-chart-smooth
A ring whose points over pointed artinian algebras lift along surjections is formally smooth at its characteristic $0$ points. At such a point it has a complete chart, defined on the quotient by a minimal prime.

% before: sha-counts-vd-model-lift-chart-graded-change
A complete chart can be changed by an automorphism of the power series ring. If the linear terms of the images of $n$ elements of $\ker x$ are independent, the automorphism can be chosen to send these images to the variables.

% before: sha-counts-vd-model-lift-chart-graded
When the ring has a weighted presentation with origin at the point, the chart can be chosen to respect the weights. Its variables of a given weight are then counted by the derivations of that weight at the point.

% before: sha-counts-vd-model-lift-chart-factor-integral
Points of a complete local ring with values in finite extensions of $\mathbf Q_\ell$ take integral values, since an element of the maximal ideal cannot map to an element of absolute value at least $1$. This makes the image of such a ring in a finite $\mathbf Q_\ell$-algebra a finitely generated $\mathbf Z_\ell$-module, which the next three lemmas use.

% before: sha-counts-vd-model-lift-chart-generation-span
Comparing homomorphisms through generators needs two standard facts about complete local rings. Elements that separate the homomorphisms to the dual numbers generate the maximal ideal modulo its square, and homomorphisms that are continuous and agree on such generators agree everywhere.

% before: sha-counts-vd-model-lift-chart
A chart of a universal local ring at its point of characteristic $0$ comes from the deformation theory of that point over artinian algebras over the fraction field of $\mathcal O$. If every such deformation that satisfies the conditions extends along small surjections, the completion at the point is a power series ring. The torus action makes the chart graded, with as many variables of each weight as there are first-order deformations of that weight.

% before: sha-counts-vd-model-local-lifts-flag
Ordinarity at $\ell$ is defined through idempotents built from an element of inertia. For lifting it is easier to use the equivalent description by a symplectic change of basis congruent to $1$ that makes the lift upper triangular, with diagonal characters equal to the $\delta_i$ on inertia.

% before: sha-counts-vd-model-local-lifts-symplectic
The change of basis of the previous lemma must be lifted along a surjection of pointed artinian algebras. Over a field of characteristic $0$, symplectic matrices congruent to $1$ lift along surjections whose kernels consist of nilpotent elements.

% before: sha-counts-vd-model-local-lifts-unramified
The diagonal characters of an ordinary lift differ from the $\delta_i$ by unramified characters with unipotent values. Such characters lift along surjections with nilpotent kernel, and they stay unramified.

% before: sha-counts-vd-model-local-lifts-borel
Once the diagonal characters are lifted, the off-diagonal entries of an upper triangular lift must be lifted too. This needs the diagonal characters to differ from the cyclotomic character, and it rests on the published theorems on local Galois cohomology.

% before: sha-counts-vd-model-local-lifts-commuting
At the primes of $S_b$ the lifts factor through an abelian quotient of the decomposition group. A lift with commuting values and distinct residual characters is diagonalizable, and it lifts by lifting its diagonal characters.

% before: sha-counts-vd-model-local-lifts-borel-cocycle
The entries above the diagonal of an upper triangular lift are cocycles for products of the diagonal characters, up to correction terms in the entries already lifted. Such a cocycle lifts along a surjection when the second cohomology of its kernel vanishes, which local Tate duality gives as soon as the character differs from the cyclotomic one.

% before: sha-counts-vd-model-local-lifts-commuting-diagonal
A lift with commuting values whose residual characters are distinct can be diagonalized. A generic linear combination of its values has distinct eigenvalues residually, and its eigenvectors, normalized for the symplectic form, give a symplectic change of basis congruent to $1$.

% before: sha-counts-vd-model-local-lifts-commuting-characters
Once the lift is diagonal, its diagonal entries are characters congruent to the residual ones. A character with values in $1$ plus a nilpotent ideal is the exponential of an additive character, and the latter lifts along any additive section.

% before: sha-counts-vd-model-local-lifts
The chart at a prime of $T$ needs the extension property for the local conditions there. At $\ell$ these are the ordinary symplectic lifts, and at a prime of $S_b$ the symplectic lifts trivial on the inertia elements in the kernel of $\psi$.

% before: sha-counts-vd-model-local-tangent-conditions
A first-order deformation of $r_0$ is $r_0+\varepsilon\gamma r_0$ for a function $\gamma$ with values in matrices. The conditions at a prime of $T$ become linear conditions on $\gamma$: values in $\mathfrak{sp}_4$, entries below the diagonal (at $\ell$) or off the diagonal (at the primes of $S_b$) given by one coboundary, and diagonal entries vanishing on inertia.

% before: sha-counts-vd-model-local-tangent-cocycles
Above the diagonal at $\ell$, the entries of a first-order deformation are cocycles for the ratios of the characters $\delta_a$. The dimension of the space of these cocycles rests on the published theorems on local Galois cohomology.

% before: sha-counts-vd-model-local-tangent-unramified-cocycle
A cocycle of the local Galois group that vanishes on inertia is determined by its value at a Frobenius element. Every value occurs, since a homomorphism trivial on inertia can send a Frobenius element to any element of a profinite group. Applied to the semidirect product of a stable lattice by the Galois group, this gives the cocycle.

% before: sha-counts-vd-model-local-tangent-unramified
On the diagonal, the entries of a first-order deformation are homomorphisms to $E$ that vanish on inertia. Such a homomorphism is determined by its value at a Frobenius element, and every value occurs.

% before: sha-counts-vd-model-local-tangent-root
Each weight other than $(0,0)$ occurs at one pair of entries related by the symplectic condition. So a matrix of $\mathfrak{sp}_4$ supported on that weight is determined by one entry.

% before: sha-counts-vd-model-local-tangent
The number of variables of each weight in the local chart is the dimension of a space of first-order deformations of $r_0$. These dimensions are the numbers $d(j)$ at $\ell$ and $d'(j)$ at the primes of $S_b$ that the local base requires.

% before: sha-counts-vd-model-local-charts
With these inputs, each universal local ring at a prime of $T$ has the chart required by the local base. The conditions at the prime are polynomial identities in the entries, so they pass to images and are detected after an injective map.

% before: sha-counts-vd-model-descent
The framed ring is built over the Galois group of the maximal extension of $\mathbf Q$ unramified outside $S$, through which the points of a Siegel datum factor. Through this descent the symplectic condition and the triviality on inertia away from $T$ become conditions on lifts of that Galois group, the second one at the primes of $S$ outside $T$.

% before: sha-counts-vd-model-framed-cohomology
The numbers of variables and of relations of the framed ring are dimensions of cohomology groups of the weight spaces of $\mathfrak{sp}_4$. The groups are framed: cochains come with local trivializations at the primes of $T$, normalized at $\ell$, and the cocycles vanish on the subgroups $N_q$ at the primes of $S$ outside $T$.

% before: sha-counts-vd-model-place-datum
The counts of the framed cohomology groups, and the presentation of the framed ring, use only a group with maps from local groups at two finite sets of places. Stating the definitions for such a datum separates the linear algebra of cochains from the Galois groups, which enter only through the published theorems on their cohomology.

% before: sha-counts-vd-model-selmer
The framed cocycles differ from the classes of the Selmer group with trivial restrictions at $T$ by the choices of the local trivializations. Counting these choices gives the dimension of the framed cocycles in terms of a Selmer group and invariants.

% before: sha-counts-vd-model-obstruction
The framed obstruction group sits between the local classes that the global classes do not explain and the global classes of degree $2$ that vanish locally. Its dimension is therefore a sum of local and global dimensions, which the Euler characteristic formulas then evaluate.

% before: sha-counts-vd-model-unramified
At the primes of $S$ outside $T$ the local condition asks for classes that vanish on the subgroup $N_q$ of the inertia group. Since $N_q$ has index prime to $\ell$, this is the unramified condition, whose dimension is that of the invariants.

% before: sha-counts-vd-model-euler
For the Galois group of $\mathbf Q$ unramified outside $S$, the local and global Euler characteristic formulas and Poitou-Tate duality turn the two previous dimension formulas into one. The obstruction group exceeds the framed cocycles by the invariants of complex conjugation and the twisted dual invariants.

% before: sha-counts-vd-model-counts
For the count of the model, each weight must carry as many relations as variables, up to a correction at the weights $(0,0)$ and $\pm(0,1)$. The correction comes from the global Euler characteristic, and the last condition of a Siegel datum removes the variables of the other weights.

% before: sha-counts-vd-model-form-datum
The presentation of the framed ring holds for a general group with local conditions, and it is proved in that generality. The data are a bilinear form preserved up to scalars, a grading of the matrix entries by an abelian group, and conditions at finitely many places that behave like the conditions of a deformation problem: they are cut out by equations, stable under conjugation and under the torus, liftable along small extensions, and with a tangent space described by a subgroup.

% before: sha-counts-vd-model-lemma-p-situation
Four of the lemmas below and the presentation itself share one setup, to which each adds the hypotheses it needs. It fixes the group, the residual representation and its lift with values in $\mathcal O$, both normalising the form, the place datum with its local conditions, and the local lifts over the base ring.

% before: sha-counts-vd-model-lemma-p-universal
The first step is the existence of a universal framing. It is built from the universal lift of Mazur by adjoining the entries of the framing matrices and dividing by the equations of the framing and of the local conditions.

% before: sha-counts-vd-model-lemma-p-grading
The torus of the weights acts on framings by conjugation, and the universal property turns this action into a grading of the universal ring. This gives a weighted presentation in which the entries of the universal framing are homogeneous of the weights of their positions.

% before: sha-counts-vd-model-lemma-p-tangent-abstract
The relative cotangent space of $R$ over $B$ is dual to the homomorphisms to the dual numbers that are trivial on $B$. When these homomorphisms are described, weight by weight, by vector spaces with linear forms given by elements of $\mathfrak m_R$, finitely many of these elements form a basis of the relative cotangent space.

% before: sha-counts-vd-model-lemma-p-tangent
The variables of the presentation come from the tangent space of the universal ring relative to the base. Homomorphisms to the dual numbers over the residue field are framings with values there, and these correspond to framed cocycles, weight by weight.

% before: sha-counts-vd-model-lemma-p-relations-parts
A relative presentation replaces the variables of a given weighted presentation by other homogeneous elements. The parts of the ring must not change in the process, so that the weights of the entries of the framing are still read off from the new presentation.

% before: sha-counts-vd-model-lemma-p-relations-kernel
The relations of a weighted presentation must be counted weight by weight. This needs the kernel to be generated by finitely many series, each of a single weight, which follows from noetherianity and the theorem of Krull.

% before: sha-counts-vd-model-lemma-p-relations
Given generators of the maximal ideal over the base, the kernel of the corresponding presentation is generated by the kernel of the base and finitely many series. They can be chosen homogeneous and minimal, which is what allows the relations to be counted.

% before: sha-counts-vd-model-lemma-p-unitary
Framings and local conditions are stated with $\Phi$-unitary matrices, and lifting along a small extension changes such a matrix by an element of the Lie algebra. Since $2$ is a unit, every $\Phi$-unitary matrix over the smaller ring lifts.

% before: sha-counts-vd-model-lemma-p-classes-section-lift
The obstruction classes are computed from lifts to $A$ of the universal framing, and these lifts must have homogeneous entries that depend continuously on the group element. A continuous function with values in the image of a part lifts to a continuous function with values in that part.

% before: sha-counts-vd-model-lemma-p-classes-section
The values of a framing are $\Phi$-unitary up to a constant matrix, and a homogeneous lift of a $\Phi$-unitary matrix need not be $\Phi$-unitary. The correction of Lemma~\ref{sha-counts-vd-model-lemma-p-unitary} is a polynomial in the entries of the lift, so it keeps their weights and their continuity.

% before: sha-counts-vd-model-lemma-p-classes-local-split
A lift with the local condition and a homogeneous lift differ by elements of $\ker p$ of several weights. To compare them, a square-zero extension of $A$ by copies of $\ker p$ is embedded in a power series ring over $A$, by a torus element whose values are independent over $A$. In that ring an element of the part of weight $j$ is $T(j)$ times its image under a twisted homomorphism.

% before: sha-counts-vd-model-lemma-p-classes-local
At the places of $Q$ the lift must also keep the local condition. The lifting property of a local condition gives some lift with the condition, and the torus property carries the condition over to a homogeneous lift built from it.

% before: sha-counts-vd-model-lemma-p-classes-coefficients
The kernel of the small extension is a vector space over the residue field with a basis given by the relations. The coordinates in this basis vanish on the homogeneous elements of the kernel whose weight is not that of the coordinate.

% before: sha-counts-vd-model-lemma-p-classes-cocycle
A coordinate of the failure of the lifts to form a framing is a framed $2$-cocycle. Its values lie in the weight space of the coordinate, because the lifts are homogeneous.

% before: sha-counts-vd-model-lemma-p-classes-correction
When the framed $2$-cocycle given by a combination of coordinates is a framed coboundary, the coboundary corrects the lifts. The corrected lifts form a framing over the quotient on which only that combination survives.

% before: sha-counts-vd-model-lemma-p-classes
The relations are counted through the obstruction to lifting the universal framing along the small extension that they define. Each relation gives a class in the framed obstruction group, and a vanishing combination of classes lets the framing lift.

% before: sha-counts-vd-model-lemma-p
The framed deformation ring over a local base has a presentation over it whose variables and relations are counted by the framed cohomology groups. The torus action makes the presentation homogeneous, and the counts hold weight by weight. The lemma holds for the general data of the situation above; the chapter uses it only for the Galois group $G_S$ and the symplectic form, in Corollary~\ref{sha-counts-vd-model-lemma-p-siegel}.

% before: sha-counts-vd-model-lemma-p-kill-section
To lift a homomorphism along a small extension, its values are first lifted one at a time. The lifts can be chosen so that they depend continuously on the values, because the kernel of the extension meets a power of the maximal ideal only in $0$.

% before: sha-counts-vd-model-lemma-p-kill-coboundary
The defect of such a lift is a $2$-cocycle with values in the Lie algebra over the kernel of the extension. This kernel is a vector space over the residue field, so the defect splits into cocycles with values in the weight spaces over the residue field, where the hypothesis on $2$-cocycles applies.

% before: sha-counts-vd-model-lemma-p-kill
At the primes of $S$ outside $T$ the local condition asks that the lift be trivial on a subgroup of inertia. Such a condition has the properties of a local condition as soon as the second cohomology of the quotient by the subgroup vanishes for the adjoint coefficients.

% before: sha-counts-vd-model-lemma-p-siegel-transfer
A $2$-cocycle that is a coboundary on a subgroup of finite index is a coboundary on the whole group, when the index is invertible in the coefficients. The trivialization on the whole group is an average, over the cosets of the subgroup, of translates of the given one.

% before: sha-counts-vd-model-lemma-p-siegel-tame
At a prime of $Q$ the condition is triviality on the subgroup $N_q$ of inertia, and Lemma~\ref{sha-counts-vd-model-lemma-p-kill} asks for trivializations of the $2$-cocycles that are invariant under $N_q$. Modulo a small open subgroup, such a cocycle is trivialized along the powers of a Frobenius element. The subgroup they generate has index prime to $\ell$, as $[I_q:N_q]$ has, so the previous lemma extends the trivialization.

% before: sha-counts-vd-model-lemma-p-siegel
For the Galois group $G_S$ the form is the symplectic form, the residual and base representations are diagonal, and the conditions at the primes of $Q$ are triviality on the subgroups $N_q$. The general lemma then gives the presentation of the framed deformation ring of $G_S$ used to build the framed ring of the model.

% before: sha-counts-vd-model-framed
These ingredients give the framed ring of Definition~\ref{sha-counts-vd-model-framing}. The descent turns framings over $G_S$ into framings over $G$, the universal property supplies the origin, and the counts of the two previous lemmas give the numbers of variables and series of each weight.

% before: sha-counts-vd-model-exists
The bound applies to a Siegel datum once it has a graded model. As for the unitary group in Proposition~\ref{sha-model-exists}, the model is built from presentations of the local deformation rings at the primes of $T$ and of the global ring over them.

% before: sha-counts-vd-bound
The bound for $V_D$ in part (b) of Proposition~\ref{sha-counts} comes from the commutative algebra applied to a graded model. Each element of a finite set meeting the two-signed orbits gives an irreducible $\bar{\mathbf Q}_\ell$-point, and the valuation at it of the element of weight $0$ is that of a difference of traces of a Frobenius element.

% before: sha-counts
Each case has its own bound. Cases I, II, III and V use the automorphy of residually reducible representations for $U(3)$ or $GSp_4$; case Z uses the vanishing for the trivial character; case IV, at the bad primes, uses the bounds on the fixed layers that hold on controlled windows.

% before: sha-plus-part
Adding the bounds of the cases over all characters, twists and primes gives the plus bound on controlled windows. The result is polynomial in the conductor and the twist, plus the error $\exp(x^{4\varepsilon/5})$ coming from the bad primes.

% before: sha-weakening-from-galois-representations
The minus bound, the bound at $2$ and the plus bound now combine on a common window. Only the plus part uses the Galois representations of Chapter~\ref{chapter-galois-representations}; the minus part and the prime $2$ are unconditional.
